% Personnalité et personne — Philosophie 1re Tech — Cours signé OnBuch+
% Programme officiel MINESEC. Compiler avec : tectonic N05-personnalite-personne.tex
\newcommand{\DOCMATIERE}{Philosophie}
\newcommand{\DOCNIVEAU}{1re Tech}
\newcommand{\DOCMODULE}{Philosophie – classe de Première technique}
\newcommand{\DOCLECON}{N05}
\newcommand{\DOCTITRE}{Personnalité et personne}
% OnBuch+ — préambule « cours signés » ENRICHI (compilé avec tectonic / XeLaTeX)
% Système visuel « Pop » d'OnBuch+ : fond crème (jamais blanc), bordures encre
% épaisses, ombres « dures » décalées, titres Archivo Black en capitales.
% Version MATHÉMATIQUES : graphes (pgfplots/tikz), boîtes pédagogiques
% (définition, propriété, méthode, attention, le savais-tu, exercice résolu,
% exercice, TP, bilan), page de garde et signature OnBuch+ ancrée en bas.
%
% Macros attendues dans main.tex AVANT \input{preamble} :
%   \DOCMATIERE  \DOCNIVEAU  \DOCMODULE  \DOCLECON (numéro)  \DOCTITRE
\documentclass[11pt]{article}

\usepackage{fontspec}
\usepackage[french]{babel}
\usepackage[a4paper,margin=2cm,top=3.4cm,bottom=2.8cm,headheight=1.4cm,footskip=1.3cm]{geometry}
\usepackage{amsmath,amssymb,amsfonts}
\usepackage{mathtools} % \xmapsto, \coloneqq... souvent utilisés par les modèles
\usepackage{cancel} % simplifications barrées (bilans de forces)
% \abs{z} (module, valeur absolue) et \norm{u} (norme), utilisés par les modèles
\providecommand{\abs}{}\let\abs\relax\DeclarePairedDelimiter{\abs}{\lvert}{\rvert}
\providecommand{\norm}{}\let\norm\relax\DeclarePairedDelimiter{\norm}{\lVert}{\rVert}
\usepackage[table]{xcolor} % [table] : fournit aussi \rowcolors (alternance de lignes)
\usepackage{pagecolor}
\usepackage{tikz}
\usetikzlibrary{arrows.meta,positioning,calc,shapes.geometric,shapes.misc,shapes.arrows,decorations.pathmorphing,decorations.pathreplacing,decorations.markings,patterns,shadows,mindmap,trees,fit,backgrounds,matrix,intersections,angles,quotes}
\usepackage{pgfplots}
\usepackage[version=4]{mhchem} % équations nucléaires \ce{^{235}_{92}U}
\usepackage[european,siunitx]{circuitikz} % schémas électriques
\pgfplotsset{compat=1.18}
\usepgfplotslibrary{fillbetween,groupplots} % aires entre courbes (calcul intégral)
\usepackage{enumitem}
\usepackage{fancyhdr}
% [most] charge toutes les bibliothèques tcolorbox (listings, minted, raster,
% documentation...) et gonfle inutilement la mémoire TeX sur un document long
% (~300 boîtes) : on ne charge que ce qui est réellement utilisé.
\usepackage[skins,breakable]{tcolorbox}
\usepackage{graphicx}
\usepackage{colortbl}
\usepackage{booktabs}
\usepackage{tabularx}
\usepackage{multirow}
\usepackage{diagbox} % case d'en-tête diagonale des tableaux à double entrée
% Colonnes extensibles centrée / gauche / droite (C, L, R), souvent utilisées par les modèles
\newcolumntype{C}{>{\centering\arraybackslash}X}
\newcolumntype{L}{>{\raggedright\arraybackslash}X}
\newcolumntype{R}{>{\raggedleft\arraybackslash}X}
\usepackage{array}
\usepackage{multicol}
\usepackage{siunitx}
% parse-numbers=false : accepte aussi \SI{2,5 \times 10^{-3}}{...}
\sisetup{locale=FR,output-decimal-marker={,},group-minimum-digits=4,parse-numbers=false}
\DeclareSIUnit\ohm{\ensuremath{\symup{\Omega}}} % U+2126 absent de la police
\DeclareSIUnit\division{div} % réglages d'oscilloscope (ms/div, V/div)
\DeclareSIUnit\tr{tr} % tours (vitesse de rotation, ex. \SI{1500}{\tr\per\minute})
\DeclareSIUnit\molar{M} % concentration molaire (ex. \SI{3}{\molar}, \SI{1}{\micro\molar})
\DeclareSIUnit\an{an} % année (le modèle confond parfois avec \year, un compteur TeX) — ex. \SI{5}{\kilo\meter\per\an}
\DeclareSIUnit\FCFA{FCFA} % franc CFA (ex. \SI{500}{\FCFA\per\kilo\gram})
\DeclareSIUnit\degreeBrix{\textdegree Brix} % teneur en sucre d'un sirop/jus (réfractométrie)
\DeclareSIUnit\month{mois} % le modèle utilise parfois \month, un compteur TeX, comme unité de durée
\DeclareSIUnit\microliter{\micro\liter} % le modèle écrit parfois \microliter en un seul token
\DeclareSIUnit\million{millions} % dénombrement (ex. \SI{15}{\million\per\mL} spermatozoïdes)
\usepackage{qrcode}
% \degree : commande fréquemment utilisée par les modèles pour un angle ou une
% température en dehors de \SI{}{} (non fournie par les paquets chargés).
\providecommand{\degree}{^{\circ}}
% \qed (fin de démonstration), utilisé par les modèles sans amsthm
\providecommand{\qed}{\hfill$\square$}
% \barre : double barre des valeurs interdites dans un tableau de variations (tkz-tab)
\providecommand{\barre}{\ensuremath{\|}}
% environnement proof (amsthm non chargé) utilisé par les modèles
\ifdefined\proof\else\newenvironment{proof}[1][Démonstration]{\par\noindent\textit{#1.}\ }{\qed\par}\fi
\usepackage{lastpage}
\usepackage{hyperref}
\hypersetup{hidelinks}

% ============================================================
% Palette « Pop » OnBuch+ — mêmes hex que la classe P (plus_ui.dart)
% ============================================================
\definecolor{poporange}{HTML}{F59321}
\definecolor{popdark}{HTML}{E07A0C}
\definecolor{popcream}{HTML}{FFF6E8}
\definecolor{popink}{HTML}{1C1714}
\definecolor{poppurple}{HTML}{7A5AE0}
\definecolor{popgreen}{HTML}{2E9E6A}
\definecolor{poppink}{HTML}{E0457B}
\definecolor{popblue}{HTML}{2F7DE1}
\definecolor{popgold}{HTML}{C98A12}
\definecolor{popmuted}{HTML}{8A7F76}
\definecolor{popline}{HTML}{EADFD2}
\definecolor{popgray}{HTML}{8A7F76}
\colorlet{popgrey}{popgray}
% Alias tolérants (le modèle nomme parfois les couleurs différemment)
\colorlet{popwhite}{white}
\colorlet{popblack}{popink}
\colorlet{popred}{poppink}
\colorlet{popyellow}{popgold}
% Teintes claires (fonds de tableaux/encadrés) — le modèle nomme parfois
% les couleurs avec un suffixe L (ex. popblueL) sans les avoir définies.
\colorlet{poporangeL}{poporange!20}
\colorlet{popdarkL}{popdark!20}
\colorlet{poppurpleL}{poppurple!20}
\colorlet{popgreenL}{popgreen!20}
\colorlet{poppinkL}{poppink!20}
\colorlet{popblueL}{popblue!20}
\colorlet{popgoldL}{popgold!20}
\colorlet{popredL}{popred!20}
\colorlet{popgrayL}{popgray!20}
% Teintes claires pour fonds de tableaux / zones
\colorlet{popblueL}{popblue!10!white}
\colorlet{poporangeL}{poporange!14!white}
\colorlet{popgreenL}{popgreen!12!white}
\colorlet{poppurpleL}{poppurple!12!white}
\colorlet{poppinkL}{poppink!12!white}
\colorlet{popgoldL}{popgold!14!white}

\pagecolor{popcream}

% ============================================================
% Typographie
% ============================================================
\defaultfontfeatures{Ligatures=TeX}
\setmainfont{PlusJakartaSans-Medium.ttf}[Path=../../../fonts/,BoldFont=PlusJakartaSans-Bold.ttf]
% Maths en sans-serif (Fira Math) avec chiffres et lettres droites de Plus
% Jakarta, pour que les formules se fondent dans le texte.
\usepackage[math-style=ISO,bold-style=ISO]{unicode-math}
\setmathfont{FiraMath-Regular.otf}[Path=../../../fonts/]
\setmathfont{PlusJakartaSans-Medium.ttf}[Path=../../../fonts/,range={up/num,up/{Latin,latin},\mathup}]
\usepackage{icomma} % virgule décimale sans espace en mode math
% Caractères mathématiques Unicode absents des polices de texte (Plus Jakarta,
% Archivo Black) : rendus par leur équivalent mathématique, en texte comme en
% formule — sinon ils s'affichent en carré vide (ex. « Arithmétique dans ℤ »).
\usepackage{newunicodechar}
\newunicodechar{ℝ}{\ensuremath{\mathbb{R}}}
\newunicodechar{ℤ}{\ensuremath{\mathbb{Z}}}
\newunicodechar{ℂ}{\ensuremath{\mathbb{C}}}
\newunicodechar{ℕ}{\ensuremath{\mathbb{N}}}
\newunicodechar{ℚ}{\ensuremath{\mathbb{Q}}}
\newunicodechar{𝔻}{\ensuremath{\mathbb{D}}}
\newunicodechar{α}{\ensuremath{\alpha}}
\newunicodechar{β}{\ensuremath{\beta}}
\newunicodechar{γ}{\ensuremath{\gamma}}
\newunicodechar{δ}{\ensuremath{\delta}}
\newunicodechar{ε}{\ensuremath{\varepsilon}}
\newunicodechar{θ}{\ensuremath{\theta}}
\newunicodechar{λ}{\ensuremath{\lambda}}
\newunicodechar{μ}{\ensuremath{\mu}}
\newunicodechar{σ}{\ensuremath{\sigma}}
\newunicodechar{τ}{\ensuremath{\tau}}
\newunicodechar{φ}{\ensuremath{\varphi}}
\newunicodechar{ψ}{\ensuremath{\psi}}
\newunicodechar{ω}{\ensuremath{\omega}}
\newunicodechar{Σ}{\ensuremath{\Sigma}}
% majuscules produites par \MakeUppercase dans les titres (α-aminés -> Α)
\newunicodechar{Α}{\ensuremath{\alpha}}
\newunicodechar{Β}{\ensuremath{\beta}}
\newunicodechar{Γ}{\ensuremath{\Gamma}}
\newunicodechar{Δ}{\ensuremath{\Delta}}
\newunicodechar{Ε}{\ensuremath{\varepsilon}}
\newunicodechar{Θ}{\ensuremath{\Theta}}
\newunicodechar{Λ}{\ensuremath{\Lambda}}
\newunicodechar{Μ}{\ensuremath{\mu}}
\newunicodechar{Τ}{\ensuremath{\tau}}
\newunicodechar{Φ}{\ensuremath{\Phi}}
\newunicodechar{Ψ}{\ensuremath{\Psi}}
\newunicodechar{Ω}{\ensuremath{\Omega}}
\newunicodechar{↦}{\ensuremath{\mapsto}}
\newunicodechar{∈}{\ensuremath{\in}}
\newunicodechar{∉}{\ensuremath{\notin}}
\newunicodechar{∧}{\ensuremath{\wedge}}
\newunicodechar{⇔}{\ensuremath{\Leftrightarrow}}
\newunicodechar{⟺}{\ensuremath{\Longleftrightarrow}}
\newunicodechar{⇒}{\ensuremath{\Rightarrow}}
\newunicodechar{⟹}{\ensuremath{\Longrightarrow}}
\newunicodechar{∘}{\ensuremath{\circ}}
\newunicodechar{∪}{\ensuremath{\cup}}
\newunicodechar{∩}{\ensuremath{\cap}}
\newunicodechar{≡}{\ensuremath{\equiv}}
\newunicodechar{⊂}{\ensuremath{\subset}}
\newunicodechar{⊥}{\ensuremath{\perp}}
\newunicodechar{∃}{\ensuremath{\exists}}
\newunicodechar{∀}{\ensuremath{\forall}}
\newunicodechar{∛}{\ensuremath{\sqrt[3]{\,}}}
\newunicodechar{∜}{\ensuremath{\sqrt[4]{\,}}}
\newunicodechar{⃗}{\ensuremath{{}^{\to}}}
\newunicodechar{ⁿ}{\ifmmode^{n}\else\textsuperscript{\ensuremath{n}}\fi}
\newunicodechar{ˣ}{\ifmmode^{x}\else\textsuperscript{\ensuremath{x}}\fi}
\newunicodechar{ᵇ}{\ifmmode^{b}\else\textsuperscript{\ensuremath{b}}\fi}
\newunicodechar{ʳ}{\ifmmode^{r}\else\textsuperscript{\ensuremath{r}}\fi}
\newunicodechar{ᵘ}{\ifmmode^{u}\else\textsuperscript{\ensuremath{u}}\fi}
\newunicodechar{ʸ}{\ifmmode^{y}\else\textsuperscript{\ensuremath{y}}\fi}
\newunicodechar{ᵃ}{\ifmmode^{a}\else\textsuperscript{\ensuremath{a}}\fi}
\newunicodechar{ᵉ}{\ifmmode^{e}\else\textsuperscript{\ensuremath{e}}\fi}
\newunicodechar{ᵏ}{\ifmmode^{k}\else\textsuperscript{\ensuremath{k}}\fi}
\newunicodechar{ᵖ}{\ifmmode^{p}\else\textsuperscript{\ensuremath{p}}\fi}
\newunicodechar{ᴺ}{\ifmmode^{N}\else\textsuperscript{\ensuremath{N}}\fi}
\newunicodechar{ᶜ}{\ifmmode^{c}\else\textsuperscript{\ensuremath{c}}\fi}
\newunicodechar{ʰ}{\ifmmode^{h}\else\textsuperscript{\ensuremath{h}}\fi}
\newunicodechar{ᵠ}{\ifmmode^{q}\else\textsuperscript{\ensuremath{q}}\fi}
\newunicodechar{ₙ}{\ifmmode_{n}\else\textsubscript{\ensuremath{n}}\fi}
\newunicodechar{ₐ}{\ifmmode_{a}\else\textsubscript{\ensuremath{a}}\fi}
\newunicodechar{ᵢ}{\ifmmode_{i}\else\textsubscript{\ensuremath{i}}\fi}
\newunicodechar{ᵣ}{\ifmmode_{r}\else\textsubscript{\ensuremath{r}}\fi}
\newunicodechar{ₖ}{\ifmmode_{k}\else\textsubscript{\ensuremath{k}}\fi}
\newunicodechar{ᵦ}{\ifmmode_{\beta}\else\textsubscript{\ensuremath{\beta}}\fi}
\newunicodechar{ₓ}{\ifmmode_{x}\else\textsubscript{\ensuremath{x}}\fi}
\newfontfamily\popdisplay{ArchivoBlack-Regular.ttf}[Path=../../../fonts/]
\newfontfamily\poplogo{SpaceGrotesk-Bold.ttf}[Path=../../../fonts/]
\newfontfamily\popsemi{PlusJakartaSans-SemiBold.ttf}[Path=../../../fonts/]
\newfontfamily\popxbold{PlusJakartaSans-ExtraBold.ttf}[Path=../../../fonts/]
\newfontfamily\popxboldls{PlusJakartaSans-ExtraBold.ttf}[Path=../../../fonts/,LetterSpace=3.0]

\setlist[enumerate]{leftmargin=1.7em,itemsep=5pt,topsep=4pt}
\setlist[itemize]{leftmargin=1.7em,itemsep=4pt,topsep=4pt,label={\color{poporange}\textbullet}}
\setlength{\parindent}{0pt}
\setlength{\parskip}{5pt}
\renewcommand{\arraystretch}{1.25}

% pgfplots : style Pop par défaut
\pgfplotsset{
  every axis/.append style={axis line style={popink,line width=1pt},tick style={popink},
    grid style={popline,line width=0.5pt},label style={font=\small},tick label style={font=\footnotesize}},
  popcurve/.style={poporange,line width=1.8pt,no markers,smooth},
}
% Même style disponible dans un \draw TikZ simple (hors pgfplots)
\tikzset{popcurve/.style={poporange,line width=1.8pt}}
\tikzset{popgrid/.style={popline,very thin}}
\colorlet{popcurve}{poporange} % le nom du style est parfois utilisé comme couleur

% ============================================================
% Pill / squiggle
% ============================================================
\newcommand{\poppill}[3]{%
  \tikz[baseline=-0.55ex]\node[draw=popink,line width=1.2pt,rounded corners=2.6pt,%
    fill=#2,inner xsep=6.5pt,inner ysep=3pt]{%
    \popxboldls\fontsize{7}{7}\selectfont\color{#3}\MakeUppercase{#1}};%
}
\newcommand{\popsquiggle}[1]{%
  \tikz[baseline=-0.4ex]{
    \draw[#1,line width=2.6pt,line cap=round]
      plot [smooth] coordinates {(0,0.09) (0.34,0.01) (0.68,0.21) (1.02,0.01) (1.36,0.10)};
  }%
}
% Mot-clé surligné (marqueur orange sous le texte)
\newcommand{\cle}[1]{{\popxbold\color{popdark}#1}}

% ============================================================
% En-tête / pied
% ============================================================
\pagestyle{fancy}
\fancyhf{}
\renewcommand{\headrulewidth}{0pt}
\renewcommand{\footrulewidth}{0pt}
\lhead{\raisebox{-0.15ex}{\poplogo\fontsize{13}{13}\selectfont\color{popink}OnBuch\color{poporange}+}}
\rhead{\poppill{\DOCMATIERE\ \textbullet\ \DOCNIVEAU\ \textbullet\ Leçon \DOCLECON}{white}{popink}}
\lfoot{\popxbold\fontsize{7.6}{7.6}\selectfont\color{popmuted}\MakeUppercase{Cours signé OnBuch+}\ \textbullet\ {\popsemi\DOCTITRE}}
\rfoot{\tikz[baseline=-0.6ex]\node[rounded corners=8.5pt,fill=popink,minimum height=17pt,minimum width=17pt,inner xsep=5pt,inner sep=0]{\popxbold\fontsize{8}{8}\selectfont\color{popcream}\thepage\,/\,\pageref*{LastPage}};}

% ============================================================
% Titres
% ============================================================
\newcommand{\chaptitre}[2]{%
  \begin{tcolorbox}[enhanced,colback=poporange,colframe=popink,boxrule=2.6pt,arc=11pt,
      drop shadow=popink,left=15pt,right=15pt,top=12pt,bottom=11pt,before skip=3pt,after skip=18pt]
    \poppill{#2}{popink}{popcream}\par\vspace{9pt}
    {\raggedright\hyphenpenalty=10000\exhyphenpenalty=10000\popdisplay\fontsize{22}{24}\selectfont\color{popcream}\MakeUppercase{#1}\par}
  \end{tcolorbox}
}
% Section numérotée : pastille orange avec le numéro + titre Archivo
\newcounter{popsec}
\newcommand{\coursec}[1]{%
  \stepcounter{popsec}\setcounter{popsub}{0}%
  \par\vspace{16pt}\noindent
  \begin{minipage}{\linewidth}
  \tikz[baseline=-0.7ex]\node[draw=popink,line width=1.6pt,fill=poporange,rounded corners=4pt,minimum size=22pt,inner sep=0]{\popdisplay\fontsize{12}{12}\selectfont\color{popcream}\thepopsec};%
  \hspace{8pt}{\popdisplay\fontsize{15}{17}\selectfont\color{popink}\MakeUppercase{#1}}\par
  \vspace{4pt}\noindent{\color{poporange}\rule{36pt}{3.6pt}}
  \end{minipage}\par\nopagebreak\vspace{8pt}
}
\newcounter{popsub}
\newcommand{\courssub}[1]{%
  \stepcounter{popsub}%
  \par\vspace{9pt}\noindent{\popxbold\fontsize{11.5}{13}\selectfont\color{popink}{\color{poporange}\thepopsec.\thepopsub}\hspace{6pt}#1}\par\nopagebreak\vspace{4pt}
}

% ============================================================
% Boîtes pédagogiques — même recette : fond blanc, bordure épaisse colorée,
% ombre dure encre, badge plein. Titre optionnel [#1].
% ============================================================
\tcbset{popbase/.style={breakable,enhanced,colback=white,boxrule=2.3pt,arc=9pt,
  shadow={3pt}{-3pt}{0pt}{popink},left=14pt,right=14pt,top=11pt,bottom=11pt,before skip=11pt,after skip=15pt,
  fontupper=\popsemi\fontsize{10.6}{14.5}\selectfont\color{popink},
  fontlower=\popsemi\fontsize{10.6}{14.5}\selectfont\color{popink}}}
% Coche dessinée (le glyphe ✓ n'existe pas dans les polices du document)
\newcommand{\popcheck}[1]{\tikz[baseline=-0.5ex]\draw[#1,line width=1.6pt,line cap=round,line join=round](-0.13,0.0)--(-0.04,-0.09)--(0.13,0.1);}
\providecommand{\coche}{\popcheck{popgreen}}
\providecommand{\resultat}[1]{\par\smallskip\noindent\fcolorbox{popgreen}{popgreenL}{\parbox{\dimexpr\linewidth-2\fboxsep-2\fboxrule\relax}{\textbf{Résultat.} #1}}\par\smallskip}
\newcommand{\popboxhead}[3]{\poppill{#1}{#2}{white}\ifx\relax#3\relax\else\hspace{6pt}{\popxbold\fontsize{10.6}{12}\selectfont\color{popink}#3}\fi\par\vspace{7pt}}

\newtcolorbox{aretenir}[1][]{popbase,colframe=popgreen,before upper={\popboxhead{À retenir}{popgreen}{#1}}}
\newtcolorbox{exemplebox}[1][]{popbase,colframe=poporange,before upper={\popboxhead{Exemple}{poporange}{#1}}}
\newtcolorbox{definition}[1][]{popbase,colframe=popblue,before upper={\popboxhead{Définition}{popblue}{#1}}}
\newtcolorbox{propriete}[1][]{popbase,colframe=poppurple,before upper={\popboxhead{Propriété}{poppurple}{#1}}}
\newtcolorbox{methode}[1][]{popbase,colframe=popgold,before upper={\popboxhead{Méthode}{popgold}{#1}}}
\newtcolorbox{attention}[1][]{popbase,colframe=poppink,before upper={\popboxhead{Attention}{poppink}{#1}}}
\newtcolorbox{experience}[1][]{popbase,colframe=popblue,colback=popblueL,before upper={\popboxhead{Expérience}{popblue}{#1}}}
% « Le savais-tu ? » : carte violette pleine (contexte, culture, Cameroun)
\newtcolorbox{savaistu}[1][]{popbase,colback=poppurple,colframe=popink,
  fontupper=\popsemi\fontsize{10.4}{14.2}\selectfont\color{white},
  before upper={\poppill{Le savais-tu\,?}{white}{poppurple}\ifx\relax#1\relax\else\hspace{6pt}{\popxbold\color{white}#1}\fi\par\vspace{7pt}}}
% Exercice résolu : énoncé en haut, \tcblower puis la solution sur fond crème
\newcounter{popexo}\newcounter{popexr}
\newtcolorbox{exoresolu}[1][]{popbase,colframe=poporange,colbacklower=popcream,
  segmentation style={popink,line width=1.2pt,dashed},
  before upper={\stepcounter{popexr}\popboxhead{Exercice résolu \thepopexr}{poporange}{#1}},
  before lower={\poppill{Solution}{popink}{popcream}\par\vspace{6pt}}}
% Exercice (non résolu en place) — la correction est dans la partie Corrigés
\newtcolorbox{exercice}[1][]{popbase,colframe=popink,
  before upper={\stepcounter{popexo}\popboxhead{Exercice \thepopexo}{popink}{#1}}}
% Corrigé
\newtcolorbox{corrige}[1][]{popbase,colframe=popgreen,colback=popgreenL,
  before upper={\popboxhead{Corrigé}{popgreen}{#1}}}
% Objectifs / prérequis / bilan
\newtcolorbox{objectifs}{popbase,colframe=popink,colback=poporangeL,
  before upper={\popboxhead{Objectifs de la leçon}{popink}{}}}
\newtcolorbox{prerequis}{popbase,colframe=popink,colback=popblueL,
  before upper={\popboxhead{Prérequis}{popblue}{}}}
\newtcolorbox{bilan}{popbase,colframe=popink,colback=white,boxrule=2.8pt,
  before upper={\popboxhead{Fiche bilan}{poporange}{L'essentiel à connaître pour l'examen}}}
% Figure encadrée avec légende
\newtcolorbox{popfigure}[1][]{enhanced,breakable=false,colback=white,colframe=popline,boxrule=1.4pt,arc=8pt,
  left=10pt,right=10pt,top=10pt,bottom=8pt,before skip=10pt,after skip=12pt,halign=center,
  lower separated=false,halign lower=center,
  fontlower=\popsemi\fontsize{9}{11.5}\selectfont\color{popmuted},
  #1}
\newcommand{\legende}[1]{\par\vspace{6pt}{\popsemi\fontsize{9}{11.5}\selectfont\color{popmuted}#1}}

% ============================================================
% Signature OnBuch+
% ============================================================
\newcommand{\poplogoinline}[1]{{\poplogo\fontsize{#1}{#1}\selectfont\color{popink}OnBuch\color{poporange}+}}
\newcommand{\signatureonbuch}{%
  \begin{tcolorbox}[enhanced,colback=popink,colframe=popink,boxrule=2.2pt,arc=10pt,
      drop shadow=poporange,left=14pt,right=14pt,top=10pt,bottom=10pt,before skip=0pt,after skip=0pt]
    \tikz[baseline=-0.7ex]\node[circle,fill=popgreen,draw=popcream,line width=1.3pt,minimum size=22pt,inner sep=0]{\popcheck{white}};%
    \hspace{9pt}\begin{minipage}[c]{0.62\linewidth}
      {\popxbold\fontsize{12}{13}\selectfont\color{popcream}Signé\ \textbullet\ {\poplogo OnBuch\color{poporange}+}}\par\vspace{2pt}
      {\raggedright\popsemi\fontsize{8}{10.5}\selectfont\color{popline}\MakeUppercase{Équipe pédagogique OnBuch+}\\\MakeUppercase{Conforme au programme officiel MINESEC}\par}
    \end{minipage}\hfill
    {\poplogo\fontsize{22}{22}\selectfont\color{popcream}OnBuch\color{poporange}+}
  \end{tcolorbox}%
}

% ============================================================
% Page de garde
% #1 titre · #2 matière · #3 niveau · #4 module · #5 n° leçon · #6 durée ·
% #7 description / objectifs (texte ou itemize)
% ============================================================
\newcommand{\pagegarde}[7]{%
  \thispagestyle{empty}
  \newgeometry{a4paper,margin=1.7cm,top=1.6cm,bottom=1.4cm}%
  \begin{tikzpicture}[remember picture,overlay]
    \fill[poporange!18] ([xshift=-3.2cm,yshift=-3.6cm]current page.north east) circle (4.6cm);
    \fill[poppurple!14] ([xshift=2.4cm,yshift=7.6cm]current page.south west) circle (3.8cm);
  \end{tikzpicture}%
  {\poplogo\fontsize{30}{30}\selectfont\color{popink}OnBuch\color{poporange}+}\hfill
  \raisebox{0.6ex}{\poppill{Cours signé \textbullet\ Premium}{popink}{popcream}}\par
  \vspace{1pt}\popsquiggle{poppurple}\par
  \vspace{8pt}
  \begin{tcolorbox}[enhanced,colback=poporange,colframe=popink,boxrule=2.8pt,arc=13pt,
      drop shadow=popink,left=18pt,right=18pt,top=12pt,bottom=13pt,after skip=10pt]
    \poppill{#2 \textbullet\ #3}{popink}{popcream}\hspace{5pt}\poppill{#4}{white}{popink}\par\vspace{8pt}
    {\popdisplay\fontsize{14}{14}\selectfont\color{popink}LEÇON #5}\par\vspace{4pt}
    {\raggedright\hyphenpenalty=10000\exhyphenpenalty=10000\popdisplay\fontsize{25}{27}\selectfont\color{popcream}\MakeUppercase{#1}\par}
    \vspace{8pt}
    {\popxbold\fontsize{9.5}{9.5}\selectfont\color{popink}\MakeUppercase{Durée conseillée : #6}}
  \end{tcolorbox}
  \begin{tcolorbox}[enhanced,colback=white,colframe=popink,boxrule=2.2pt,arc=10pt,
      drop shadow=popink,left=16pt,right=16pt,top=10pt,bottom=10pt,after skip=8pt]
    \poppill{Dans cette leçon}{poppurple}{white}\par\vspace{5pt}
    \popsemi\fontsize{9.3}{12.6}\selectfont\color{popink}#7
  \end{tcolorbox}
  \noindent\begin{minipage}[t]{0.487\linewidth}
    \begin{tcolorbox}[enhanced,colback=popgreen,colframe=popink,boxrule=2.2pt,arc=10pt,
        drop shadow=popink,left=11pt,right=11pt,top=10pt,bottom=10pt,height=3.1cm,valign=center]
      \tcbox[colback=white,colframe=popink,boxrule=1.3pt,arc=5pt,boxsep=0pt,nobeforeafter,
        left=3pt,right=3pt,top=3pt,bottom=3pt,valign=center]{\qrcode[height=1.9cm]{https://play.google.com/store/apps/details?id=com.luvvix.onbuch&hl=fr}}%
      \hspace{8pt}\begin{minipage}[c]{0.5\linewidth}
        {\popdisplay\fontsize{11}{12.5}\selectfont\color{white}\MakeUppercase{Télécharge OnBuch}}\par\vspace{4pt}
        {\raggedright\popsemi\fontsize{8.4}{11}\selectfont\color{white}Gratuit \textbullet\ Google Play\\Tuteur IA, annales, concours, résultats\par}
      \end{minipage}
    \end{tcolorbox}
  \end{minipage}\hfill
  \begin{minipage}[t]{0.487\linewidth}
    \begin{tcolorbox}[enhanced,colback=poppurple,colframe=popink,boxrule=2.2pt,arc=10pt,
        drop shadow=popink,left=11pt,right=11pt,top=10pt,bottom=10pt,height=3.1cm,valign=center]
      \tikz[baseline=-0.7ex]\node[circle,fill=white,draw=popink,line width=1.3pt,minimum size=1.6cm,inner sep=0]{\popdisplay\fontsize{24}{24}\selectfont\color{poppurple}+};%
      \hspace{8pt}\begin{minipage}[c]{0.6\linewidth}
        {\popdisplay\fontsize{11}{12.5}\selectfont\color{white}\MakeUppercase{Passe à OnBuch+}}\par\vspace{4pt}
        {\raggedright\popsemi\fontsize{8.4}{11}\selectfont\color{white}Tous les cours signés, Tuteur IA illimité, fiches PDF — onglet OnBuch+ de l'app\par}
      \end{minipage}
    \end{tcolorbox}
  \end{minipage}
  \vspace{8pt}
  \popcontenu
  \vfill
  \signatureonbuch
  \restoregeometry
  \newpage
}

% Rangée « Ce cours contient » (page de garde)
\newcommand{\poptile}[3]{% #1 couleur #2 gros texte #3 légende
  \begin{tcolorbox}[enhanced,colback=#1,colframe=popink,boxrule=2pt,arc=9pt,shadow={2.5pt}{-2.5pt}{0pt}{popink},
      left=6pt,right=6pt,top=9pt,bottom=9pt,halign=center,valign=center,height=1.75cm,width=\linewidth,nobeforeafter]
    {\popdisplay\fontsize{12.5}{13}\selectfont\color{white}\MakeUppercase{#2}}\par\vspace{3pt}
    {\centering\popsemi\fontsize{7.8}{9.5}\selectfont\color{white}#3\par}
  \end{tcolorbox}}
\newcommand{\popcontenu}{%
  \par\noindent{\popxboldls\fontsize{7.5}{7.5}\selectfont\color{popmuted}CE COURS SIGNÉ CONTIENT}\par\vspace{7pt}
  \noindent\begin{minipage}[t]{0.235\linewidth}\poptile{popblue}{Cours}{complet, illustré, pas à pas}\end{minipage}\hfill
  \begin{minipage}[t]{0.235\linewidth}\poptile{poporange}{Méthodes}{et exercices résolus}\end{minipage}\hfill
  \begin{minipage}[t]{0.235\linewidth}\poptile{poppink}{Exercices}{type examen + corrigés}\end{minipage}\hfill
  \begin{minipage}[t]{0.235\linewidth}\poptile{popgreen}{Fiche bilan}{l'essentiel à retenir}\end{minipage}\par}

% Sommaire
\newtcolorbox{sommaire}{enhanced,colback=white,colframe=popink,boxrule=2.2pt,arc=10pt,
  drop shadow=popink,left=18pt,right=18pt,top=13pt,bottom=13pt,before skip=6pt,after skip=20pt,
  before upper={\poppill{Sommaire}{poppurple}{white}\par\vspace{9pt}\popsemi\fontsize{10.6}{14.5}\selectfont\color{popink}}}

% Signature de fin de cours — poussée tout en bas de la dernière page
\newcommand{\findecours}{%
  \par\vfill\nopagebreak
  \begin{center}\popsquiggle{poporange}\end{center}\vspace{4pt}
  \signatureonbuch
}
\AtBeginDocument{\def\llbracket{\mathopen{[\mkern-3.5mu[}}\def\rrbracket{\mathclose{]\mkern-3.5mu]}}} % intervalles d'entiers (glyphe ⟦ absent de la police)
% Glyphes absents de Fira Math : redéfinis à partir de glyphes disponibles
\AtBeginDocument{%
  \def\setminus{\mathbin{\backslash}}\let\smallsetminus\setminus
  \def\vdots{\mathord{\vbox{\baselineskip3.2pt\lineskiplimit0pt\kern4pt\hbox{.}\hbox{.}\hbox{.}}}}%
  \def\ddots{\mathinner{\mkern1mu\raise6.5pt\hbox{.}\mkern2mu\raise3.6pt\hbox{.}\mkern2mu\raise0.7pt\hbox{.}\mkern1mu}}%
  \def\triangle{\mathord{\scalebox{0.9}{$\symup{\Delta}$}}}%
  \def\nabla{\mathord{\raisebox{\depth}{\scalebox{1}[-1]{$\symup{\Delta}$}}}}%
  \def\checkmark{\popcheck{popdark}}%
  \def\star{\mathbin{\ast}}\def\bigstar{\mathord{\ast}}%
  \def\diamond{\mathbin{\scalebox{0.6}{$\Diamond$}}}%
  \def\bigcup{\mathop{\mathchoice{\raisebox{-0.3em}{\scalebox{1.7}{$\cup$}}}{\scalebox{1.25}{$\cup$}}{\cup}{\cup}}\displaylimits}%
  \def\bigcap{\mathop{\mathchoice{\raisebox{-0.3em}{\scalebox{1.7}{$\cap$}}}{\scalebox{1.25}{$\cap$}}{\cap}{\cap}}\displaylimits}%
  \def\lesseqqgtr{\mathrel{\lessgtr}}\def\bigoplus{\mathop{\oplus}\displaylimits}%
}
\providecommand{\ssi}{si et seulement si} % abréviation fréquente
\AtBeginDocument{\providecommand{\wideparen}{\overparen}} % arc de cercle

\begin{document}

\pagegarde{Personnalité et personne}{Philosophie}{1re Tech}{Philosophie – classe de Première technique}{N05}{2 séances (fiche de progression harmonisée)}{Cette leçon examine comment se construit la personnalité humaine à travers ses facteurs biologiques, sociaux et culturels, puis interroge la valeur absolue de la personne à travers les perspectives de Kant (la dignité) et de Mounier (le personnalisme). Elle conduit l'élève à appliquer ces notions à des situations concrètes de la vie camerounaise et à justifier une démarche argumentée.
\vspace{4pt}
\begin{multicols}{2}\small
\begin{itemize}
\item À la fin de cette leçon, je sais définir les notions de personnalité et de personne et les distinguer clairement
\item À la fin de cette leçon, je sais identifier et expliquer les facteurs biologiques, sociaux et culturels d'édification de la personnalité
\item À la fin de cette leçon, je sais analyser un cas concret camerounais (éducation familiale, scolaire, communautaire) pour montrer comment une personnalité se construit
\item À la fin de cette leçon, je sais exposer la conception kantienne de la dignité de la personne (humanité comme fin en soi)
\item À la fin de cette leçon, je sais présenter le personnalisme de Mounier et sa conception de la personne comme être relationnel et engagé
\item À la fin de cette leçon, je sais comparer les perspectives de Kant et de Mounier sur la valeur de la personne
\end{itemize}\end{multicols}}

\begin{sommaire}
\begin{multicols}{2}
\begin{enumerate}
\item Introduction : problématique et définitions
\item Les facteurs biologiques de la personnalité
\item Les facteurs sociaux et culturels de la personnalité
\item La valeur de la personne selon Kant : la dignité
\item La valeur de la personne selon Mounier : le personnalisme
\item Synthèse et méthode : appliquer et argumenter
\item Activité d'intégration
\item Méthodes et exercices résolus
\item Exercices
\item Corrigés des exercices
\item Fiche bilan
\end{enumerate}
\end{multicols}
\end{sommaire}

\begin{objectifs}
\begin{itemize}
\item À la fin de cette leçon, je sais définir les notions de personnalité et de personne et les distinguer clairement
\item À la fin de cette leçon, je sais identifier et expliquer les facteurs biologiques, sociaux et culturels d'édification de la personnalité
\item À la fin de cette leçon, je sais analyser un cas concret camerounais (éducation familiale, scolaire, communautaire) pour montrer comment une personnalité se construit
\item À la fin de cette leçon, je sais exposer la conception kantienne de la dignité de la personne (humanité comme fin en soi)
\item À la fin de cette leçon, je sais présenter le personnalisme de Mounier et sa conception de la personne comme être relationnel et engagé
\item À la fin de cette leçon, je sais comparer les perspectives de Kant et de Mounier sur la valeur de la personne
\item À la fin de cette leçon, je sais appliquer les notions de la leçon à des situations de la vie courante (respect de la dignité, travail, justice, réseaux sociaux)
\item À la fin de cette leçon, je sais rédiger et justifier une démarche argumentée, une réponse ou une conclusion en dissertation philosophique
\end{itemize}
\end{objectifs}

\begin{prerequis}
\begin{itemize}
\item Notion d'individu et de société (vie de classe et éducation civique du secondaire)
\item Distinction entre l'inné et l'acquis (notions de base en sciences de la vie)
\item Notion de valeur morale (bien, mal, devoir) abordée en éducation à la citoyenneté
\item Expérience quotidienne de la vie en communauté (famille, village, école, tontine)
\end{itemize}
\end{prerequis}

\newpage

\coursec{Introduction : problématique et définitions}

\courssub{Situation d'accroche : le cas des jumeaux de Briqueterie}

Au cœur du quartier populaire de Briqueterie, à Yaoundé, vivent deux frères jumeaux, \emph{Paul} et \emph{Pierre}. Ils partagent le même patrimoine génétique, ont grandi sous le même toit, mangé les mêmes plats de \emph{ndolé} et de \emph{koki}, fréquenté le même établissement secondaire — le Lycée de Briqueterie — et bénéficié de l'autorité du même père, chef de chantier respecté, et de la même mère, commerçante au marché Central. Pourtant, à 22 ans, leurs destins divergent radicalement. Paul a obtenu son Baccalauréat technique série F4, gère aujourd'hui une petite entreprise de menuiserie aluminium qui emploie trois apprentis, et est connu dans le quartier pour son sens des responsabilités et sa parole tenue. Pierre, lui, a décroché en classe de Seconde, a sombré dans la consommation de \emph{tramol} et de boissons frelatées, multiplie les vols à l'arraché sur l'axe Carrefour Warda–Briqueterie, et fait l'objet d'un mandat de dépôt à la prison centrale de Kondengui.

Comment deux individus issus de la \textbf{même hérédité} et façonnés par le \textbf{même milieu} construisent-ils des \textbf{personnalités} si radicalement opposées ? Cette situation, loin d'être anecdotique, met en lumière le mystère de l'édification de soi. Elle pose simultanément une seconde question, tout aussi cruciale : si Pierre est jugé coupable et incarcéré, pourquoi le droit camerounais — comme toute législation inspirée des Droits de l'homme — lui garantit-il encore le droit à un avocat, l'interdiction de la torture, le respect de son intégrité physique ? Pourquoi dit-on qu'il reste une \textbf{personne} dotée de \textbf{dignité}, alors même que ses actes semblent l'en exclure ?

\courssub{Problématique}

De ce double constat naissent les deux pôles de notre problématique, qui guideront toute cette leçon :

\begin{enumerate}
    \item \textbf{Qu'est-ce qui fait d'un individu humain une personnalité unique ?} Comment s'articulent l'inné (le biologique) et l'acquis (le social, le culturel, le spirituel) pour produire cette singularité irréductible qu'est Paul ou Pierre ?
    \item \textbf{La personne a-t-elle une valeur absolue ou relative ?} La dignité de Pierre est-elle suspendue à ses mérites, à son utilité sociale, à sa moralité ? Ou bien possède-t-elle un caractère inaliénable, intangible, qui commande le respect quelles que soient les circonstances ?
\end{enumerate}

Pour répondre, nous devons d'abord poser nos concepts avec rigueur : distinguer ce qui relève de l'individu biologique, de la personne morale et juridique, et de la personnalité psychologique.

\courssub{Définition de la personnalité : l'unicité concrète}

\paragraph{Intuition :} Quand nous disons de quelqu'un : « Il a une forte personnalité » ou « Elle n'a pas de personnalité », nous ne parlons pas de son statut civil, mais de son \emph{style d'être au monde} : sa façon de réagir sous la pression, ses goûts, ses valeurs, sa stabilité émotionnelle, son charisme ou sa discrétion. La personnalité, c'est ce qui fait que Paul \emph{est} Paul et non Pierre, même s'ils sont indiscernables biologiquement.

\begin{definition}[Personnalité]
La \cle{personnalité} est l'ensemble organisé et relativement stable des traits \textbf{physiques}, \textbf{psychologiques} (tempérament, intelligence, affectivité, caractère) et \textbf{moraux} (valeurs, convictions, engagements) qui distinguent un individu de tous les autres et assurent sa reconnaissance identitaire dans la durée. Elle est le \emph{contenu singulier} de l'existence.
\end{definition}

\paragraph{Structure :} On distingue classiquement trois strates imbriquées :
\begin{itemize}
    \item Le \textbf{tempérament} : fond biologique, inné, héréditaire (ex. : extraversion/introversion, réactivité nerveuse). C'est le « donné » de départ.
    \item Le \textbf{caractère} : ensemble des habitudes morales et volontaires acquises par l'éducation, l'effort personnel, l'imitation (ex. : honnêteté, courage, lâcheté, générosité). C'est le « construit ».
    \item L'\textbf{intelligence} et l'\textbf{affectivité} : capacités cognitives et émotionnelles qui médiatisent le rapport au monde.
\end{itemize}

\begin{exemplebox}[Personnalité et choix de vie à Douala]
Deux jeunes diplômés en Génie Civil, même école, mêmes notes. L'un, prudent, anxieux, attaché à la sécurité (tempérament inhibé, caractère conformiste), intègre la fonction publique comme ingénieur des Travaux Publics. L'autre, audacieux, créatif, tolérant le risque (tempérament extraverti, caractère entrepreneurial), crée sa propre structure de BTP et bidouille des chantiers complexes. Leurs \emph{personnalités} orientent différemment des compétences techniques identiques.
\end{exemplebox}

\begin{attention}[Piège : personnalité $\neq$ personne]
Ne confondez pas \textbf{personnalité} (contenu psychologique et moral, variable, perfectible) et \textbf{personne} (statut ontologique et juridique, binaire : on l'est ou on ne l'est pas). Un nourrisson de trois jours \emph{est une personne} (il a des droits, on ne peut le tuer), mais il n'a \emph{pas encore de personnalité constituée} (pas de caractère, pas de valeurs choisies). Inversement, un grand criminel a une personnalité (très structurée, quoique perverse) mais sa qualité de personne est ce qui interdit de le torturer.
\end{attention}

\courssub{Définition de la personne : le sujet de droits et de dignité}

\paragraph{Intuition :} Dans le Code civil camerounais (inspiré du Code Napoléon et de la Déclaration de 1789), l'article 1er dispose : « La personnalité juridique commence à la naissance et s'éteint à la mort. » Le terme « personnalité » y est employé au sens juridique de « capacité d'être sujet de droit ». Mais philosophiquement, la \textbf{personne} dépasse le droit : c'est l'être humain élevé au rang de \emph{fin en soi}, qui ne peut être utilisé comme un simple moyen.

\begin{definition}[Personne]
La \cle{personne} est l'être humain considéré comme \textbf{sujet de droits, de devoirs et de dignité}, c'est-à-dire comme une fin en soi, et non comme une chose ou un instrument. Ce statut est \textbf{universel} (tout homme est une personne) et \textbf{inaliénable} (on ne peut le perdre, même par le crime).
\end{definition}

\paragraph{Mécanisme de reconnaissance :} La personne émerge par la \textbf{reconnaissance mutuelle} (Hegel) : je suis une personne parce que les autres me traitent comme telle (ils me parlent, me jugent, me sanctionnent, me soignent) et que je me reconnais moi-même comme responsable de mes actes. C'est un statut \emph{relationnel} et \emph{normatif}.

\begin{exemplebox}[La personne face à la justice à Ngaoundéré]
Un prévenu comparaît pour vol aggravé. Le juge lui demande : « Avez-vous un avocat ? » — « Non. » — « Je vous en commets un d'office. » Pourquoi ? Parce que, \emph{quoi qu'il ait fait}, il reste une \textbf{personne} : il a le droit à la défense, à la présomption d'innocence, à ne pas être frappé. Si on le traitait en « chose » (objet de répression pure), on le jetterait en prison sans procès. Le procès \emph{est} la procédure qui respecte sa qualité de personne.
\end{exemplebox}

\courssub{Distinction fondamentale : individu, personne, personnalité}

Ces trois notions sont souvent confondues dans le langage courant. Elles désignent trois plans de réalité distincts, bien qu'inséparables chez l'être humain concret.

\begin{center}
\begin{tikzpicture}[
    circle1/.style={draw=popblue, thick, fill=popblue!10, minimum width=6cm},
    circle2/.style={draw=poporange, thick, fill=poporange!10, minimum width=9cm},
    circle3/.style={draw=popgreen, thick, fill=popgreen!10, minimum width=12cm},
    labelstyle/.style={font=\small\bfseries, text=popdark},
    examplestyle/.style={font=\footnotesize, text=popmuted, align=center}
]
% Cercle 3 : Personnalité (extérieur)
\node[circle3] (perso) at (0,0) {};
\node[labelstyle] at (0, 5.2) {PERSONNALITÉ};
\node[examplestyle] at (0, 4.2) {Singularité psychique \& morale};
\node[examplestyle] at (0, 3.5) {Paul : entrepreneur, rigoureux, catholique pratiquant};
\node[examplestyle] at (0, 2.8) {Pierre : impulsif, rebelle, en rupture};

% Cercle 2 : Personne (milieu)
\node[circle2] (pers) at (0,0) {};
\node[labelstyle] at (0, 1.8) {PERSONNE};
\node[examplestyle] at (0, 1.0) {Statut moral \& juridique : Sujet de droits/devoirs};
\node[examplestyle] at (0, 0.3) {Paul \& Pierre : tous deux citoyens camerounais,};
\node[examplestyle] at (0, -0.4) {justiciables, patients, électeurs (si majeurs)};

% Cercle 1 : Individu (intérieur)
\node[circle1] (ind) at (0,0) {};
\node[labelstyle] at (0, -1.4) {INDIVIDU};
\node[examplestyle] at (0, -2.2) {Exemplaire biologique de l'espèce \emph{Homo sapiens}};
\node[examplestyle] at (0, -2.9) {Paul \& Pierre : même ADN, même groupe sanguin O+,};
\node[examplestyle] at (0, -3.6) {même morphologie, même prédisposition au diabète};

% Flèches explicatives
\draw[->, popdark, thick] (-6.5, 4) -- (-4.5, 4) node[right, labelstyle, xshift=1cm] {Se construit toute la vie};
\draw[->, popdark, thick] (-6.5, 1) -- (-4.5, 1) node[right, labelstyle, xshift=1cm] {Reconnue à la naissance};
\draw[->, popdark, thick] (-6.5, -2) -- (-4.5, -2) node[right, labelstyle, xshift=1cm] {Donné à la conception};
\end{tikzpicture}
\end{center}
\begin{popfigure}
\legende{Schéma des trois cercles concentriques : l'individu (biologique) est le socle ; la personne (statut moral/juridique) l'enveloppe et le protège ; la personnalité (singularité) est l'accomplissement variable au sommet. Exemples camerounais : Paul et Pierre, jumeaux de Briqueterie.}
\end{popfigure}

\begin{center}
\begin{tabularx}{\linewidth}{|l|X|X|}\hline
\rowcolor{popblueL}
\textbf{Notion} & \textbf{Définition rigoureuse} & \textbf{Exemple camerounais (Paul \& Pierre)} \\ \hline
\textbf{Individu} & Exemplaire comptable de l'espèce humaine ; réalité biologique, génétique, anatomique. « Un » parmi d'autres. & Jumeaux monozygotes : même caryotype 46,XY, même empreintes digitales (quasi), même susceptibilité génétique au paludisme. \\ \hline
\textbf{Personne} & Sujet de droit et de morale ; être doué de \textbf{dignité} (valeur intrinsèque) ; titulaire de droits (vie, intégrité, expression) et de devoirs (respect d'autrui, lois). Statut \textbf{binaire} (on l'est / on ne l'est pas). & Tous deux sont « personnes » aux yeux de la Constitution du 18 janvier 1996 (art. 1er : « La personne humaine est sacrée »). Pierre, en prison, garde le droit à la santé, aux visites, à un procès équitable. \\ \hline
\textbf{Personnalité} & Organisation unique et dynamique de traits physiques, psychiques, moraux ; \textbf{singularité} irréductible ; se \textbf{construit} dans le temps (éducation, choix, épreuves). & Paul : personnalité \emph{intégrée}, \emph{responsable}, \emph{projetante}. Pierre : personnalité \emph{fragmentée}, \emph{réactive}, \emph{en quête de sens} (ou en rupture). \\ \hline
\end{tabularx}
\end{center}
\begin{popfigure}
\legende{Tableau comparatif des trois notions : individu (biologique, identique pour les jumeaux), personne (juridique/moral, identique pour les jumeaux), personnalité (psychologique/moral, différente pour les jumeaux).}
\end{popfigure}

\begin{savaistu}[Droit camerounais et personnalité]
En droit positif camerounais (Code civil, Livre I, Titre I), la \textbf{personnalité juridique} commence à la \textbf{naissance viable} (enfant né vivant et viable) et s'éteint à la \textbf{mort}. Elle confère la \emph{capacité de jouissance} (être titulaire de droits : propriété, héritage, honneur) et, sous réserve de la majorité ou de l'émancipation, la \emph{capacité d'exercice} (agir seul valablement). 

Mais la \textbf{personnalité morale} (au sens philosophique de notre leçon) ne se décrète pas par la loi : elle \emph{se construit} par l'éducation, la culture, les choix éthiques, l'engagement citoyen. C'est pourquoi le programme d'Éducation à la Citoyenneté et à la Morale (ECM) dans les lycées techniques vise explicitement la \emph{formation de la personnalité} de l'élève, bien au-delà de sa simple inscription à l'état civil.
\end{savaistu}

\courssub{Annonce du plan}

Fort de ces distinctions, nous allons désormais déployer l'analyse en deux temps majeurs :

\begin{enumerate}
    \item \textbf{Les facteurs d'édification de la personnalité (Chapitre I).} Nous examinerons comment le biologique (tempérament, hérédité), le social (famille, école, pairs, médias, culture bamiléké, bassa, peule, etc.), le culturel (valeurs, rites, langues) et le spirituel (liberté, conscience, projet de vie) s'entrelacent pour faire de l'individu une personnalité singulière. Cela éclairera le cas Paul/Pierre.
    \item \textbf{La valeur de la personne : deux perspectives philosophiques majeures (Chapitre II).} Nous confronterons \textbf{Emmanuel Kant} (la personne comme \emph{fin en soi}, dignité = prix infini, devoir de respect inconditionnel) et \textbf{Emmanuel Mounier} (le personnalisme : la personne comme \emph{être en relation}, engagement, communion, refus de l'individualisme comme du collectivisme). Ces deux géants nous donneront les armes conceptuelles pour répondre : Pierre a-t-il encore une dignité intouchable ?
\end{enumerate}

\coursec{Les facteurs biologiques de la personnalité}

La section précédente a posé la distinction fondamentale entre la \cle{personnalité} (ensemble dynamique et structuré de traits psychologiques, fruit d'une construction) et la \cle{personne} (sujet de droit et de morale, doté d'une valeur intrinsèque). Nous devons maintenant enquêter sur les fondations mêmes de cet édifice : avant d'être façonnée par la culture, l'éducation ou les choix personnels, la personnalité repose sur un socle biologique. Ce socle ne la \emph{détermine} pas au sens strict (ce serait nier la liberté), mais il la \emph{conditionne} : il offre des potentialités, des seuils de réactivité, une « matière première » que l'existence viendra modeler. Comprendre ce facteur biologique, c'est refuser à la fois l'innéisme naïf (tout est dans les gènes) et le culturalisme radical (le corps n'est qu'une page blanche).

\courssub{L'hérédité : transmission des traits physiques et des dispositions temperamentales}

L'hérédité est le premier facteur biologique. Elle opère par la transmission du patrimoine génétique (l'ADN) lors de la fécondation. Ce patrimoine détermine des \cle{traits physiques} objectifs : la taille, la pigmentation de la peau (teint), la morphologie du visage, la couleur des yeux, la texture des cheveux. Mais au-delà de cette anatomie visible, la génétique transmet des \cle{dispositions} internes, souvent regroupées sous le terme de \cle{tempérament}.

\begin{definition}[Tempérament]
Ensemble des dispositions nerveuses, physiologiques et hormonales \emph{innées} qui constituent la base biologique du caractère. Le tempérament se manifeste par la vitesse, l'intensité et la durée des réactions affectives et motrices (ex. : réactivité forte/faible, impulsivité, stabilité émotionnelle). Il est relativement stable tout au long de la vie, contrairement au caractère qui s'élabore par l'éducation et l'expérience.
\end{definition}

Historiquement, la médecine hippocratique (V\textsuperscript{e} s. av. J.-C.) puis Galien (II\textsuperscript{e} s. ap. J.-C.) ont lié ces dispositions à la prépondérance de l'un des quatre \emph{humeurs} dans l'organisme. Bien que la physiologie moderne ait abandonné cette théorie des humeurs, la typologie qui en découle reste un repère pédagogique utile pour décrire des profils de réactivité de base :

\begin{center}
\begin{tabularx}{\linewidth}{|l|X|X|}\hline
\textbf{Tempérament} & \textbf{Humeur dominante (théorie antique)} & \textbf{Profil de réactivité (lecture moderne)} \\ \hline
\cle{Sanguin} & Sang & Réactivité forte, rapide, courte durée. Extraversion, enthousiasme, instabilité de l'attention. \\ \hline
\cle{Colérique} & Bile jaune & Réactivité forte, rapide, longue durée. Énergie, dominance, irritabilité, persévérance. \\ \hline
\cle{Mélancolique} & Bile noire & Réactivité faible, lente, longue durée. Introversion, sensibilité profonde, tendance à l'inhibition, rigueur. \\ \hline
\cle{Flegmatique} & Lymphe & Réactivité faible, lente, courte durée. Calme, lenteur, régularité, faible expressivité émotionnelle. \\ \hline
\end{tabularx}
\end{center}

\noindent
\textbf{Remarque :} Nul n'est un « pur » type. Tout individu présente un \emph{tempérament dominant} mélangé à des traits secondaires. Le tempérament n'est pas le caractère : il en est la \emph{condition de possibilité}. Un colérique peut devenir un leader juste (caractère vertueux) ou un tyran (caractère vicieux) selon son éducation et ses choix.

\courssub{Le rôle du corps : santé, handicap, apparence et construction de l'image de soi}

Le corps n'est pas seulement un support génétique passif ; c'est le \emph{véhicule} de notre présence au monde. La \cle{phénoménologie} (Merleau-Ponty) insiste sur le « corps propre » (\emph{Leib} vs \emph{Körper}) : je ne \emph{possède} pas un corps, je \emph{suis} mon corps. Trois dimensions illustrent l'influence du corps sur la personnalité :

\begin{enumerate}
    \item \textbf{La santé et la vitalité :} Un capital santé solide favorise l'exploration, la prise de risque, la confiance en ses capacités. À l'inverse, la maladie chronique ou la fatigue permanente peuvent orienter la personnalité vers le repli, la prudence excessive, voire l'anxiété.
    \item \textbf{Le handicap :} Une déficience sensorielle (cécité, surdité), motrice ou cognitive restructure l'espace de l'expérience. La personnalité ne se réduit pas au handicap, mais ce dernier impose des adaptations cognitives et sociales spécifiques qui sculptent l'identité (résilience, créativité dans la compensation, ou au contraire sentiment d'impuissance appris si l'entourage est surprotecteur ou excluant).
    \item \textbf{L'apparence et le regard d'autrui :} La taille, la beauté (selon les canons culturels), les stigmates visibles (cicatrices, particularités génétiques) déclenchent des réactions sociales (admiration, moquerie, pitié, rejet) que l'individu intériorise. C'est le miroir social (Cooley, Mead) : je deviens ce que je crois que les autres voient en moi.
\end{enumerate}

\begin{exemplebox}[Exemple camerounais : un enfant albinos à Douala]
Imaginons \cle{Paul}, 10 ans, né albinos (absence de mélanine) dans un quartier populaire de Douala. Biologiquement, son patrimoine génétique (mutation sur le gène \emph{TYR} ou \emph{OCA2}) détermine sa peau très claire, ses cheveux blonds, sa photophobie et son risque accru de cancers cutanés. Ces \emph{traits physiques} sont des données brutes.
Mais la \emph{personnalité} de Paul se construit dans l'interaction :
\begin{itemize}
    \item \textbf{Le regard social :} Dans certains milieux, des croyances traditionnelles (sorcellerie, « enfant de la lune », porte-bonheur ou malédiction) génèrent soit une surprotection mystique, soit un rejet violent, des moqueries (\emph{« mbia », « blanc »}) ou des tentatives d'enlèvement pour rituels.
    \item \textbf{L'intériorisation :} Si Paul intègre le rejet, il développera une personnalité anxieuse, méfiante, à faible estime de soi (repli mélancolique). Si sa famille et l'école le valorisent, lui expliquent sa condition biologiquement, et lui donnent des moyens de se protéger (crèmes, chapeaux, lunettes), il pourra développer une personnalité résiliente, affirmée, transformant la différence en force (ex. : militant pour les droits des albinos).
    \item \textbf{Le corps comme contrainte et ressource :} Sa photophobie l'oblige à éviter le soleil de midi ; il organise différemment son temps, ses loisirs (lecture, jeux d'intérieur, musique), ce qui oriente ses centres d'intérêt et ses compétences.
\end{itemize}
\emph{Conclusion :} L'hérédité a donné les \emph{cartes} (albinisme) ; le milieu social et les choix de Paul (et de son entourage) jouent la \emph{partie}. La personnalité émerge de cette dialectique.
\end{exemplebox}

\courssub{Les limites du déterminisme biologique : leçons des études sur les jumeaux}

L'argument massue du déterminisme biologique (« tout est dans les gènes ») repose souvent sur la ressemblance frappante des jumeaux \cle{monozygotes} (vrais jumeaux, même ADN). Pourtant, l'étude rigoureuse de jumeaux élevés \emph{séparément} (méthode des \emph{adoption studies}) révèle la part irréductible de l'environnement et de l'épigénétique.

\begin{exoresolu}[Étude de document : Jumeaux homozygotes élevés séparément — Données comparées]
\textbf{Énoncé :} Une méta-analyse fictive mais réaliste (inspirée des études de Bouchard, Minnesota / Plomin, UK) compare le coefficient de corrélation ($r$) pour divers traits de personnalité entre deux groupes :
\begin{itemize}
    \item Groupe A : Jumeaux monozygotes élevés \emph{ensemble} (mêmes gènes + même famille + même école + même quartier).
    \item Groupe B : Jumeaux monozygotes élevés \emph{séparément} dès la naissance (mêmes gènes + familles, écoles, quartiers différents).
\end{itemize}
Les résultats (corrélation moyenne pondérée) sont donnés dans le tableau ci-dessous. Un $r = 1$ signifie identité parfaite ; $r = 0$ signifie absence de lien ; $r = 0,50$ signifie que 50\% de la variance est partagée.

\begin{center}
\begin{tabular}{|l|c|c|c|}\hline
\textbf{Trait de personnalité (Big Five)} & \textbf{Groupe A (Ensemble)} & \textbf{Groupe B (Séparés)} & \textbf{Écart (A - B)} \\ \hline
Extraversion & 0,52 & 0,44 & 0,08 \\ \hline
Agréabilité & 0,45 & 0,32 & 0,13 \\ \hline
Conscience (Rigueur) & 0,49 & 0,38 & 0,11 \\ \hline
Stabilité émotionnelle (inverse Névrosisme) & 0,41 & 0,35 & 0,06 \\ \hline
Ouverture à l'expérience & 0,48 & 0,42 & 0,06 \\ \hline
\textbf{Moyenne globale} & \textbf{0,47} & \textbf{0,38} & \textbf{0,09} \\ \hline
\end{tabular}
\end{center}

\begin{popfigure}
\begin{tikzpicture}
\begin{axis}[
    ybar,
    bar width=12pt,
    width=\linewidth,
    height=9cm,
    enlargelimits=0.25,
    ylabel={Coefficient de corrélation ($r$)},
    symbolic x coords={Extraversion, Agréabilité, Conscience, Stabilité émotionnelle, Ouverture, Moyenne globale},
    xtick=data,
    nodes near coords={\pgfmathprintnumber[fixed,precision=2]{\pgfplotspointmeta}},
    nodes near coords align={vertical},
    ymin=0, ymax=0.65,
    ytick={0,0.1,0.2,0.3,0.4,0.5,0.6},
    legend style={at={(0.5,-0.25)},anchor=north,legend columns=-1},
    title style={align=center},
    title={Ressemblance de personnalité : Jumeaux MZ élevés ensemble vs séparés},
    popblue!80!popdark,
    grid=major,
]
\addplot[fill=popblue!60, draw=popblue!80!popdark] coordinates {
    (Extraversion,0.52) (Agréabilité,0.45) (Conscience,0.49) (Stabilité émotionnelle,0.41) (Ouverture,0.48) (Moyenne globale,0.47)
};
\addplot[fill=poporange!60, draw=poporange!80!popdark] coordinates {
    (Extraversion,0.44) (Agréabilité,0.32) (Conscience,0.38) (Stabilité émotionnelle,0.35) (Ouverture,0.42) (Moyenne globale,0.38)
};
\legend{Élevés ensemble, Élevés séparément}
\end{axis}
\end{tikzpicture}
\legende{Figure 1 : Degré de ressemblance (corrélation $r$) pour les traits du Big Five chez des jumeaux monozygotes. Même avec 100\% d'ADN commun, l'éducation séparée réduit la ressemblance, prouvant l'impact majeur de l'environnement non partagé. Données fictives réalistes basées sur la littérature comportementale.}
\end{popfigure}

\textbf{Analyse et interprétation :}
\begin{enumerate}
    \item \textbf{L'hérédité pèse lourd :} Même élevés séparément, les jumeaux se ressemblent significativement ($r \approx 0,38$ en moyenne). Cela confirme que le \cle{tempérament} (base biologique) structure fortement la personnalité. Les gènes créent des \emph{dispositions} (seuils de réactivité, sensibilité à la récompense, réactivité au stress).
    \item \textbf{L'environnement non partagé est décisif :} L'écart moyen de $0,09$ (et jusqu'à $0,13$ pour l'Agréabilité) montre que grandir dans des familles différentes, avec des pairs, des enseignants, des traumatismes ou des opportunités distincts, creuse des différences réelles. L'environnement \emph{non partagé} (expériences uniques de chaque jumeau) explique une part massive de la variance (environ 50\% au total, dont une partie mesurée ici par l'écart A-B).
    \item \textbf{L'interaction Gène $\times$ Environnement :} Un même génotype (ex. : sensibilité au stress élevée) s'exprime différemment selon le milieu : dans un milieu sécurisant $\to$ vigilance adaptative ; dans un milieu violent $\to$ trouble anxieux. C'est l'\cle{épigénétique} : l'environnement « allume » ou « éteint » certains gènes.
\end{enumerate}
\end{exoresolu}

\begin{aretenir}[Exemple chiffré : la part du biologique]
Les études comportementales convergent : l'héritabilité ($h^2$) des grands traits de personnalité (Big Five) se situe généralement entre \textbf{40\% et 50\%}. Cela signifie que \emph{environ la moitié} des différences observées entre individus dans une population donnée est corrélée à des différences génétiques. L'autre moitié (environ 50\%) relève de l'environnement (partagé : famille, culture ; et surtout \emph{non partagé} : expériences uniques, hasard développemental, choix personnels). \textbf{La biologie donne des dispositions, non un destin.}
\end{aretenir}

\courssub{Schéma synthétique : de l'hérédité à la personnalité}

Le diagramme ci-dessous résume l'articulation entre le biologique (nécessaire mais non suffisant) et le construit (éducation, culture, choix).

\begin{popfigure}
\begin{tikzpicture}[
    node distance=1.2cm and 1.5cm,
    box/.style={draw=popblue!80!popdark, thick, rounded corners, fill=popblue!10, minimum width=3.2cm, minimum height=1.2cm, align=center, font=\small},
    arrow/.style={-Stealth, thick, popblue!80!popdark},
    plus/.style={font=\Large\bfseries, poporange!80!popdark}
]
% Nœuds
\node[box, align=center] (hered){\textbf{HÉRÉDITÉ}\\(Gènes, ADN)\\\emph{Donné de fait}};
\node[box, right=of hered, align=center] (dispo){\textbf{DISPOSITIONS}\\(Tempérament, Métabolisme,\\Réactivité nerveuse)\\\emph{Potentiels innés}};
\node[box, below right=of dispo, align=center] (educ){\textbf{ÉDUCATION}\\(Famille, École, Pairs)\\\emph{Apports culturels}};
\node[box, below left=of dispo, align=center] (milieu){\textbf{MILIEU DE VIE}\\(Société, Culture, Événements,\\Hasard)\\\emph{Contraintes \& Opportunités}};
\node[box, right=of dispo, xshift=1.5cm, align=center] (choix){\textbf{CHOIX PERSONNELS}\\(Liberté, Réflexivité,\\Engagements)\\\emph{Agency}};
\node[box, below=of dispo, yshift=-0.5cm, fill=popgreen!15, draw=popgreen!80!popdark, minimum width=5cm, align=center] (perso){\textbf{PERSONNALITÉ}\\(Structure unique, stable,\\évolutive)\\\emph{Résultat singulier}};

% Flèches
\draw[arrow] (hered) -- node[above, font=\tiny, popdark] {Transmission} (dispo);
\draw[arrow] (dispo) -- node[above right, font=\tiny, popdark] {Base} (perso);
\draw[arrow] (educ) -- node[left, font=\tiny, popdark] {Façonne} (perso);
\draw[arrow] (milieu) -- node[right, font=\tiny, popdark] {Sculpte} (perso);
\draw[arrow] (choix) -- node[above, font=\tiny, popdark] {Oriente} (perso);

% Symboles "+" pour montrer la convergence
\node[plus] at ($(dispo)!0.5!(educ)$) {+};
\node[plus] at ($(dispo)!0.5!(milieu)$) {+};
\node[plus] at ($(dispo)!0.5!(choix)$) {+};
\end{tikzpicture}
\legende{Figure 2 : Modèle interactionniste. L'hérédité fournit les dispositions (tempérament) ; l'éducation, le milieu et les choix personnels s'y ajoutent pour construire la personnalité unique. Aucune flèche unique ne va de l'hérédité à la personnalité sans médiation.}
\end{popfigure}

\begin{attention}[Piège à éviter : le réductionnisme génétique]
\emph{Erreur fréquente :} « Puisque les jumeaux monozygotes se ressemblent à 50\%, la personnalité est à 50\% génétique, donc immuable. »
\emph{Correction :}
\begin{itemize}
    \item L'héritabilité ($h^2$) est une statistique de \emph{population}, pas une prévision \emph{individuelle}. Elle ne s'applique pas à Paul (notre albinos) pris isolément.
    \item « Disposition » $\neq$ « Détermination ». Une disposition (ex. : impulsivité colérique) peut être canalisée (sport, art, méditation) ou exacerbée (milieu violent). Le \emph{phénotype} (personnalité observable) résulte d'une interaction dynamique Gène $\times$ Environnement tout au long de la vie (plasticité cérébrale).
    \item Dire « c'est génétique » pour excuser un comportement (violence, paresse, addiction) est une \emph{faute morale} (mauvaise foi sartrienne) et une \emph{erreur scientifique} (négligence de l'épigénétique et de la neuroplasticité).
\end{itemize}
\end{attention}

\begin{savaistu}[Le masque et la personne]
Le mot \textbf{personne} vient du latin \emph{persona}, qui désignait le \textbf{masque} porté par les acteurs au théâtre antique. Ce masque amplifiait la voix (\emph{per-sonare} : « résonner à travers ») et indiquait le rôle (le vieillard, le jeune premier, la servante).
\emph{Double sens philosophique :}
\begin{enumerate}
    \item \textbf{Le rôle social :} La personnalité est aussi une « persona » : nous jouons des rôles (élève, fils, ami, citoyen) attendus par la société. Le tempérament biologique facilite ou rend plus difficile certains rôles.
    \item \textbf{Le sujet de droit :} En droit romain, la \emph{persona} est le sujet capable de droits et de devoirs (par opposition à l'esclave, \emph{res}, chose). C'est là la racine de la \cle{dignité} (Kant) et de la \cle{valeur de la personne} (Mounier) que nous étudierons dans les sections 4 et 5.
\end{enumerate}
La biologie forge le visage ; la liberté choisit le masque… ou ose le retirer.
\end{savaistu}

\coursec{Les facteurs sociaux et culturels de la personnalité}

Après avoir examiné les \cle{facteurs biologiques} (tempérament, hérédité, maturation nerveuse) qui constituent le \emph{substrat naturel} de la personnalité, nous devons maintenant explorer comment ce substrat est modelé, orienté et enrichi par l'environnement humain. La personnalité ne se déploie pas dans le vide : elle se construit dans et par la \cle{socialisation}, ce processus continu par lequel l'individu intériorise les normes, les valeurs, les langages et les pratiques de son groupe d'appartenance. Au Cameroun comme ailleurs, c'est au carrefour de la famille, de l'école, des pairs, de la culture, du milieu professionnel et des choix personnels que s'élabore l'identité unique de chaque élève de Première technique.

\begin{definition}[Socialisation]
La \textbf{socialisation} est le processus par lequel l'individu, dès la naissance et tout au long de sa vie, intériorise les normes, les valeurs, les croyances, les langages et les modes de comportement de la société ou du groupe auquel il appartient. Elle transforme un être biologique en un \cle{sujet social} capable d'interagir, de communiquer et de se reconnaître comme membre d'une culture. On distingue la \emph{socialisation primaire} (famille, petite enfance) et la \emph{socialisation secondaire} (école, travail, pairs, médias).
\end{definition}

\courssub{La famille : premier milieu d'éducation et de transmission}

La famille est le \textbf{premier agent de socialisation}. C'est le creuset où se forge le \cle{soi} naissant. Bien avant l'entrée à l'école, l'enfant a déjà intégré une vision du monde, une langue maternelle (ou un répertoire langagier), des interdits (tabous alimentaires, respect des aînés, gestion du regard), des rituels quotidiens et une place dans la lignée.

\begin{exemplebox}[Éducation traditionnelle dans trois aires culturelles camerounaises]
\begin{itemize}
\item \textbf{Zone Bamiléké (Ouest) :} L'éducation est souvent collective, confiée à la \emph{chefferie} et aux \emph{sociétés secrètes} (ex. \emph{Kwifo}, \emph{Ngumba}). L'enfant apprend le sens de l'effort, l'entrepreneuriat, le respect de la hiérarchie et le culte des ancêtres. Le \emph{nom} donné à la naissance (ex. \emph{Tchinda} « celui qui revient », \emph{Ngouafack} « la force du rocher ») programme déjà une attente identitaire.
\item \textbf{Zone Béti (Centre, Sud, Est) :} L'éducation passe par l'\emph{esprit de clan} (\emph{ayong}), le respect scrupuleux des aînés (\emph{nkùm}), l'apprentissage des proverbes (\emph{mvet}) et des danses initiatiques. La \emph{case des palabres} est le lieu où l'enfant observe la résolution des conflits et l'exercice de l'autorité légitime.
\item \textbf{Zone Peule (Nord, Adamaoua) :} Le \emph{pulaaku} (code de conduite peul) structure la personnalité : \emph{semteende} (pudeur, réserve), \emph{munyal} (patience, endurance), \emph{hakkille} (prudence, sagesse), \emph{ngorgu} (courage). L'apprentissage du pastoralisme et la transmission orale (contes, poésie) forgent une identité forte, mobile et communautaire.
\end{itemize}
\end{exemplebox}

Dans tous les cas, la famille transmet un \textbf{capital culturel} (langue, savoir-faire, normes) et un \textbf{capital symbolique} (nom, réputation, honneur) qui constituent le socle de la personnalité. La structure familiale (monoparentale, élargie, polygame, nucléaire) module aussi l'apprentissage de l'autonomie, de la solidarité ou de la rivalité fraternelle.

\courssub{L'école : socialisation secondaire, règles, savoirs et citoyenneté}

L'école technique au Cameroun opère une \textbf{socialisation secondaire} systématique. Elle ne se contente pas d'instruire ; elle \emph{formate} au rôle d'élève, puis de futur technicien ou professionnel.

\begin{itemize}
\item \textbf{Apprentissage des règles impersonnelles :} Ponctualité, port de l'uniforme, respect du règlement intérieur, évaluation par notes anonymes. L'élève découvre l'universalité de la norme : la même règle s'applique à tous, indépendamment de l'origine sociale.
\item \textbf{Acquisition de savoirs codifiés :} Mathématiques, physique, technologie, mais aussi français, anglais, éducation à la citoyenneté. Ces savoirs offrent des \cle{outils cognitifs} pour comprendre le monde et s'y insérer.
\item \textbf{Construction de la citoyenneté :} Levée des couleurs, chant de l'hymne national, cours d'éducation civique, participation aux clubs (environnement, droits de l'homme, entrepreneuriat). L'école vise à faire émerger un \emph{citoyen camerounais} conscient de ses droits et devoirs, au-delà de son appartenance ethnique ou religieuse.
\item \textbf{Socialisation professionnelle :} Les stages, les travaux pratiques en atelier, les visites d'entreprises initient aux codes du monde du travail : hiérarchie, sécurité, travail en équipe, rapport au temps productif.
\end{itemize}

\begin{attention}[Piège : croire que l'école « efface » la famille]
L'école ne remplace pas la famille ; elle \emph{superpose} un nouveau registre de normes. Les conflits de loyauté sont fréquents : l'élève peut vivre une \cle{dissonance cognitive} entre la valorisation de l'oralité et du respect silencieux de l'aîné (famille) et l'exigence de prise de parole argumentée, critique et individuelle (école). La personnalité se construit \emph{dans} cette tension, pas en dehors.
\end{attention}

\courssub{Le groupe des pairs et les réseaux sociaux : l'influence des modèles contemporains}

À l'adolescence, le \textbf{groupe des pairs} devient un agent de socialisation majeur, souvent concurrent de la famille et de l'école. C'est le laboratoire de l'\cle{identité juvénile} : on y teste des styles vestimentaires, des langages (argot, \emph{camfranglais}), des valeurs (solidarité, virilité, modernité), des limites (transgression, prise de risque).

\begin{itemize}
\item \textbf{Influence des amis :} Le besoin d'appartenance pousse à la conformité (effet de mode, consommation d'alcool, absentéisme) mais aussi à l'émulation positive (révision en groupe, projets associatifs, sport).
\item \textbf{Modèles médiatiques et influenceurs :} Artistes urbains (Stanley Enow, Locko, Mimie, Ko-C), footballeurs (André Onana, Vincent Aboubakar), youtubeurs, tiktokeurs camerounais ou de la diaspora. Ils incarnent la \emph{réussite}, la \emph{modernité}, la \emph{liberté}. L'adolescent construit son \emph{idéal du moi} en partie sur ces figures.
\item \textbf{Réseaux sociaux (WhatsApp, Facebook, TikTok, Instagram) :} Ils démultiplient l'espace des pairs. La \cle{reputation numérique} (likes, vues, commentaires) devient un enjeu d'estime de soi. Le cyberharcèlement, l'exposition à des contenus violents ou pornographiques, la pression de l'image parfaite sont de nouveaux risques pour la construction identitaire.
\end{itemize}

\begin{exemplebox}[Enquête de classe — Exemple chiffré]
Dans une classe de 30 élèves de Première technique à Yaoundé, on a demandé : « \emph{Quelles sont les trois influences qui te semblent les plus fortes sur ta façon de penser et d'agir aujourd'hui ?} »
\begin{center}
\begin{tabularx}{\linewidth}{|l|X|c|}\hline
\textbf{Rang} & \textbf{Facteur cité} & \textbf{Nombre de citations (sur 30)} \\ \hline
1 & Famille (parents, oncles/tantes, grands-parents) & 27 \\ \hline
2 & Amis / Groupe de pairs (réel + réseaux sociaux) & 22 \\ \hline
3 & École (professeurs, règlement, cours) & 18 \\ \hline
4 & Médias / Influenceurs / Artistes & 12 \\ \hline
5 & Religion / Église / Mosquée & 9 \\ \hline
6 & Choix personnels / Réflexion propre & 7 \\ \hline
\end{tabularx}
\end{center}
\legende{Enquête informelle, non scientifique, mais illustratrice de la hiérarchie perçue par les élèves. La famille reste la référence majeure, mais le duo « Amis + Réseaux » talonne l'école.}
\end{exemplebox}

\courssub{La culture et la religion : rites d'initiation, valeurs communautaires, appartenance religieuse}

La \cle{culture} au sens large (coutumes, langues, arts, rites, visions du monde) et la \cle{religion} fournissent des \textbf{grands récits} qui donnent du sens à l'existence et structurent la personnalité en profondeur.

\begin{itemize}
\item \textbf{Rites d'initiation :} Circoncision collective (Bassa, Bamiléké, Peuls), \emph{Ngondo} (Sawa), \emph{So} (Béti), \emph{Mvett} (Fang), retraite monastique chrétienne, mémorisation du Coran. Ces rites marquent le passage de l'enfance à l'âge adulte, confèrent un nouveau statut, transmettent un savoir secret et lient l'individu à la lignée des ancêtres ou à la communauté des croyants. Ils sont des \emph{ancrages identitaires} puissants.
\item \textbf{Valeurs communautaires :} L'\emph{ubuntu} (« Je suis parce que nous sommes »), la solidarité de clan (\emph{tontine}, \emph{njangi}), le respect des aînés, l'hospitalité, la honte comme régulateur social (\emph{« Cela ne se fait pas chez nous »}). La personnalité se définit \emph{par rapport aux autres}, non en isolation.
\item \textbf{Appartenance religieuse :} Catholicisme, protestantisme, islam, églises de réveil, religions traditionnelles. La foi offre un cadre éthique (dix commandements, piliers de l'islam), une espérance (salut, paradis), une communauté de soutien (paroisse, oumma) et des rituels quotidiens (prière, jeûne) qui disciplinent le corps et l'esprit. Pour beaucoup d'élèves, Dieu est le \emph{premier référent} de leur identité : « Je suis enfant de Dieu » ou « Je suis serviteur d'Allah ».
\end{itemize}

\begin{savaistu}[Le nom, première inscription identitaire au Cameroun]
Dans de nombreuses cultures camerounaises, le \textbf{nom attribué à l'enfant} n'est pas une simple étiquette : il exprime une \emph{personnalité attendue}, un \emph{destin symbolique} ou les circonstances de la naissance.
\begin{itemize}
\item \emph{Bamiléké :} \textbf{Ngouanet} (« Dieu m'a donné »), \textbf{Tchouakeu} (« Celui qui sait attendre »), \textbf{Kamdem} (« Le chef est arrivé »).
\item \emph{Béti :} \textbf{Ondoa} (« Le chemin »), \textbf{Mvondo} (« La pluie »), \textbf{Essomba} (« La force »).
\item \emph{Peul :} \textbf{Amadou} (« Très loué »), \textbf{Hapsatou} (« Celle qui juge équitablement »), \textbf{Boubacar} (« Père de Bakar »).
\item \emph{Chrétiens :} \textbf{Dieudonné}, \textbf{Bénédicte}, \textbf{Grâce}, \textbf{Prospérité}.
\end{itemize}
Porter ce nom, c'est porter une \emph{exigence} et une \emph{promesse}. L'élève de Première technique porte souvent en lui cette injonction symbolique : « Deviens ce que ton nom dit que tu es. »
\end{savaistu}

\courssub{Le milieu professionnel et économique : le métier, le quartier, les conditions de vie}

L'environnement matériel et professionnel agit comme un \textbf{moule silencieux} sur la personnalité. Ce n'est pas seulement le métier futur (technicien, ouvrier, ingénieur) qui compte, mais le \emph{milieu de vie présent}.

\begin{itemize}
\item \textbf{Le quartier (village, cité universitaire, quartier populaire, résidence huppée) :} L'espace physique détermine les rencontres, les dangers, les ressources, les modèles visibles. Un élève grandissant à \emph{Nkolbisson} (Yaoundé) ou \emph{Bépanda} (Douala) développe des stratégies de survie, de débrouillardise (\emph{Article 15}), de vigilance, différentes de celui qui grandit à \emph{Bastos} ou \emph{Akwa} résidentiel.
\item \textbf{Les conditions de vie :} Accès à l'eau, à l'électricité, à internet, à la santé, aux transports. La précarité économique peut forge la résilience, l'endurance, le sens du sacrifice — ou générer du fatalisme, de la rancœur, de l'agressivité.
\item \textbf{L'apprentissage informel / stage / petit métier :} Beaucoup d'élèves techniques travaillent les vacances (menuiserie, soudure, coiffure, commerce, cybercafé). Ils y apprennent la \emph{réalité du terrain} : relation client, gestion de l'argent, respect du patron, fierté du travail bien fait. Cette expérience module leur rapport à l'effort, à l'autorité, à l'argent.
\end{itemize}

\begin{attention}[Erreur fréquente : opposer inné et acquis comme s'ils s'excluaient]
\textbf{La personnalité ne naît pas de l'addition « biologique + social », mais de leur \emph{interaction permanente}.}
\begin{itemize}
\item Un tempérament \emph{extraverti} (biologique) s'épanouira différemment dans une famille valorisant la parole publique (Bamiléké) ou dans une famille valorisant la réserve (Peul).
\item Une prédisposition génétique à l'anxiété peut être amplifiée par un milieu insécurisant ou atténuée par un entourage sécurisant et une foi solide.
\item \textbf{Conclusion :} Il n'y a pas de « part d'inné » et de « part d'acquis » séparables. Il y a un \emph{devenir} singulier où le biologique offre des potentialités que le social actualise, oriente, inhibe ou sublime.
\end{itemize}
\end{attention}

\courssub{La liberté personnelle : la personnalité comme projet et non comme simple produit}

Si la personnalité était \emph{seulement} le résultat des facteurs précédents, l'homme serait un \emph{produit déterminé} de son hérédité et de son milieu. Or, l'expérience vécue et la philosophie (Sartre, Mounier, Kant) nous enseignent que l'être humain est \textbf{liberté}.

\begin{propriete}[La personnalité comme projet (Mounier, Sartre)]
La personnalité n'est pas un \emph{état} figé, mais un \emph{mouvement} : elle est ce \textbf{projet} par lequel le sujet s'approprie son héritage biologique et social, le trie, l'interprète, le refuse ou le transcende.
\begin{itemize}
\item \textbf{Appropriation :} Je ne subis pas ma timidité ; je \emph{choisis} de la travailler en m'inscrivant au club de théâtre.
\item \textbf{Interprétation :} Mon échec au Probatoire n'est pas une fatalité ; je le \emph{lis} comme un signal pour changer de méthode.
\item \textbf{Transcendance :} Né dans un quartier difficile, je \emph{décide} de devenir ingénieur pour transformer mon milieu.
\end{itemize}
\emph{Exigible au Bac :} Savoir expliquer que la liberté ne nie pas les conditionnements, mais qu'elle est la capacité de \emph{prendre de la distance} par rapport à eux pour poser un acte signifiant.
\end{propriete}

\begin{methode}[Analyser un cas concret : identifier les facteurs en interaction]
Pour comprendre la personnalité d'un individu réel ou fictif, suivez ces 4 étapes :
\begin{enumerate}
\item \textbf{Repérer les facteurs biologiques :} Tempérament apparent (calme, nerveux, leader né), santé, aptitudes physiques/intellectuelles innées, hérédité connue.
\item \textbf{Repérer les facteurs sociaux et culturels :} Famille (structure, valeurs, style éducatif), école (parcours, réussites/échecs, orientation), pairs (groupe d'appartenance, influence), culture/religion (rites, croyances, langue), milieu pro/éco (quartier, travail, ressources).
\item \textbf{Repérer les choix personnels (liberté) :} Décisions marquantes (orientation, rupture, engagement, conversion, effort soutenu), projets d'avenir, discours de justification (« Pourquoi j'ai fait ça ? »).
\item \textbf{Conclure sur l'interaction :} Montrez comment le biologique a été \emph{saisi} par le social, et comment la liberté a \emph{répondu} à cette double détermination. Évitez le déterminisme (tout est joué d'avance) comme le volontarisme naïf (je peux tout, n'importe quand).
\end{enumerate}
\end{methode}

\courssub{Étude de cas guidée : portrait croisé de deux élèves de Première technique}

Appliquons la méthode à deux profils types, inspirés de la réalité camerounaise.

\begin{exoresolu}[Portrait croisé : Amadou (village) vs Stéphanie (ville)]
\textbf{Énoncé :} Amadou, 17 ans, Première F (Mécanique) à Ngaoundéré. Stéphanie, 17 ans, Première G (Génie Civil) à Douala. Analysez les facteurs qui façonnent leur personnalité.

\textbf{Solution détaillée :}

\begin{center}
\begin{tabularx}{\linewidth}{|l|X|X|}\hline
\textbf{Facteur} & \textbf{Amadou (Ngaoundéré, milieu peul rural/urbain)} & \textbf{Stéphanie (Douala, milieu sawa/urbain cosmopolite)} \\ \hline
\textbf{Biologique} & Tempérament \emph{calme, endurant} (élevage), bonne constitution physique. & Tempérament \emph{vif, nerveux}, myopie corrigée, douée en dessin technique. \\ \hline
\textbf{Famille} & Père éleveur, mère au foyer, 3 épouses, 12 enfants. Éducation \emph{pulaaku} : réserve, respect absolu des aînés, apprentissage par observation. Frères aînés = modèles. & Parents fonctionnaires (mère infirmière, père comptable), 2 enfants. Éducation \emph{dialoguée}, valorisation de l'école, autonomie encouragée, prière familiale catholique. \\ \hline
\textbf{École} & Collège public, effectifs pléthoriques (80/classe), peu de matériel. Professeurs exigeants mais distants. Orientation technique par défaut (notes moyennes). Fierté d'être « technicien ». & Lycée privé bilingue, effectifs réduits (35), ateliers équipés. Orientation choisie (Génie Civil). Pression pour l'excellence (mention au Probatoire visée). \\ \hline
\textbf{Pairs / Réseaux} & Groupe d'amis du quartier : foot, thé, discussions sur l'émigration (Europe/USA), WhatsApp pour suivre actualités. Modèles : footballeurs locaux, grands frères revenus du « bush ». & Groupe mixte (filles/garçons) : révisions, sorties centre commercial, TikTok/Insta. Influenceuses mode/lifestyle (camerounaises + ivoiriennes). Modèles : ingénieures femmes, architectes. \\ \hline
\textbf{Culture / Religion} & Islam (sunnite), 5 prières/jour, Coran le week-end, Ramadan strict. Rite de circoncision passé à 12 ans (entrée dans l'âge adulte). Solidarité clanique (\emph{njangi}). & Catholicisme (paroisse dynamique), aumônerie jeunes, scoutisme. Ngondo suivi à la TV. Culture \emph{camfranglais} fluide. Identité « Doualaise » fière. \\ \hline
\textbf{Milieu pro / éco} & Aide au troupeau vacances. Connaissance du bétail, réparation motos. Quartier populaire, eau au forage, électricité irrégulière. Débrouillardise (\emph{Article 15}) quotidienne. & Stages en BTP (entreprise père ami). Visite chantiers (Pont Wouri, tours). Quartier résidentiel, fibre optique, climatisation. Gestion budget de poche. \\ \hline
\textbf{Choix personnels (Liberté)} & \textbf{Projet :} « Devenir chef d'atelier maintenance engins agricoles, moderniser l'élevage paternel. » Choix : refuser l'exode rural, valoriser le savoir local + technique école. & \textbf{Projet :} « Intégrer ENSTP Yaoundé, puis Master Urbanisme durable. » Choix : s'engager dans association « Filles du BTP », refuser mariage précoce proposé par oncle. \\ \hline
\textbf{Synthèse interaction} & Personnalité \emph{ancrée, collective, patiente, pragmatique}. La technique est un \emph{outil au service de la communauté}. Liberté = fidélité créative. & Personnalité \emph{ouverte, individuelle, ambitieuse, critique}. La technique est un \emph{levier d'ascension et de transformation}. Liberté = rupture assumée. \\ \hline
\end{tabularx}
\end{center}
\legende{Analyse comparative montrant comment des facteurs similaires (famille, école, religion) produisent des personnalités distinctes par leur \emph{interaction} avec le biologique et les \emph{choix personnels}.}
\end{exoresolu}

\begin{aretenir}[Synthèse : Les facteurs sociaux et culturels]
\begin{itemize}
\item La \textbf{socialisation} (primaire : famille ; secondaire : école, pairs, travail, médias) est le moteur principal de la construction de la personnalité.
\item Au Cameroun, \textbf{famille élargie, rites d'initiation, religion, quartier, langue} sont des facteurs majeurs, souvent plus prégnants que dans les sociétés individualistes occidentales.
\item Les \textbf{réseaux sociaux} introduisent des modèles globaux qui entrent en tension/résonance avec les modèles locaux.
\item La personnalité n'est pas un \emph{produit fini} de ces déterminismes : elle est un \textbf{projet de liberté} qui s'en saisit, les trie et les dépasse.
\item \textbf{Méthode d'analyse :} Biologique $\rightarrow$ Social/Culturel $\rightarrow$ Choix personnels $\rightarrow$ Interaction dynamique.
\end{itemize}
\end{aretenir}

% --- ILLUSTRATION 1 : Diagramme circulaire (pgfplots) ---
\begin{popfigure}
\begin{tikzpicture}
% Data: Famille 35, Ecole 20, Pairs 20, Culture/Religion 15, Choix 10
% Colors: popblue, poporange, popgreen, poppurple, poppink
\fill[popblue] (0,0) -- (0.00:2.3) arc (0.00:126.00:2.3) -- cycle;
\node[white,font=\scriptsize\bfseries] at (63.00:1.55) {35\,\%};
\fill[popblue] (3.0,1.90) rectangle ++(0.28,0.28); \node[anchor=west,font=\scriptsize] at (3.35,2.04) {Famille (35\,\%)};
\fill[poporange] (0,0) -- (126.00:2.3) arc (126.00:198.00:2.3) -- cycle;
\node[white,font=\scriptsize\bfseries] at (162.00:1.55) {20\,\%};
\fill[poporange] (3.0,1.48) rectangle ++(0.28,0.28); \node[anchor=west,font=\scriptsize] at (3.35,1.62) {\'Ecole (20\,\%)};
\fill[popgreen] (0,0) -- (198.00:2.3) arc (198.00:270.00:2.3) -- cycle;
\node[white,font=\scriptsize\bfseries] at (234.00:1.55) {20\,\%};
\fill[popgreen] (3.0,1.06) rectangle ++(0.28,0.28); \node[anchor=west,font=\scriptsize] at (3.35,1.20) {Pairs \& R\'eseaux (20\,\%)};
\fill[poppurple] (0,0) -- (270.00:2.3) arc (270.00:324.00:2.3) -- cycle;
\node[white,font=\scriptsize\bfseries] at (297.00:1.55) {15\,\%};
\fill[poppurple] (3.0,0.64) rectangle ++(0.28,0.28); \node[anchor=west,font=\scriptsize] at (3.35,0.78) {Culture \& Religion (15\,\%)};
\fill[poppink] (0,0) -- (324.00:2.3) arc (324.00:360.00:2.3) -- cycle;
\node[white,font=\scriptsize\bfseries] at (342.00:1.55) {10\,\%};
\fill[poppink] (3.0,0.22) rectangle ++(0.28,0.28); \node[anchor=west,font=\scriptsize] at (3.35,0.36) {Choix personnels (10\,\%)};
\draw[white,thick] (0,0) circle (2.3);
\end{tikzpicture}
\legende{Poids relatif estimé des facteurs dans la construction de la personnalité adolescente (enquête indicative, 30 élèves). La famille domine, mais l'addition « Pairs + Réseaux » (20\%) + « École » (20\%) + « Culture/Religion » (15\%) + « Choix » (10\%) = 65\% montre que la personnalité est largement \emph{co-construite} hors du seul foyer.}
\end{popfigure}

% Note: The \pie command is not standard pgfplots/TikZ. It requires a specific library (pgf-pie).
% The prompt says "Diagramme circulaire (camembert pgfplots)". Standard pgfplots does not have a built-in \pie command.
% I must implement a pie chart using standard TikZ/pgfplots (polar axis) to ensure compilability without extra libraries.
% Let's replace the above with a standard TikZ pie chart using \foreach.

% Replacing the popfigure content with a compilable TikZ pie chart.
% I will rewrite the popfigure block below.

% Actually, I need to output the *final* code. I will write the correct TikZ code for the pie chart directly in the final output.
% The prompt allows TikZ. I'll use a \foreach loop with cumulative angles.

% Let's prepare the TikZ code for the pie chart.
% Values: 35, 20, 20, 15, 10. Total 100.
% Colors: popblue, poporange, popgreen, poppurple, poppink.
% Labels in legend.

% I will write the final LaTeX code\coursec{Les facteurs sociaux et culturels de la personnalité}

Après avoir examiné les \cle{facteurs biologiques} (tempérament, hérédité, maturation nerveuse) qui constituent le \emph{substrat naturel} de la personnalité, nous devons maintenant explorer comment ce substrat est modelé, orienté et enrichi par l'environnement humain. La personnalité ne se déploie pas dans le vide : elle se construit dans et par la \cle{socialisation}, ce processus continu par lequel l'individu intériorise les normes, les valeurs, les langages et les pratiques de son groupe d'appartenance. Au Cameroun comme ailleurs, c'est au carrefour de la famille, de l'école, des pairs, de la culture, du milieu professionnel et des choix personnels que s'élabore l'identité unique de chaque élève de Première technique.

\begin{definition}[Socialisation]
La \textbf{socialisation} est le processus par lequel l'individu, dès la naissance et tout au long de sa vie, intériorise les normes, les valeurs, les croyances, les langages et les modes de comportement de la société ou du groupe auquel il appartient. Elle transforme un être biologique en un \cle{sujet social} capable d'interagir, de communiquer et de se reconnaître comme membre d'une culture. On distingue la \emph{socialisation primaire} (famille, petite enfance) et la \emph{socialisation secondaire} (école, travail, pairs, médias).
\end{definition}

\courssub{La famille : premier milieu d'éducation et de transmission}

La famille est le \textbf{premier agent de socialisation}. C'est le creuset où se forge le \cle{soi} naissant. Bien avant l'entrée à l'école, l'enfant a déjà intégré une vision du monde, une langue maternelle (ou un répertoire langagier), des interdits (tabous alimentaires, respect des aînés, gestion du regard), des rituels quotidiens et une place dans la lignée.

\begin{exemplebox}[Éducation traditionnelle dans trois aires culturelles camerounaises]
\begin{itemize}
\item \textbf{Zone Bamiléké (Ouest) :} L'éducation est souvent collective, confiée à la \emph{chefferie} et aux \emph{sociétés secrètes} (ex. \emph{Kwifo}, \emph{Ngumba}). L'enfant apprend le sens de l'effort, l'entrepreneuriat, le respect de la hiérarchie et le culte des ancêtres. Le \emph{nom} donné à la naissance (ex. \emph{Tchinda} « celui qui revient », \emph{Ngouafack} « la force du rocher ») programme déjà une attente identitaire.
\item \textbf{Zone Béti (Centre, Sud, Est) :} L'éducation passe par l'\emph{esprit de clan} (\emph{ayong}), le respect scrupuleux des aînés (\emph{nkùm}), l'apprentissage des proverbes (\emph{mvet}) et des danses initiatiques. La \emph{case des palabres} est le lieu où l'enfant observe la résolution des conflits et l'exercice de l'autorité légitime.
\item \textbf{Zone Peule (Nord, Adamaoua) :} Le \emph{pulaaku} (code de conduite peul) structure la personnalité : \emph{semteende} (pudeur, réserve), \emph{munyal} (patience, endurance), \emph{hakkille} (prudence, sagesse), \emph{ngorgu} (courage). L'apprentissage du pastoralisme et la transmission orale (contes, poésie) forgent une identité forte, mobile et communautaire.
\end{itemize}
\end{exemplebox}

Dans tous les cas, la famille transmet un \textbf{capital culturel} (langue, savoir-faire, normes) et un \textbf{capital symbolique} (nom, réputation, honneur) qui constituent le socle de la personnalité. La structure familiale (monoparentale, élargie, polygame, nucléaire) module aussi l'apprentissage de l'autonomie, de la solidarité ou de la rivalité fraternelle.

\courssub{L'école : socialisation secondaire, règles, savoirs et citoyenneté}

L'école technique au Cameroun opère une \textbf{socialisation secondaire} systématique. Elle ne se contente pas d'instruire ; elle \emph{formate} au rôle d'élève, puis de futur technicien ou professionnel.

\begin{itemize}
\item \textbf{Apprentissage des règles impersonnelles :} Ponctualité, port de l'uniforme, respect du règlement intérieur, évaluation par notes anonymes. L'élève découvre l'universalité de la norme : la même règle s'applique à tous, indépendamment de l'origine sociale.
\item \textbf{Acquisition de savoirs codifiés :} Mathématiques, physique, technologie, mais aussi français, anglais, éducation à la citoyenneté. Ces savoirs offrent des \cle{outils cognitifs} pour comprendre le monde et s'y insérer.
\item \textbf{Construction de la citoyenneté :} Levée des couleurs, chant de l'hymne national, cours d'éducation civique, participation aux clubs (environnement, droits de l'homme, entrepreneuriat). L'école vise à faire émerger un \emph{citoyen camerounais} conscient de ses droits et devoirs, au-delà de son appartenance ethnique ou religieuse.
\item \textbf{Socialisation professionnelle :} Les stages, les travaux pratiques en atelier, les visites d'entreprises initient aux codes du monde du travail : hiérarchie, sécurité, travail en équipe, rapport au temps productif.
\end{itemize}

\begin{attention}[Piège : croire que l'école « efface » la famille]
L'école ne remplace pas la famille ; elle \emph{superpose} un nouveau registre de normes. Les conflits de loyauté sont fréquents : l'élève peut vivre une \cle{dissonance cognitive} entre la valorisation de l'oralité et du respect silencieux de l'aîné (famille) et l'exigence de prise de parole argumentée, critique et individuelle (école). La personnalité se construit \emph{dans} cette tension, pas en dehors.
\end{attention}

\courssub{Le groupe des pairs et les réseaux sociaux : l'influence des modèles contemporains}

À l'adolescence, le \textbf{groupe des pairs} devient un agent de socialisation majeur, souvent concurrent de la famille et de l'école. C'est le laboratoire de l'\cle{identité juvénile} : on y teste des styles vestimentaires, des langages (argot, \emph{camfranglais}), des valeurs (solidarité, virilité, modernité), des limites (transgression, prise de risque).

\begin{itemize}
\item \textbf{Influence des amis :} Le besoin d'appartenance pousse à la conformité (effet de mode, consommation d'alcool, absentéisme) mais aussi à l'émulation positive (révision en groupe, projets associatifs, sport).
\item \textbf{Modèles médiatiques et influenceurs :} Artistes urbains (Stanley Enow, Locko, Mimie, Ko-C), footballeurs (André Onana, Vincent Aboubakar), youtubeurs, tiktokeurs camerounais ou de la diaspora. Ils incarnent la \emph{réussite}, la \emph{modernité}, la \emph{liberté}. L'adolescent construit son \emph{idéal du moi} en partie sur ces figures.
\item \textbf{Réseaux sociaux (WhatsApp, Facebook, TikTok, Instagram) :} Ils démultiplient l'espace des pairs. La \cle{réputation numérique} (likes, vues, commentaires) devient un enjeu d'estime de soi. Le cyberharcèlement, l'exposition à des contenus violents ou pornographiques, la pression de l'image parfaite sont de nouveaux risques pour la construction identitaire.
\end{itemize}

\begin{exemplebox}[Enquête de classe — Exemple chiffré]
Dans une classe de 30 élèves de Première technique à Yaoundé, on a demandé : « \emph{Quelles sont les trois influences qui te semblent les plus fortes sur ta façon de penser et d'agir aujourd'hui ?} »
\begin{center}
\begin{tabularx}{\linewidth}{|l|X|c|}\hline
\textbf{Rang} & \textbf{Facteur cité} & \textbf{Nombre de citations (sur 30)} \\ \hline
1 & Famille (parents, oncles/tantes, grands-parents) & 27 \\ \hline
2 & Amis / Groupe de pairs (réel + réseaux sociaux) & 22 \\ \hline
3 & École (professeurs, règlement, cours) & 18 \\ \hline
4 & Médias / Influenceurs / Artistes & 12 \\ \hline
5 & Religion / Église / Mosquée & 9 \\ \hline
6 & Choix personnels / Réflexion propre & 7 \\ \hline
\end{tabularx}
\end{center}
\legende{Enquête informelle, non scientifique, mais illustratrice de la hiérarchie perçue par les élèves. La famille reste la référence majeure, mais le duo « Amis + Réseaux » talonne l'école.}
\end{exemplebox}

\courssub{La culture et la religion : rites d'initiation, valeurs communautaires, appartenance religieuse}

La \cle{culture} au sens large (coutumes, langues, arts, rites, visions du monde) et la \cle{religion} fournissent des \textbf{grands récits} qui donnent du sens à l'existence et structurent la personnalité en profondeur.

\begin{itemize}
\item \textbf{Rites d'initiation :} Circoncision collective (Bassa, Bamiléké, Peuls), \emph{Ngondo} (Sawa), \emph{So} (Béti), \emph{Mvett} (Fang), retraite monastique chrétienne, mémorisation du Coran. Ces rites marquent le passage de l'enfance à l'âge adulte, confèrent un nouveau statut, transmettent un savoir secret et lient l'individu à la lignée des ancêtres ou à la communauté des croyants. Ils sont des \emph{ancrages identitaires} puissants.
\item \textbf{Valeurs communautaires :} L'\emph{ubuntu} (« Je suis parce que nous sommes »), la solidarité de clan (\emph{tontine}, \emph{njangi}), le respect des aînés, l'hospitalité, la honte comme régulateur social (\emph{« Cela ne se fait pas chez nous »}). La personnalité se définit \emph{par rapport aux autres}, non en isolation.
\item \textbf{Appartenance religieuse :} Catholicisme, protestantisme, islam, églises de réveil, religions traditionnelles. La foi offre un cadre éthique (dix commandements, piliers de l'islam), une espérance (salut, paradis), une communauté de soutien (paroisse, oumma) et des rituels quotidiens (prière, jeûne) qui disciplinent le corps et l'esprit. Pour beaucoup d'élèves, Dieu est le \emph{premier référent} de leur identité : « Je suis enfant de Dieu » ou « Je suis serviteur d'Allah ».
\end{itemize}

\begin{savaistu}[Le nom, première inscription identitaire au Cameroun]
Dans de nombreuses cultures camerounaises, le \textbf{nom attribué à l'enfant} n'est pas une simple étiquette : il exprime une \emph{personnalité attendue}, un \emph{destin symbolique} ou les circonstances de la naissance.
\begin{itemize}
\item \emph{Bamiléké :} \textbf{Ngouanet} (« Dieu m'a donné »), \textbf{Tchouakeu} (« Celui qui sait attendre »), \textbf{Kamdem} (« Le chef est arrivé »).
\item \emph{Béti :} \textbf{Ondoa} (« Le chemin »), \textbf{Mvondo} (« La pluie »), \textbf{Essomba} (« La force »).
\item \emph{Peul :} \textbf{Amadou} (« Très loué »), \textbf{Hapsatou} (« Celle qui juge équitablement »), \textbf{Boubacar} (« Père de Bakar »).
\item \emph{Chrétiens :} \textbf{Dieudonné}, \textbf{Bénédicte}, \textbf{Grâce}, \textbf{Prospérité}.
\end{itemize}
Porter ce nom, c'est porter une \emph{exigence} et une \emph{promesse}. L'élève de Première technique porte souvent en lui cette injonction symbolique : « Deviens ce que ton nom dit que tu es. »
\end{savaistu}

\courssub{Le milieu professionnel et économique : le métier, le quartier, les conditions de vie}

L'environnement matériel et professionnel agit comme un \textbf{moule silencieux} sur la personnalité. Ce n'est pas seulement le métier futur (technicien, ouvrier, ingénieur) qui compte, mais le \emph{milieu de vie présent}.

\begin{itemize}
\item \textbf{Le quartier (village, cité universitaire, quartier populaire, résidence huppée) :} L'espace physique détermine les rencontres, les dangers, les ressources, les modèles visibles. Un élève grandissant à \emph{Nkolbisson} (Yaoundé) ou \emph{Bépanda} (Douala) développe des stratégies de survie, de débrouillardise (\emph{Article 15}), de vigilance, différentes de celui qui grandit à \emph{Bastos} ou \emph{Akwa} résidentiel.
\item \textbf{Les conditions de vie :} Accès à l'eau, à l'électricité, à internet, à la santé, aux transports. La précarité économique peut forge la résilience, l'endurance, le sens du sacrifice — ou générer du fatalisme, de la rancœur, de l'agressivité.
\item \textbf{L'apprentissage informel / stage / petit métier :} Beaucoup d'élèves techniques travaillent les vacances (menuiserie, soudure, coiffure, commerce, cybercafé). Ils y apprennent la \emph{réalité du terrain} : relation client, gestion de l'argent, respect du patron, fierté du travail bien fait. Cette expérience module leur rapport à l'effort, à l'autorité, à l'argent.
\end{itemize}

\begin{attention}[Erreur fréquente : opposer inné et acquis comme s'ils s'excluaient]
\textbf{La personnalité ne naît pas de l'addition « biologique + social », mais de leur \emph{interaction permanente}.}
\begin{itemize}
\item Un tempérament \emph{extraverti} (biologique) s'épanouira différemment dans une famille valorisant la parole publique (Bamiléké) ou dans une famille valorisant la réserve (Peul).
\item Une prédisposition génétique à l'anxiété peut être amplifiée par un milieu insécurisant ou atténuée par un entourage sécurisant et une foi solide.
\item \textbf{Conclusion :} Il n'y a pas de « part d'inné » et de « part d'acquis » séparables. Il y a un \emph{devenir} singulier où le biologique offre des potentialités que le social actualise, oriente, inhibe ou sublime.
\end{itemize}
\end{attention}

\courssub{La liberté personnelle : la personnalité comme projet et non comme simple produit}

Si la personnalité était \emph{seulement} le résultat des facteurs précédents, l'homme serait un \emph{produit déterminé} de son hérédité et de son milieu. Or, l'expérience vécue et la philosophie (Sartre, Mounier, Kant) nous enseignent que l'être humain est \textbf{liberté}.

\begin{propriete}[La personnalité comme projet (Mounier, Sartre)]
La personnalité n'est pas un \emph{état} figé, mais un \emph{mouvement} : elle est ce \textbf{projet} par lequel le sujet s'approprie son héritage biologique et social, le trie, l'interprète, le refuse ou le transcende.
\begin{itemize}
\item \textbf{Appropriation :} Je ne subis pas ma timidité ; je \emph{choisis} de la travailler en m'inscrivant au club de théâtre.
\item \textbf{Interprétation :} Mon échec au Probatoire n'est pas une fatalité ; je le \emph{lis} comme un signal pour changer de méthode.
\item \textbf{Transcendance :} Né dans un quartier difficile, je \emph{décide} de devenir ingénieur pour transformer mon milieu.
\end{itemize}
\emph{Exigible au Bac :} Savoir expliquer que la liberté ne nie pas les conditionnements, mais qu'elle est la capacité de \emph{prendre de la distance} par rapport à eux pour poser un acte signifiant.
\end{propriete}

\begin{methode}[Analyser un cas concret : identifier les facteurs en interaction]
Pour comprendre la personnalité d'un individu réel ou fictif, suivez ces 4 étapes :
\begin{enumerate}
\item \textbf{Repérer les facteurs biologiques :} Tempérament apparent (calme, nerveux, leader né), santé, aptitudes physiques/intellectuelles innées, hérédité connue.
\item \textbf{Repérer les facteurs sociaux et culturels :} Famille (structure, valeurs, style éducatif), école (parcours, réussites/échecs, orientation), pairs (groupe d'appartenance, influence), culture/religion (rites, croyances, langue), milieu pro/éco (quartier, travail, ressources).
\item \textbf{Repérer les choix personnels (liberté) :} Décisions marquantes (orientation, rupture, engagement, conversion, effort soutenu), projets d'avenir, discours de justification (« Pourquoi j'ai fait ça ? »).
\item \textbf{Conclure sur l'interaction :} Montrez comment le biologique a été \emph{saisi} par le social, et comment la liberté a \emph{répondu} à cette double détermination. Évitez le déterminisme (tout est joué d'avance) comme le volontarisme naïf (je peux tout, n'importe quand).
\end{enumerate}
\end{methode}

\courssub{Étude de cas guidée : portrait croisé de deux élèves de Première technique}

Appliquons la méthode à deux profils types, inspirés de la réalité camerounaise.

\begin{exoresolu}[Portrait croisé : Amadou (village) vs Stéphanie (ville)]
\textbf{Énoncé :} Amadou, 17 ans, Première F (Mécanique) à Ngaoundéré. Stéphanie, 17 ans, Première G (Génie Civil) à Douala. Analysez les facteurs qui façonnent leur personnalité.

\textbf{Solution détaillée :}

\begin{center}
\begin{tabularx}{\linewidth}{|l|X|X|}\hline
\textbf{Facteur} & \textbf{Amadou (Ngaoundéré, milieu peul rural/urbain)} & \textbf{Stéphanie (Douala, milieu sawa/urbain cosmopolite)} \\ \hline
\textbf{Biologique} & Tempérament \emph{calme, endurant} (élevage), bonne constitution physique. & Tempérament \emph{vif, nerveux}, myopie corrigée, douée en dessin technique. \\ \hline
\textbf{Famille} & Père éleveur, mère au foyer, 3 épouses, 12 enfants. Éducation \emph{pulaaku} : réserve, respect absolu des aînés, apprentissage par observation. Frères aînés = modèles. & Parents fonctionnaires (mère infirmière, père comptable), 2 enfants. Éducation \emph{dialoguée}, valorisation de l'école, autonomie encouragée, prière familiale catholique. \\ \hline
\textbf{École} & Collège public, effectifs pléthoriques (80/classe), peu de matériel. Professeurs exigeants mais distants. Orientation technique par défaut (notes moyennes). Fierté d'être « technicien ». & Lycée privé bilingue, effectifs réduits (35), ateliers équipés. Orientation choisie (Génie Civil). Pression pour l'excellence (mention au Probatoire visée). \\ \hline
\textbf{Pairs / Réseaux} & Groupe d'amis du quartier : foot, thé, discussions sur l'émigration (Europe/USA), WhatsApp pour suivre actualités. Modèles : footballeurs locaux, grands frères revenus du « bush ». & Groupe mixte (filles/garçons) : révisions, sorties centre commercial, TikTok/Insta. Influenceuses mode/lifestyle (camerounaises + ivoiriennes). Modèles : ingénieures femmes, architectes. \\ \hline
\textbf{Culture / Religion} & Islam (sunnite), 5 prières/jour, Coran le week-end, Ramadan strict. Rite de circoncision passé à 12 ans (entrée dans l'âge adulte). Solidarité clanique (\emph{njangi}). & Catholicisme (paroisse dynamique), aumônerie jeunes, scoutisme. Ngondo suivi à la TV. Culture \emph{camfranglais} fluide. Identité « Doualaise » fière. \\ \hline
\textbf{Milieu pro / éco} & Aide au troupeau vacances. Connaissance du bétail, réparation motos. Quartier populaire, eau au forage, électricité irrégulière. Débrouillardise (\emph{Article 15}) quotidienne. & Stages en BTP (entreprise père ami). Visite chantiers (Pont Wouri, tours). Quartier résidentiel, fibre optique, climatisation. Gestion budget de poche. \\ \hline
\textbf{Choix personnels (Liberté)} & \textbf{Projet :} « Devenir chef d'atelier maintenance engins agricoles, moderniser l'élevage paternel. » Choix : refuser l'exode rural, valoriser le savoir local + technique école. & \textbf{Projet :} « Intégrer ENSTP Yaoundé, puis Master Urbanisme durable. » Choix : s'engager dans association « Filles du BTP », refuser mariage précoce proposé par oncle. \\ \hline
\textbf{Synthèse interaction} & Personnalité \emph{ancrée, collective, patiente, pragmatique}. La technique est un \emph{outil au service de la communauté}. Liberté = fidélité créative. & Personnalité \emph{ouverte, individuelle, ambitieuse, critique}. La technique est un \emph{levier d'ascension et de transformation}. Liberté = rupture assumée. \\ \hline
\end{tabularx}
\end{center}
\legende{Analyse comparative montrant comment des facteurs similaires (famille, école, religion) produisent des personnalités distinctes par leur \emph{interaction} avec le biologique et les \emph{choix personnels}.}
\end{exoresolu}

\begin{aretenir}[Synthèse : Les facteurs sociaux et culturels]
\begin{itemize}
\item La \textbf{socialisation} (primaire : famille ; secondaire : école, pairs, travail, médias) est le moteur principal de la construction de la personnalité.
\item Au Cameroun, \textbf{famille élargie, rites d'initiation, religion, quartier, langue} sont des facteurs majeurs, souvent plus prégnants que dans les sociétés individualistes occidentales.
\item Les \textbf{réseaux sociaux} introduisent des modèles globaux qui entrent en tension/résonance avec les modèles locaux.
\item La personnalité n'est pas un \emph{produit fini} de ces déterminismes : elle est un \textbf{projet de liberté} qui s'en saisit, les trie et les dépasse.
\item \textbf{Méthode d'analyse :} Biologique $\rightarrow$ Social/Culturel $\rightarrow$ Choix personnels $\rightarrow$ Interaction dynamique.
\end{itemize}
\end{aretenir}

% --- ILLUSTRATION 1 : Diagramme circulaire (TikZ) ---
\begin{popfigure}
\begin{tikzpicture}
% Pie chart data: Famille 35, Ecole 20, Pairs 20, Culture/Religion 15, Choix 10
% Colors
\colorlet{cFam}{popblue}
\colorlet{cEco}{poporange}
\colorlet{cPair}{popgreen}
\colorlet{cCul}{poppurple}
\colorlet{cChoix}{poppink}
% Radius
\def\radius{3cm}
% Starting angle
\def\startangle{90}
% Calculate cumulative angles
% Famille 35% -> 126 deg
% Ecole 20% -> 72 deg
% Pairs 20% -> 72 deg
% Culture 15% -> 54 deg
% Choix 10% -> 36 deg
% Draw slices
\foreach \val/\col/\lab [count=\i] in {35/cFam/Famille, 20/cEco/École, 20/cPair/Pairs \& Réseaux, 15/cCul/Culture \& Religion, 10/cChoix/Choix personnels} {
    \pgfmathsetmacro{\angle}{\val * 3.6}
    \pgfmathsetmacro{\midangle}{\startangle - \angle/2}
    \draw[fill=\col, draw=white, thick] (0,0) -- (\startangle:\radius) arc (\startangle:\startangle-\angle:\radius) -- cycle;
    % Legend marker (placed manually below)
    \pgfmathsetmacro{\startangle}{\startangle - \angle}
}
% Legend
\begin{scope}[yshift=-4.5cm, xshift=-5cm]
\foreach \val/\col/\lab [count=\i] in {35/cFam/Famille (35\%), 20/cEco/École (20\%), 20/cPair/Pairs \& Réseaux (20\%), 15/cCul/Culture \& Religion (15\%), 10/cChoix/Choix personnels (10\%)} {
    \fill[\col] (0, -\i*0.6) rectangle (0.5, -\i*0.6+0.4);
    \node[anchor=west] at (0.7, -\i*0.6+0.2) {\lab};
}
\end{scope}
\end{tikzpicture}
\legende{Poids relatif estimé des facteurs dans la construction de la personnalité adolescente (enquête indicative, 30 élèves). La famille domine (35\%), mais l'addition « Pairs + Réseaux » (20\%) + « École » (20\%) + « Culture/Religion » (15\%) + « Choix » (10\%) = 65\% montre que la personnalité est largement \emph{co-construite} hors du seul foyer.}
\end{popfigure}

% --- ILLUSTRATION 2 : Carte mentale TikZ ---
\begin{popfigure}
\begin{tikzpicture}[
    mindmap,
    concept color=popblue,
    level 1 concept/.append style={level distance=4.5cm, sibling angle=60},
    level 2 concept/.append style={level distance=3cm, sibling angle=45},
    every node/.style={font=\small},
    text=white
]
\node[concept, scale=1.2] {Ma personnalité}
    child[concept color=poporange] { node[concept] {Famille}
        child { node[concept, scale=0.6] {Valeurs} }
        child { node[concept, scale=0.6] {Langue} }
        child { node[concept, scale=0.6] {Interdits} }
        child { node[concept, scale=0.6] {Nom/Identité} }
    }
    child[concept color=popgreen] { node[concept] {École}
        child { node[concept, scale=0.6] {Règles} }
        child { node[concept, scale=0.6] {Savoirs} }
        child { node[concept, scale=0.6] {Citoyenneté} }
        child { node[concept, scale=0.6] {Métier} }
    }
    child[concept color=poppurple] { node[concept] {Culture}
        child { node[concept, scale=0.6] {Rites} }
        child { node[concept, scale=0.6] {Communauté} }
        child { node[concept, scale=0.6] {Ubuntu} }
        child { node[concept, scale=0.6] {Arts/Proverbes} }
    }
    child[concept color=poppink] { node[concept] {Religion}
        child { node[concept, scale=0.6] {Éthique} }
        child { node[concept, scale=0.6] {Rituels} }
        child { node[concept, scale=0.6] {Communauté foi} }
        child { node[concept, scale=0.6] {Espérance} }
    }
    child[concept color=popgold] { node[concept] {Amis/Pairs}
        child { node[concept, scale=0.6] {Modèles} }
        child { node[concept, scale=0.6] {Réseaux soc.} }
        child { node[concept, scale=0.6] {Argot/Style} }
        child { node[concept, scale=0.6] {Appartenance} }
    }
    child[concept color=popdark] { node[concept] {Choix pers.}
        child { node[concept, scale=0.6] {Projet} }
        child { node[concept, scale=0.6] {Responsabilité} }
        child { node[concept, scale=0.6] {Créativité} }
        child { node[concept, scale=0.6] {Transcendance} }
    };
\end{tikzpicture}
\legende{Carte mentale des facteurs de la personnalité. Au centre, le sujet unificateur ; six branches principales (Famille, École, Culture, Religion, Amis/Pairs, Choix personnels) avec des sous-branches illustrant les contenus transmis ou vécus. La personnalité émerge au carrefour de ces six flux.}
\end{popfigure}

\coursec{La valeur de la personne selon Kant : la dignité}

Nous avons vu dans la section précédente que la personnalité se construit sous l'influence de facteurs sociaux et culturels : la famille, l'école, la langue, les traditions du village ou du quartier façonnent ce que chacun devient. Mais une question demeure en suspens : si la personnalité varie d'une culture à l'autre et d'un individu à l'autre, y a-t-il quelque chose qui vaut \emph{également} en tout être humain, quelles que soient ses origines, sa richesse ou sa place dans la société ? Le domestique et son patron, le mendiant de carrefour et le ministre, ont-ils la même valeur ? C'est à cette question que répond Emmanuel Kant avec une notion capitale : la \cle{dignité}.

\courssub{Kant et le contexte des Lumières}

\begin{definition}[Emmanuel Kant]
Emmanuel \cle{Kant} (1724-1804) est un philosophe allemand, figure majeure du mouvement des \cle{Lumières}, ce courant du XVIII\textsuperscript{e} siècle qui affirme la primauté de la raison contre les préjugés, les superstitions et l'autorité aveugle. Son ouvrage moral fondamental est le \emph{Fondement de la métaphysique des mœurs} (1785), où il établit les principes d'une morale fondée sur la raison.
\end{definition}

Le XVIII\textsuperscript{e} siècle est le siècle où l'on commence à affirmer que tous les hommes naissent libres et égaux en droits. Pourtant, à la même époque, l'esclavage et la traite négrière font encore commerce des êtres humains : on achète et on vend des personnes comme on achète du tissu ou du sucre. Kant va donner à la révolte contre cette situation son fondement philosophique le plus solide : montrer, par la raison, \emph{pourquoi} aucun être humain ne peut jamais être traité comme une marchandise.

\begin{savaistu}[Le philosophe sédentaire]
Kant a mené une vie d'une régularité proverbiale : on raconte que les habitants de Königsberg (sa ville natale, aujourd'hui Kaliningrad) réglaient leurs montres sur son passage lors de sa promenade quotidienne. Il n'a presque jamais quitté sa ville. Pourtant, sa défense de la dignité humaine a fait le tour du monde : elle inspire aujourd'hui les constitutions et les tribunaux du monde entier, y compris en Afrique et au Cameroun. Preuve que les idées voyagent plus loin que les hommes.
\end{savaistu}

\courssub{La distinction fondamentale : le prix et la dignité}

Partons d'une intuition simple. Au marché Mokolo à Yaoundé, un sac de riz a un \cle{prix} : si le vendeur est trop cher, je peux acheter un autre sac, identique, chez le voisin. Le sac de riz a un \emph{équivalent} : n'importe quoi de même valeur peut le remplacer. Prenons maintenant un être humain : votre mère, votre ami, vous-même. Peut-on le remplacer par un « équivalent » ? Peut-on dire : « celui-ci est trop cher, donnez-m'en un autre » ? L'idée même révolte. Voilà l'intuition que Kant va transformer en concept rigoureux.

\begin{definition}[Prix et dignité]
Le \cle{prix} est la valeur d'une chose qui peut être échangée contre un équivalent (argent, autre chose). Ce qui a un prix est \emph{remplaçable}. La \cle{dignité} est la valeur absolue et inaliénable de la personne, qui la place \emph{au-dessus de tout prix} : ce qui a une dignité n'a pas d'équivalent et ne peut être remplacé par rien.
\end{definition}

Kant formule cette distinction dans le \emph{Fondement de la métaphysique des mœurs} : dans le règne des fins, tout a soit un prix, soit une dignité. Ce qui a un prix peut être mis à la place de quelque chose d'autre, comme équivalent ; ce qui est au-dessus de tout prix, et n'admet par conséquent aucun équivalent, possède une dignité.

\begin{popfigure}
\begin{center}
\begin{tikzpicture}[scale=1]
\draw[fill=popblueL, rounded corners=4pt] (0,0) rectangle (5.4,4.2);
\draw[fill=poporangeL, rounded corners=4pt] (6.0,0) rectangle (11.4,4.2);
\node[font=\bfseries, text=popblue] at (2.7,3.75) {Ce qui a un prix};
\node[font=\bfseries, text=poporange!70!black] at (8.7,3.75) {Ce qui a une dignit\'e};
\node[align=left, text width=4.6cm, font=\small] at (2.7,1.9) {\textbullet\ Une marchandise (riz, t\'el\'ephone)\\[2pt] \textbullet\ Un service (transport, coupe de cheveux)\\[2pt] \textbullet\ A un \'equivalent : peut \^etre remplac\'e\\[2pt] \textbullet\ Se compare, se n\'egocie};
\node[align=left, text width=4.6cm, font=\small] at (8.7,1.9) {\textbullet\ La personne humaine\\[2pt] \textbullet\ N'a aucun \'equivalent : irrempla\c{c}able\\[2pt] \textbullet\ Ne se vend ni ne s'ach\`ete\\[2pt] \textbullet\ Exige le respect inconditionnel};
\end{tikzpicture}
\legende{La distinction kantienne entre le prix (valeur relative, échangeable) et la dignité (valeur absolue, sans équivalent).}
\end{center}
\end{popfigure}

\begin{attention}[La dignité ne se mérite pas]
Erreur fréquente : croire que la dignité se \emph{mérite} par les bonnes actions ou se \emph{perd} par les mauvaises. Pour Kant, c'est faux : la dignité est \emph{inhérente} à l'humanité. Elle appartient à tout être humain du seul fait qu'il est un être raisonnable et libre — même au paresseux, même au coupable, même au criminel. On peut punir le criminel (il doit répondre de ses actes), mais on ne peut ni le torturer, ni l'humilier, ni le traiter comme une bête : sa dignité demeure intacte. Ne confondez donc jamais dignité et honneur, mérite ou réputation, qui eux se gagnent et se perdent.
\end{attention}

\courssub{La personne comme fin en soi}

De la distinction précédente découle le principe moral le plus célèbre de Kant, l'une des formules de l'\emph{impératif catégorique} (le commandement moral inconditionnel) :

\begin{aretenir}[La formule de l'humanité]
« Agis de telle sorte que tu traites l'humanité, en ta personne et en celle d'autrui, toujours en même temps comme une \cle{fin}, et jamais simplement comme un \cle{moyen}. »
\end{aretenir}

Décomposons cette formule. Un \emph{moyen} est ce qu'on utilise pour atteindre un but extérieur à lui : le marteau est un moyen pour planter le clou ; une fois utilisé, on le range, on s'en débarrasse quand il est usé. Une \emph{fin en soi} est ce qui vaut par lui-même, ce qui est sa propre fin : on ne l'utilise pas, on le respecte. Dire que la personne est une fin en soi, c'est dire qu'on ne peut jamais la réduire au rôle d'instrument dans les projets d'autrui.

\begin{propriete}[Tout être raisonnable est une fin en soi]
Pour Kant, tout être raisonnable est une fin en soi ; nul ne peut être traité \emph{uniquement} comme un moyen. (Propriété admise, sans démonstration exigible.)
\end{propriete}

\begin{attention}[« Jamais simplement comme un moyen » : le mot décisif]
Kant n'interdit pas d'utiliser les services d'autrui ! Il écrit « jamais \emph{simplement} comme un moyen ». Quand je paie le moto-taxi qui me conduit au lycée, il est bien un moyen de transport pour moi — mais pas \emph{simplement} un moyen : je le paie librement, il a consenti à ce service, je le salue et le respecte. Son humanité est en même temps traitée comme une fin. Ce qui est interdit, c'est de réduire l'autre au statut d'outil : le tromper, le contraindre, l'exploiter, c'est-à-dire l'utiliser sans tenir compte de sa volonté et de sa valeur propres.
\end{attention}

\courssub{Le fondement de la dignité : la raison et la liberté}

Pourquoi l'être humain a-t-il une dignité, et non un prix ? La réponse de Kant tient en deux mots : la \cle{raison} et la \cle{liberté} morale. L'animal agit par instinct ; la chose n'agit pas du tout. Seul l'être humain peut se demander : « que dois-je faire ? », distinguer le bien du mal, et choisir librement d'obéir à la loi morale plutôt qu'à ses désirs. Cette capacité — la raison pratique, c'est-à-dire la faculté de se donner à soi-même une loi morale (l'\emph{autonomie}) — est ce qui élève l'homme au-dessus de toute la nature et de tout prix.

Conséquence capitale : ce fondement est le \emph{même} en tous. La raison morale ne dépend ni de la richesse, ni de l'ethnie, ni du sexe, ni du statut social, ni même de l'intelligence scolaire. Le planteur de Mbalmayo et le professeur d'université, le Bamiléké et le Béti, le riche et le pauvre possèdent également cette capacité de se conduire moralement. La dignité est donc rigoureusement \emph{égale} pour tous les êtres humains.

\begin{experience}[Une observation pour établir le raisonnement]
\textbf{Document observé.} Dans une cour de justice camerounaise, un voleur reconnu coupable est condamné à une peine de prison. Pendant l'audience, le juge s'adresse à lui avec respect, il est défendu par un avocat, il peut faire appel du jugement.\par
\textbf{Questions d'analyse.}
\begin{enumerate}
\item Le coupable est-il traité comme une chose que l'on jette ? Non : on lui applique une \emph{loi}, c'est-à-dire qu'on le considère comme un être raisonnable, capable de comprendre la règle qu'il a violée et responsable de ses actes.
\item Que montre le fait qu'il soit défendu et entendu ? Que même en le punissant, la société reconnaît en lui une personne, sujet de droits, et non un objet.
\item Conclusion : la punition elle-même, lorsqu'elle est juste, est un hommage rendu à la dignité du coupable : on le traite comme un être libre qui devait choisir le bien, non comme un animal nuisible qu'on élimine. Ce que Kant résume ainsi : le respect est dû \emph{même au criminel}.
\end{enumerate}
\end{experience}

\courssub{Les conséquences pratiques du principe de dignité}

Le principe kantien n'est pas une belle phrase abstraite : il interdit concrètement tout ce qui réduit un être humain à l'état de moyen.

\begin{itemize}
\item \textbf{L'esclavage} est la négation absolue de la dignité : l'esclave est une chose possédée, achetée et vendue, purement moyen du profit de son maître.
\item \textbf{L'exploitation} : payer un travailleur un salaire de misère qui ne lui permet pas de vivre, c'est utiliser sa force de travail sans respecter sa personne ; le profit passe avant l'humain.
\item \textbf{L'humiliation} : insulter, frapper, rabaisser publiquement quelqu'un (un élève, un employé, un époux ou une épouse), c'est nier en actes la valeur absolue de sa personne.
\item \textbf{La tromperie et la manipulation} : mentir à quelqu'un pour obtenir quelque chose de lui, c'est se servir de sa raison comme d'un instrument, sans qu'il puisse consentir en connaissance de cause.
\end{itemize}

\begin{exemplebox}[Application au Cameroun : quand l'humain devient un simple moyen]
\textbf{Le travail des enfants dans les cacaoyères.} Dans certaines exploitations de la zone forestière du Centre et du Sud, des enfants, parfois âgés de moins de douze ans, travaillent toute la journée à la place d'aller à l'école, pour un salaire dérisoire ou pour aider sans compensation. L'enfant est alors traité \emph{simplement comme un moyen} de production du cacao : sa santé, son éducation, son avenir — autant de dimensions de sa personne qui est une fin en soi — sont sacrifiés au profit. La loi camerounaise fixe d'ailleurs l'âge minimal d'admission au travail et prohibe les pires formes de travail des enfants.\par
\textbf{La mendicité forcée.} Dans les grandes villes (Douala, Yaoundé), certains enfants sont envoyés mendier aux carrefours par des adultes qui collectent le soir le produit de la mendicité. L'enfant n'est plus une personne à éduquer, mais un instrument de rapport d'argent. Les deux cas illustrent exactement ce que la formule de Kant condamne : traiter l'humanité, en la personne d'autrui, simplement comme un moyen.
\end{exemplebox}

\begin{exemplebox}[Un héritage universel et national]
L'article 1\textsuperscript{er} de la Déclaration universelle des droits de l'homme de 1948 proclame : « Tous les êtres humains naissent libres et égaux en dignité et en droits. » Le préambule de la Constitution camerounaise du 18 janvier 1996 reprend cette affirmation en déclarant l'attachement du peuple camerounais aux droits fondamentaux inscrits dans cette Déclaration. Ainsi, l'idée kantienne de dignité, pensée à Königsberg en 1785, est devenue un principe juridique qui protège chaque Camerounais.
\end{exemplebox}

\begin{popfigure}
\begin{center}
\begin{tikzpicture}[scale=1]
\draw[->, thick, popdark] (0,0) -- (12.4,0);
\foreach \x/\y in {1.2/1724, 3.2/1785, 5.6/1804, 8.2/1948, 10.8/1996}{\draw[thick] (\x,-0.12) -- (\x,0.12); \node[below, font=\scriptsize] at (\x,-0.15) {\y};}
\node[above, align=center, font=\scriptsize, text width=2.1cm] at (1.2,0.15) {Naissance de Kant \`a K\"onigsberg};
\node[above, align=center, font=\scriptsize, text width=2.3cm] at (3.2,0.15) {\emph{Fondement de la m\'etaphysique des m\oe urs}};
\node[above, align=center, font=\scriptsize, text width=1.8cm] at (5.6,0.15) {Mort de Kant};
\node[above, align=center, font=\scriptsize, text width=2.3cm, text=popblue] at (8.2,0.15) {D\'eclaration universelle des droits de l'homme};
\node[above, align=center, font=\scriptsize, text width=2.1cm, text=poporange!70!black] at (10.8,0.15) {Constitution camerounaise};
\draw[popgreen, <->] (0.6,1.35) -- (6.2,1.35) node[midway, above, font=\scriptsize\itshape, text=popgreen!60!black] {Si\`ecle des Lumi\`eres};
\end{tikzpicture}
\legende{Frise chronologique : la pensée de Kant (1724-1804), héritière des Lumières, prépare les grands textes modernes sur la dignité humaine (1948, 1996).}
\end{center}
\end{popfigure}

\courssub{Discussion guidée : le salaire achète-t-il la personne ?}

\begin{exoresolu}[Le salaire d'un employé de maison à Yaoundé rémunère-t-il son travail ou achète-t-il sa personne ?]
\textbf{Énoncé.} À Yaoundé, une famille emploie une aide de maison pour le ménage, la cuisine et la garde des enfants, moyennant un salaire mensuel. Un élève affirme : « Puisqu'on la paie, on peut exiger d'elle tout ce qu'on veut. » Ce raisonnement est-il conforme à la philosophie de Kant ?\par
\tcblower
\textbf{Solution détaillée.}
\begin{enumerate}
\item \textbf{Dégager l'enjeu.} La question est de savoir ce que le salaire achète : un \emph{travail} (une prestation, qui a un prix) ou une \emph{personne} (qui a une dignité, donc pas de prix).
\item \textbf{Appliquer la distinction prix/dignité.} Selon Kant, le salaire rémunère légitimement le \emph{service} rendu : le travail est une chose qui a un prix et peut s'échanger. Mais la personne qui accomplit ce travail garde sa dignité : elle n'est pas vendue avec ses bras.
\item \textbf{Appliquer la formule de la fin en soi.} L'employé de maison est un moyen (elle rend un service), mais elle ne doit jamais être traitée \emph{simplement} comme un moyen. Sont donc légitimes : un contrat clair, un salaire décent, des horaires raisonnables, le respect et la politesse. Sont illégitimes : les insultes, les coups, le travail sans repos, le refus de payer, les avances imposées — car ces conduites réduisent la personne à un instrument.
\item \textbf{Conclure.} Le raisonnement de l'élève est donc \emph{faux} : payer quelqu'un ne donne aucun droit sur sa personne. Le salaire achète un travail, jamais un être humain. La dignité de l'employée demeure égale à celle de ses employeurs.
\end{enumerate}
\end{exoresolu}

\begin{exercice}[Pour s'entraîner]
Un chef d'entreprise de Douala licencie une ouvrière parce qu'elle est tombée enceinte, en déclarant : « Elle ne m'est plus utile. » Analysez cette situation à l'aide de la formule de Kant : en quoi l'ouvrière est-elle traitée « simplement comme un moyen » ? Quelles exigences concrètes le respect de sa dignité imposait-il à l'employeur ? Rédigez votre réponse en un paragraphe argumenté (thèse, application du principe, conclusion).
\end{exercice}

\begin{corrige}[Exercice]
\textbf{Éléments de correction.} \textbf{Thèse} : en licenciant l'ouvrière au seul motif qu'elle « ne lui est plus utile », le chef d'entreprise la traite simplement comme un moyen de production, ce que Kant interdit. \textbf{Application} : la grossesse ne supprime ni la raison, ni la liberté, ni la dignité de la travailleuse ; sa valeur de personne est indépendante de son « utilité » momentanée pour l'entreprise. Le mot « utile » révèle précisément la faute morale : l'être humain est mesuré à l'aune du profit, comme une machine qu'on jette quand elle ralentit. \textbf{Exigences concrètes} : le respect de la dignité commandait de maintenir l'emploi, d'aménager les conditions de travail, de respecter le congé de maternité prévu par le droit du travail. \textbf{Conclusion} : une entreprise peut rechercher le profit, mais jamais au prix de la réduction d'une personne à un simple instrument ; la dignité prime sur l'utilité.
\end{corrige}

\begin{aretenir}[L'essentiel de la section]
\begin{itemize}
\item Pour Kant (1724-1804), philosophe des Lumières, tout a \emph{un prix} ou \emph{une dignité} : ce qui a un prix a un équivalent ; ce qui a une dignité est au-dessus de tout prix.
\item La personne humaine est une \emph{fin en soi} : il faut traiter l'humanité, en soi et en autrui, toujours en même temps comme une fin, jamais simplement comme un moyen.
\item Le fondement de la dignité est la raison et la liberté morale, présentes en tout être humain sans distinction de richesse, d'ethnie ou de statut : la dignité est égale pour tous et ne se perd jamais, même chez le criminel.
\item Conséquences : condamnation de l'esclavage, de l'exploitation (travail des enfants, mendicité forcée) et de l'humiliation ; héritage dans la Déclaration universelle de 1948 et la Constitution camerounaise de 1996.
\end{itemize}
\end{aretenir}

\coursec{La valeur de la personne selon Mounier : le personnalisme}

Après avoir analysé la conception kantienne qui fonde la valeur de la personne sur l'autonomie rationnelle et la dignité intrinsèque, nous nous tournons maintenant vers une philosophie qui, tout en affirmant la valeur absolue de la personne, déplace le regard : la personne n'est pas un atome moral isolé, mais un être de relation et d'engagement. C'est le pari du \cle{personnalisme} d'\cle{Emmanuel Mounier} (1905-1950), figure majeure de la philosophie française du XX\textsuperscript{e} siècle, fondateur de la revue \emph{Esprit} (1932). Face aux crises de son temps (crise économique de 1929, montée des totalitarismes, déshumanisation du travail ouvrier), Mounier refuse l'alternative stérile entre un individualisme bourgeois qui isole et un collectivisme totalitaire qui écrase. Il propose une troisième voie : la personne comme \emph{être en devenir}, qui ne se réalise qu'en se donnant.

\courssub{Contexte et naissance du personnalisme}

Emmanuel Mounier élabore sa pensée dans l'entre-deux-guerres, une période où l'Europe cherche ses repères. L'individualisme libéral a montré son incapacité à créer du lien social solide (crise économique, chômage de masse), tandis que le fascisme, le nazisme et le stalinisme proposent une « solution » collective au prix de la suppression de la personne. Mounier, chrétien non conformiste, réunit autour de la revue \emph{Esprit} des intellectuels de toutes obédiences pour penser une \cle{révolution personnaliste et communautaire}. Son ouvrage majeur, \emph{Le Personnalisme} (1949), expose cette doctrine de manière synthétique.

\begin{definition}[Le personnalisme]
Doctrine philosophique qui fait de la \cle{personne} — être singulier, relationnel, libre et responsable — la valeur centrale et le point de départ de toute réflexion éthique, sociale et politique. Contre l'individualisme (l'individu comme atome fermé) et le collectivisme (la masse qui dissout le singulier), le personnalisme affirme que la personne ne s'accomplit que dans la \cle{communion} (relation d'amour, d'amitié, de service) et l'\cle{engagement} dans l'histoire.
\end{definition}

L'intuition de départ est simple : \emph{je ne suis vraiment quelqu'un que devant quelqu'un}. Un Robinson Crusoé seul sur son île reste un individu biologique ; il devient pleinement personne quand Vendredi arrive, car il doit répondre, parler, choisir pour et avec un autre.

\courssub{La personne comme être relationnel : la disponibilité}

Pour Mounier, la structure fondamentale de la personne est la \cle{relation}. L'individu a des besoins ; la personne a une \cle{vocation}. Cette vocation ne s'invente pas en chambre close, elle se découvre dans la rencontre.

\begin{propriete}[Réalisation de la personne chez Mounier]
Chez Mounier, la personne ne se réalise pleinement que dans la \cle{communion} avec les autres et l'\cle{engagement} dans le monde. La \cle{disponibilité} — attitude d'ouverture active à l'appel de l'autre et de la situation — est le mode propre de l'existence personnelle. Il n'y a pas de personne « en soi », il n'y a que des personnes « pour autrui » et « pour le monde ».
\end{propriete}

La \cle{disponibilité} n'est pas passivité. C'est une tension vigilante : être « à la disposition de » sans se perdre. Mounier distingue trois niveaux de relation :
\begin{enumerate}
    \item La \textbf{relation fonctionnelle} : je te considère comme un moyen (le caissier, l'outil). Niveau de l'individu.
    \item La \textbf{relation interpersonnelle} : je te reconnais comme une fin, un « toi » irréductible (amitié, amour, dialogue). Niveau de la personne.
    \item La \textbf{communion} : union des libertés dans un projet commun sans absorption réciproque (communauté de travail, famille, engagement politique juste). C'est l'idéal personnaliste.
\end{enumerate}

\begin{popfigure}
\begin{tikzpicture}[
    node distance=2.6cm and 2.2cm,
    main/.style={circle, draw=popblue, thick, fill=popblue!10, minimum width=2.4cm, font=\bfseries\small, align=center},
    dim/.style={rectangle, rounded corners, draw=poporange, thick, fill=poporange!10, minimum width=3cm, minimum height=1.4cm, font=\small, align=center},
    arrow/.style={->, thick, popdark, shorten >=2pt, shorten <=2pt}
]
\node[main, align=center] (personne){PERSONNE\\ (être relationnel)};
\node[dim, align=center] (incarnation) [below left=of personne]{INCARNATION\\ corps, temps, histoire,\\ enracinement charnel};
\node[dim, align=center] (relation) [below right=of personne]{RELATION À AUTRUI\\ disponibilité, dialogue,\\ communion, amour};
\node[dim, align=center] (vocation) [below=4.2cm of personne]{VOCATION / ENGAGEMENT\\ appel à se dépasser, action,\\ responsabilité, service du bien commun};
\draw[arrow] (personne.south west) -- (incarnation.north) node[midway, left, font=\footnotesize] {s'enracine dans};
\draw[arrow] (personne.south east) -- (relation.north) node[midway, right, font=\footnotesize] {s'ouvre à};
\draw[arrow] (incarnation.south) |- (vocation.west);
\draw[arrow] (relation.south) |- (vocation.east);
\end{tikzpicture}
\legende{Les trois dimensions constitutives de la personne chez Emmanuel Mounier. La personne est le nœud vivant où s'articulent l'incarnation charnelle, l'ouverture relationnelle et l'élan vocationnel.}
\end{popfigure}

\courssub{Les trois dimensions de la personne : incarnation, relation, vocation}

Mounier refuse le spiritualisme désincarné. La personne n'est pas une âme pure ; elle est \textbf{incarnation}. Le corps n'est pas une prison ni un outil, c'est \emph{mon} corps : mon visage exprime mon intériorité, mes mains agissent sur le monde, ma fatigue et ma maladie me rappellent ma finitude. L'incarnation, c'est l'ancrage dans le temps, l'histoire, une culture, un sol.

La \textbf{relation à autrui} est le lieu où la personne naît à elle-même. Le visage de l'autre m'appelle et me fait sortir de mon repli sur moi-même. La \cle{disponibilité} est cette capacité à « faire place » en soi à l'autre sans l'assimiler. C'est le refus de l'emprise.

La \textbf{vocation} (ou engagement) couronne l'édifice. La personne n'est pas un statut acquis une fois pour toutes, c'est une \cle{tâche}. Se faire personne, c'est répondre à l'appel de la situation historique : lutter contre l'injustice, bâtir une œuvre, élever des enfants, soigner, enseigner. L'engagement n'est pas optionnel ; il est la \emph{forme} même de l'existence personnelle.

\begin{attention}[Personnalisme vs Individualisme : le piège majeur]
\emph{Erreur fréquente :} confondre « personne » et « individu » et croire que le personnalisme est un individualisme raffiné.

\emph{Réalité :} Mounier est un critique virulent de l'individualisme bourgeois. L'\cle{individu} est défini par la \cle{séparation} (atome, nombre, droits subjectifs, repli sur son intérêt). La \cle{personne} est définie par la \cle{relation} (communion, responsabilité, don de soi). L'individualisme dit : « Je suis moi, les autres sont des obstacles ou des outils. » Le personnalisme dit : « Je ne suis moi que grâce aux autres et pour les autres. »
\end{attention}

\courssub{Critique de l'individualisme bourgeois et du collectivisme totalitaire}

La force du personnalisme de Mounier réside dans sa \cle{critique sociale} à deux fronts.

\begin{itemize}
    \item \textbf{Contre l'individualisme bourgeois (capitalisme libéral) :} il réduit l'homme à un \cle{consommateur} et un \cle{producteur} isolé, mesurant toute valeur à l'aune de l'utilité et du profit. La société devient une « collection d'atomes » régie par le contrat et la concurrence. L'homme est aliéné par l'argent, la technique, la bureaucratie. Mounier dénonce l'\cle{esprit bourgeois} : confort, sécurité, refus du risque et de l'engagement.
    \item \textbf{Contre le collectivisme totalitaire (fascisme, nazisme, stalinisme) :} il sacrifie la personne à la \cle{masse}, à la \cle{race}, à la \cle{classe} ou à l'\cle{État}. L'individu devient un rouage ; la personne est niée au nom d'une pseudo-unité collective. Mounier refuse la « personne collective » (l'État-personne) qui absorbe les personnes réelles.
\end{itemize}

La \cle{révolution personnaliste} vise à instaurer une société de \cle{communautés de personnes} (familles, équipes de travail, quartiers, syndicats, mouvements) où le pouvoir est partagé, où l'économie est au service de l'homme (coopératives, économie sociale), où la politique est service du bien commun.

\courssub{Application concrète au Cameroun : la tontine (njangi) comme personnalisme vécu}

La philosophie de Mounier n'est pas une abstraction occidentale ; elle trouve une résonance profonde dans les cultures africaines, notamment camerounaises, où l'adage « Je suis parce que nous sommes » structure le lien social. La \cle{tontine} (ou \cle{njangi} dans la région du Nord-Ouest, \cle{djangui} chez les Sawa du Littoral) en est l'illustration parlante.

\begin{exemplebox}[La tontine de Bafoussam : 12 membres, une communauté de personnes]
À Bafoussam, un groupe de 12 femmes commerçantes du marché forme une tontine mensuelle. Chaque membre verse 50 000 FCFA par mois. Le « pot » de 600 000 FCFA revient à tour de rôle à l'une d'elles.
\begin{itemize}
    \item \textbf{Incarnation :} ces femmes sont des corps au travail (port de charges, vente au marché, fatigue, maternité). Leur engagement n'est pas désincarné.
    \item \textbf{Relation :} la confiance mutuelle prime sur le contrat écrit. On connaît la vie de l'autre (maladie, deuil, rentrée scolaire). Si l'une faillit, le groupe soutient, ne juge pas seulement. C'est de la \cle{disponibilité} effective.
    \item \textbf{Vocation :} l'argent du pot sert à agrandir la boutique, payer une opération, financer une scolarité. Chaque femme s'accomplit \emph{en} faisant réussir le groupe. La réussite individuelle est inséparable de la solidité du collectif.
\end{itemize}
C'est du personnalisme en acte : la personne s'édifie dans la communion solidaire.
\end{exemplebox}

\begin{exoresolu}[Analyse personnaliste d'une situation de tontine]
\textbf{Énoncé :} Dans une tontine de 10 membres à Yaoundé, l'un des membres, M. K., perd son emploi et ne peut plus cotiser pendant trois mois. Le groupe décide de cotiser à sa place pendant cette période, puis il remboursera plus tard sans intérêt. Analysez cette décision à la lumière des concepts de Mounier (disponibilité, communion, critique de l'individualisme).
\tcblower
\textbf{Solution détaillée :}
\begin{enumerate}
    \item \textbf{Identification de la structure :} la tontine n'est pas une simple association d'individus calculateurs (logique contractualiste individualiste). C'est une \cle{communauté de personnes}.
    \item \textbf{Disponibilité :} les 9 autres membres font \emph{place} à la détresse de M. K. dans leur budget propre. Ils sortent de leur intérêt immédiat (garder leur argent) pour répondre à l'appel de l'autre. C'est l'attitude personnaliste par excellence : « Je suis responsable de l'autre. »
    \item \textbf{Communion vs Contrat :} un contrat bancaire aurait exclu M. K. (défaut de paiement). La communion personnaliste \emph{intègre} la faille. Le groupe se renforce en sauvant l'un des siens. La personne de M. K. est soutenue \emph{par} la communauté.
    \item \textbf{Critique de l'individualisme :} cette décision est inintelligible pour l'agent purement égoïste et calculateur (perte sèche, risque de non-remboursement). Elle n'est intelligible que si la valeur suprême n'est pas le profit individuel mais la \cle{vie personnelle commune}.
    \item \textbf{Conclusion :} le groupe agit en sujet responsable : il décide, il assume, il engage ses ressources pour l'un des siens. C'est un acte de solidarité qui fait grandir à la fois la personne secourue et les personnes secourantes.
\end{enumerate}
\end{exoresolu}

\courssub{Comparaison Kant / Mounier : convergence et divergence}

Si Kant et Mounier s'accordent sur la \cle{valeur absolue de la personne} (elle ne saurait être réduite à un moyen), ils divergent sur le \cle{fondement} et la \cle{modalité} de cette valeur.

\begin{center}
\rowcolors{1}{popcream}{white}
\begin{tabularx}{\linewidth}{|l|X|X|}
\hline
\rowcolor{popblueL} \textbf{Critère} & \textbf{KANT (morale de la raison)} & \textbf{MOUNIER (personnalisme)} \\
\hline
\textbf{Fondement de la valeur} & La \cle{raison pratique} : l'autonomie, capacité de se donner à soi-même la loi morale. La dignité vient de la loi morale en moi. & L'\cle{existence incarnée} : l'acte de se faire personne (création continue). La valeur vient de la relation et de l'engagement. \\
\hline
\textbf{Conception de la personne} & \cle{Sujet moral universel} : même structure rationnelle en tous les hommes, abstraction faite du corps et de l'histoire. & \cle{Être singulier, incarné, historique} : un « je » unique, un corps, un temps, une culture. La personne est une tâche à accomplir. \\
\hline
\textbf{Relation à autrui} & \cle{Respect} de la dignité : ne jamais traiter autrui seulement comme un moyen. Universalisable, exigeant, à distance respectueuse. & \cle{Disponibilité} et \cle{communion} : s'engager positivement pour autrui. Proximité, amour, amitié, service. \\
\hline
\textbf{Politique} & \cle{Droit} et \cle{État de droit} : cadre juridique protégeant la liberté de chacun. & \cle{Communautés de personnes} : famille, équipe, quartier ; économie au service de l'homme. \\
\hline
\textbf{Exemple : la tontine} & Respect des règles convenues ; si l'un faillit, on applique la règle prévue (justice stricte). & Solidarité effective : le groupe porte le faible. La personne du membre prime sur la règle comptable. \\
\hline
\end{tabularx}
\end{center}

\begin{aretenir}[Synthèse Kant / Mounier]
\begin{itemize}
    \item \textbf{Convergence :} tous deux placent la \cle{personne} au sommet de l'échelle des valeurs. Tous deux refusent l'instrumentalisation (traiter autrui comme simple moyen). Tous deux fondent une éthique de l'exigence.
    \item \textbf{Divergence majeure :}
    \begin{itemize}
        \item \textbf{Kant :} personne = \cle{sujet moral} universel, autonome. Valeur = \cle{dignité} (sans prix, incomparable). Relation = \cle{respect}. Politique = \cle{État de droit}.
        \item \textbf{Mounier :} personne = \cle{être historique} singulier, incarné, relationnel. Valeur = \cle{vocation} (accomplissement dans le don). Relation = \cle{disponibilité / communion}. Politique = \cle{communautés de personnes}.
    \end{itemize}
    \item \textbf{Complémentarité :} Kant donne le \cle{garde-fou} juridico-moral (le « minimum » intangible : ne pas violer la dignité). Mounier donne l'\cle{idéal} existentiel (le « maximum » vivant : devenir personne par l'amour et l'action). Une société juste a besoin des deux : le droit kantien pour protéger l'espace de liberté, l'esprit personnaliste pour l'habiter fraternellement.
\end{itemize}
\end{aretenir}

\begin{savaistu}[Le personnalisme, source d'inspiration pour l'Afrique]
Le personnalisme de Mounier a profondément marqué la pensée africaine francophone des indépendances. \cle{Léopold Sédar Senghor} (Sénégal) s'en est inspiré pour penser un socialisme africain centré sur la personne humaine comme fin de la société. Au Cameroun, des penseurs comme \cle{Fabien Eboussi Boulaga} ou \cle{Marcien Towa} ont dialogué avec cette tradition pour penser une modernité africaine non aliénée. Aujourd'hui encore, les \cle{mouvements d'économie sociale et solidaire} (coopératives agricoles, GIC — Groupements d'Initiative Commune, mutuelles de santé villageoises) opèrent souvent selon une logique personnaliste : la rentabilité économique est subordonnée à l'épanouissement des personnes membres. Mounier n'est pas un auteur de musée ; il est un outil pour penser le \cle{développement communautaire} au Cameroun.
\end{savaistu}

\courssub{Méthode : appliquer les notions à une situation concrète (savoir-faire)}

\begin{methode}[Analyser une situation au prisme du personnalisme de Mounier]
\textbf{Consigne type :} « En vous appuyant sur la philosophie de Mounier, analysez la situation suivante : [grève, conflit familial, choix d'orientation, engagement associatif, gestion d'une crise au travail...] »

\textbf{Étapes de la démarche :}
\begin{enumerate}
    \item \textbf{Identifier les acteurs :} sont-ils traités comme des \cle{individus} (numéros, fonctions, coûts) ou comme des \cle{personnes} (singularité, histoire, visage) ?
    \item \textbf{Repérer la structure de la relation :} est-elle \cle{fonctionnelle} (instrumentale), \cle{contractuelle} (droit, devoir formel) ou \cle{personnelle/communautaire} (disponibilité, confiance, don) ?
    \item \textbf{Évaluer l'engagement :} y a-t-il \cle{vocation} (sens, dépassement de soi, service) ou simple \cle{exécution} (contrainte, routine, intérêt) ?
    \item \textbf{Qualifier le collectif :} est-ce une \cle{masse} (anonyme, manipulable), une \cle{société contractuelle} (somme d'individus) ou une \cle{communauté de personnes} (communion, responsabilité partagée) ?
    \item \textbf{Conclure en problématisant :} la situation révèle-t-elle une \cle{régression individualiste} (repli, égoïsme) ou \cle{totalitaire} (écrasement du singulier) ? Ou au contraire une \cle{dynamique personnaliste} (croissance des personnes par la relation) ? Proposer une issue personnaliste.
\end{enumerate}
\end{methode}

\begin{exercice}[Entraînement : le chef d'équipe et l'erreur]
M. D. est chef d'équipe dans une usine de transformation de bois à Douala. Un jeune apprenti, K., commet une erreur qui abîme une pièce coûteuse. Le règlement intérieur prévoit un avertissement formel et une retenue sur salaire. M. D. choisit de ne pas appliquer la sanction. Il appelle K. dans son bureau, lui montre l'erreur calmement, lui demande comment il va, s'il a des soucis à la maison, puis ils réparent la pièce ensemble. M. D. dit : « On apprend en se trompant. L'important, c'est que tu grandisses. »
\begin{enumerate}
    \item Qualifiez l'attitude de M. D. au regard de la distinction de Mounier : \cle{individu} vs \cle{personne}.
    \item Identifiez la \cle{disponibilité} et la \cle{communion} dans cette scène.
    \item Que dirait Kant de cette dispense de sanction ? (Rappel : exigence de justice et loi universelle.)
    \item Rédigez un paragraphe argumenté (une quinzaine de lignes) : « La gestion personnaliste des ressources humaines est-elle compatible avec l'exigence de justice ? »
\end{enumerate}
\end{exercice}

\begin{corrige}[Exercice : le chef d'équipe et l'erreur]
\begin{enumerate}
    \item \textbf{Individu vs Personne :} le règlement traite K. comme un \cle{individu-fonction} (unité de production, source d'erreur, coût). M. D. traite K. comme une \cle{personne} : il s'adresse à son histoire (« soucis à la maison »), à son devenir (« tu grandis »), à sa dignité présente.
    \item \textbf{Disponibilité / Communion :} \cle{disponibilité} : M. D. fait place à la réalité de K. (temps d'écoute, suspension du règlement). \cle{Communion} : « on répare ensemble » — union des libertés dans un acte commun de réparation et de formation. Le chef descend de son piédestal hiérarchique pour rejoindre l'apprenti.
    \item \textbf{Regard kantien :} Kant distinguerait deux plans. Sur le plan du \emph{droit}, le règlement est une règle posée pour tous ; ne pas l'appliquer peut créer une inégalité de traitement ou un manque de respect pour l'universalité de la loi. Sur le plan de la \emph{morale}, M. D. accomplit un devoir envers autrui : aider K. à progresser. Kant dirait : « Tu as bien fait moralement, mais attention au plan juridique : la règle doit être améliorée pour tous, pas contournée pour un seul. »
    \item \textbf{Argumentation type :} la gestion personnaliste ne nie pas la justice ; elle la dépasse par le haut. La justice (Kant) pose le plancher : droits, règles, égalité formelle. Le personnalisme (Mounier) vise le plafond : épanouissement, réconciliation, croissance des personnes. M. D. n'a pas violé la justice s'il signale l'incident pour faire évoluer la règle, ou s'il assume personnellement le coût de sa décision. L'incompatibilité n'est réelle que si la « compassion » devient arbitraire (favoritisme). Le personnalisme vrai exige justice \emph{et} miséricorde : la règle protège la personne, la personne vivifie la règle.
\end{enumerate}
\end{corrige}

\coursec{Synthèse et méthode : appliquer et argumenter}

Après avoir analysé la \cle{dignité} chez Kant et l'\cle{accomplissement personnel} chez Mounier, nous disposons désormais des outils conceptuels pour saisir l'unité de la leçon : la personnalité est une \cle{construction} dynamique (biologique, sociale, culturelle, libre), tandis que la personne est une \cle{valeur absolue} qui s'impose à notre action morale. Cette section finale a pour but de consolider ces acquis, de vous entraîner à les mobiliser dans des situations concrètes — y compris camerounaises — et de vous armer méthodologiquement pour l'épreuve du Probatoire / Baccalauréat technique.

\courssub{Bilan synthétique de la leçon}

\begin{aretenir}[Bilan : personnalité et personne]
La \cle{personnalité} est l'ensemble des traits psychologiques, affectifs, intellectuels et moraux qui singularisent un individu. Elle \textbf{se construit} tout au long de l'existence par l'interaction de quatre facteurs :
\begin{itemize}
    \item \textbf{Biologiques} : hérédité, tempérament, constitution neurophysiologique (donné de départ).
    \item \textbf{Sociaux} : famille, école, groupe d'appartenance, classe sociale (socialisation primaire et secondaire).
    \item \textbf{Culturels} : langue, valeurs, rites, traditions, médias, réseaux sociaux (héritage symbolique).
    \item \textbf{Libres} : choix éthiques, engagements, projets de vie (part de responsabilité personnelle).
\end{itemize}
La \cle{personne}, quant à elle, n'est pas un « objet » parmi d'autres : elle possède une \textbf{valeur absolue}, c'est-à-dire une \cle{dignité} (Kant) qui interdit de la traiter comme un simple moyen, et appelle un \cle{respect} inconditionnel. Pour Mounier, cette valeur s'actualise dans la \cle{communion} et le \cle{don de soi} : la personne ne se réduit pas à son individualité biologique ou sociale, elle s'\textbf{accomplit} dans l'ouverture à autrui et à des valeurs transcendantes.
\end{aretenir}

\begin{popfigure}
\centering
\begin{tikzpicture}[
    node distance=1.2cm and 2.5cm,
    box/.style={draw=popblue, thick, rounded corners, fill=popblue!5, text width=3.2cm, align=center, minimum height=1.8cm, font=\small},
    arrow/.style={-Stealth, thick, popdark},
    labelbox/.style={font=\footnotesize\bfseries, text=popblue, anchor=west}
]
% Nodes
\node[box, align=center] (bio){Facteurs biologiques\\ (Tempérament, hérédité)};
\node[box, right=of bio, align=center] (socio){Facteurs sociaux\\ (Famille, école, pairs)};
\node[box, right=of socio, align=center] (cult){Facteurs culturels\\ (Valeurs, médias, rites)};
\node[box, right=of cult, align=center] (libre){Liberté \& Choix\\ (Engagements, éthique)};
\node[box, below=1.5cm of socio, fill=poporange!5, draw=poporange, text width=5cm] (perso) {\textbf{PERSONNALITÉ}\\\emph{Construction dynamique,}\\\emph{ unique, en devenir}};
\node[box, right=1.5cm of perso, fill=popgreen!5, draw=popgreen, text width=5cm] (persn) {\textbf{PERSONNE}\\\emph{Valeur absolue : Dignité (Kant)}\\\emph{\& Accomplissement (Mounier)}};
% Arrows
\draw[arrow] (bio.south) -- ++(0,-0.5) -| (perso.north west);
\draw[arrow] (socio.south) -- ++(0,-0.5) -| (perso.north);
\draw[arrow] (cult.south) -- ++(0,-0.5) -| (perso.north east);
\draw[arrow] (libre.south) -- ++(0,-0.5) -| (perso.east);
\draw[arrow, poppurple, very thick] (perso.east) -- node[above, font=\small\bfseries, text=poppurple] {Fonde / Exige} (persn.west);
% Labels
\node[labelbox] at (bio.north west) {1.};
\node[labelbox] at (socio.north west) {2.};
\node[labelbox] at (cult.north west) {3.};
\node[labelbox] at (libre.north west) {4.};
\end{tikzpicture}
\legende{Schéma récapitulatif : de la construction de la personnalité (facteurs 1 à 4) à la valeur absolue de la personne (dignité kantienne, accomplissement mouniériste).}
\end{popfigure}

\courssub{Application à des situations concrètes camerounaises}

Le programme exige de « appliquer les notions de l'unité à des situations concrètes de la vie courante ». Voici quatre cas types, analysés avec les outils de la leçon.

\begin{exemplebox}[Cas 1 : Le respect des personnes âgées au village]
\textbf{Situation :} Dans un village du Nord-Ouest, un chef de famille âgé, atteint de démence, est isolé dans une case à l'écart ; ses enfants gèrent ses biens sans le consulter.
\textbf{Analyse :}
\begin{itemize}
    \item \textbf{Personnalité altérée :} La maladie neurodégénérative (facteur biologique) détruit la mémoire, le langage, l'autonomie : la personnalité \emph{empirique} s'effondre.
    \item \textbf{Personne intouchable :} Pour \textbf{Kant}, la dignité ne dépend \emph{pas} des facultés actuelles. L'humanité en sa personne (sa rationalité passée, sa qualité de fin en soi) subsiste. Le traiter comme un « fardeau » ou un « objet » gérable est un crime contre l'humanité en sa personne (\emph{traiter comme moyen}).
    \item \textbf{Personnalisme (Mounier) :} La communion exige de \emph{porter} l'autre dans sa vulnérabilité. L'exclusion brise la communauté personnelle. Le respect dû à l'ancêtre (valeur culturelle) doit se fonder sur la dignité ontologique, pas sur l'utilité sociale.
\end{itemize}
\textbf{Conclusion :} La société a le devoir de \emph{soigner} et d'\emph{inclure}, non de gérer l'indigne.
\end{exemplebox}

\begin{exemplebox}[Cas 2 : Le harcèlement scolaire (lycée de Yaoundé)]
\textbf{Situation :} Une élève de Première est moquée quotidiennement sur son physique et son origine ethnique sur un groupe WhatsApp de la classe ; elle décroche scolairement.
\textbf{Analyse :}
\begin{itemize}
    \item \textbf{Destruction de la personnalité :} Le harcèlement (facteur social toxique) anéantit l'estime de soi, l'identité narrative, la confiance — piliers de la construction personelle.
    \item \textbf{Négation de la personne :} Les harceleurs traitent la victime comme un \textbf{moyen} de divertissement ou de cohésion de groupe par l'exclusion (Kant : violation de l'impératif catégorique). Ils nient sa subjectivité, sa douleur, sa finitude.
    \item \textbf{Réseaux sociaux :} L'amplification numérique (facteur culturel) déshumanise : l'écran fait écran à la \emph{présence} personnelle (Mounier). L'anonymat relatif lève l'interdit moral immédiat.
\end{itemize}
\textbf{Action :} Sanction éducative (responsabiliser les auteurs), soutien psychologique (reconstruire la personnalité), éducation à l'empathie (reconnaître l'autre comme personne).
\end{exemplebox}

\begin{exemplebox}[Cas 3 : La traite des personnes (route Douala–Ngaoundéré)]
\textbf{Situation :} De jeunes filles sont recrutées par de fausses promesses d'emploi, transportées de force, exploitées sexuellement ou en travail forcé.
\textbf{Analyse :}
\begin{itemize}
    \item \textbf{Instrumentation totale :} La traite est le \textbf{paradigme} du traitement de la personne comme \emph{simple moyen} (marchandise, outil de profit). Prix mis sur une vie = négation radicale de la \cle{dignité} (Kant : « tout a un prix ou une dignité »).
    \item \textbf{Aliénation personnaliste :} Mounier parle d'\emph{aliénation} : la personne est réduite à l'état de chose, dépossédée de sa liberté, de son corps, de son projet. La \cle{communion} est remplacée par la \cle{domination}.
    \item \textbf{Facteurs structurels :} Pauvreté (facteur social), patriarcat, corruption (facteurs culturels/institutionnels) créent la vulnérabilité, mais n'excusent pas le crime.
\end{itemize}
\textbf{Réponse :} Justice pénale (sanctionner le crime), réinsertion (restaurer la subjectivité), prévention structurelle (éducation, emploi, égalité filles-garçons).
\end{exemplebox}

\begin{exemplebox}[Cas 4 : Usage des réseaux sociaux (TikTok, Facebook, X) par les lycéens]
\textbf{Situation :} Un élève passe 6h/jour à « scroller », construit une « persona » idéale, cherche la validation par les « likes », compare sa vie intérieure aux vitrines altrui.
\textbf{Analyse :}
\begin{itemize}
    \item \textbf{Personnalité fragmentée :} L'algorithme (facteur culturel technique) capte l'attention, morcelle le temps, impose des modèles normatifs (beauté, succès, consommation). Risque : \emph{hétérodétermination} — la personnalité se construit \emph{par} l'extérieur, non \emph{par} soi.
    \item \textbf{Personne en spectacle :} Mounier dénonce la \emph{société de masse} où l'individu « fait persona » (masque social) au lieu d'\emph{être personne}. La quête de \emph{visibilité} remplace la quête de \emph{vérité} et de \emph{communion}.
    \item \textbf{Autonomie menacée :} Kant lie dignité et \cle{autonomie} (se donner sa propre loi). L'addiction algorithmique est une \emph{hétéronomie} : la volonté est pilotée par des ingénieurs d'attention.
\end{itemize}
\textbf{Piste éducative :} Éducation aux médias (esprit critique), temps de déconnexion choisis, culture de l'intériorité (journal intime, lecture, silence).
\end{exemplebox}

\begin{methode}[Appliquer une notion à une situation concrète (3 étapes)]
\begin{enumerate}
    \item \textbf{Nommer \& Définir :} Citez la notion précise (ex. : « dignité », « aliénation », « hétéronomie », « socialisation ») et donnez sa définition rigoureuse en une phrase.
    \item \textbf{Qualifier le fait :} Montrez \emph{précisément} en quoi la situation illustre la notion (ex. : « Le fait que l'algorithme décide du contenu vu caractérise l'hétéronomie : la volonté n'est pas cause de son objet »).
    \item \textbf{Évaluer / Argumenter :} Dites ce que la notion permet de \emph{juger} ou de \emph{comprendre} (ex. : « Donc l'élève perd en autonomie, fondement de sa dignité morale ; il y a risque d'aliénation personnaliste »).
\end{enumerate}
\end{methode}

\begin{attention}[Piège fréquent : « Name-dropping » sans explication]
\emph{Écrire : « Selon Kant, la personne a une dignité » ou « Mounier dit qu'il faut la communion »} ne rapporte \textbf{aucun point} au Bac. Toute référence d'auteur doit être :
\begin{itemize}
    \item \textbf{Développée :} Expliquez \emph{en quoi} consiste la thèse (ex. : « Kant fonde la dignité sur l'autonomie : être fin en soi, c'est se donner sa loi morale »).
    \item \textbf{Relée au sujet :} Montrez \emph{pourquoi} cet auteur éclaire \emph{ce} problème précis (ex. : « Ainsi, le harcèlement nie l'autonomie de la victime en la soumettant à la loi d'autrui »).
\end{itemize}
\textbf{Règle d'or :} 1 citation = 3 lignes d'explication minimum.
\end{attention}

\courssub{Méthodologie de la dissertation philosophique}

L'épreuve du Probatoire / Baccalauréat technique est une \textbf{dissertation} (sujet unique, 4h). Voici la méthode canonique, adaptée à nos notions.

\begin{methode}[Structure type de la dissertation (4 grandes étapes)]
\begin{enumerate}
    \item \textbf{Introduction (≈ 30 min)} :
        \begin{itemize}
            \item \textbf{Accroche :} Citation, fait de société, paradoxe, définition du sens commun. \emph{Ex. : « "Tout a un prix ou une dignité" (Kant). Mais dans nos sociétés marchandes, la personne elle-même semble avoir un cours... »}
            \item \textbf{Analyse des termes :} Définir « personne », « prix », « valeur », « dignité ».
            \item \textbf{Problématisation :} Transformer le sujet en \textbf{question philosophique} tensionnelle. \emph{Ex. : « La personne humaine a-t-elle un prix ? » $\rightarrow$ « La valeur de la personne est-elle réductible à une valeur marchande (prix) ou possède-t-elle une valeur intrinsèque (dignité) qui la met hors commerce ? »}
            \item \textbf{Annonce du plan :} « Nous montrerons d'abord que... (thèse), puis que... (antithèse), pour enfin établir que... (synthèse). »
        \end{itemize}
    \item \textbf{Développement (≈ 2h30) — Structure dialectique en 3 parties (ou 2 parties fortes + synthèse intégrée)} :
        \begin{itemize}
            \item \textbf{Partie I (Thèse) :} La personne \emph{semble} avoir un prix / être traitée comme ayant un prix.
                \begin{itemize}
                    \item Argument 1 : Réalité socio-économique (salaires, assurances, traite, GPA, marché des organes) $\rightarrow$ \textbf{fait}.
                    \item Argument 2 : Réduction utilitariste (Bentham) : l'individu compte pour un, le bonheur total prime $\rightarrow$ sacrifice possible.
                    \item Argument 3 : Aliénation / réification (Marx, Mounier) : la personne \emph{devient} chose (marchandise, force de travail).
                \end{itemize}
            \item \textbf{Partie II (Antithèse) :} La personne \emph{ne peut} avoir de prix ; elle a une \emph{dignité}.
                \begin{itemize}
                    \item Argument 1 : \textbf{Kant} — Dignité = valeur incomparable. Autonomie = fin en soi. Prix = valeur conditionnelle, remplaçable. Dignité = inconditionnelle, irremplaçable.
                    \item Argument 2 : \textbf{Mounier} — Personne = sujet spirituel, communion. Ne se possède pas, ne s'échange pas. S'accomplit dans le don gratuit.
                    \item Argument 3 : Droits de l'homme, bioéthique (interdiction clonage, commerce organes) : consensus juridique mondial sur le « hors commerce ».
                \end{itemize}
            \item \textbf{Partie III (Synthèse / Dépassement) :} Distinguer \textbf{prix de marché} et \textbf{valeur morale} ; \textbf{condition de vie} et \textbf{dignité de vie}.
                \begin{itemize}
                    \item On peut \emph{payer} des \emph{moyens} pour protéger la personne (santé, éducation, justice) sans mettre un prix \emph{sur} la personne.
                    \item La société doit \emph{garantir} les conditions de l'autonomie (revenu décent, soins) pour que la dignité ne reste pas abstraite.
                    \item Responsabilité politique : lutter contre toute marchandisation rampante (données personnelles, corps, attention).
                \end{itemize}
        \end{itemize}
    \item \textbf{Conclusion (≈ 30 min)} : Voir boîte méthode ci-dessous.
    \item \textbf{Relecture (≈ 15 min)} : Orthographe, cohérence, noms d'auteurs exacts, paragraphes aérés.
\end{enumerate}
\end{methode}

\begin{methode}[Rédiger et justifier une conclusion personnelle (3 temps)]
\begin{enumerate}
    \item \textbf{Rappeler la problématique :} Reprenez la question centrale sous forme affirmative. \emph{Ex. : « Au terme de cette analyse, il apparaît que la personne humaine n'a pas de prix, elle a une dignité. »}
    \item \textbf{Trancher en mobilisant un auteur de la leçon :} Affirmez votre position en l'adossant à la thèse la plus solide. \emph{Ex. : « Comme le démontre Kant, la dignité est cette valeur "au-dessus de tout prix" qui fonde le respect absolu ; Mounier ajoute qu'elle s'actualise dans la communion, non dans l'échange. »}
    \item \textbf{Ouvrir sur une question nouvelle / perspective concrète :} Ne fermez pas la pensée. \emph{Ex. : « Reste à savoir si nos sociétés numériques, qui monétisent l'attention et les données intimes, ne font pas courir le risque d'une "marchandisation douce" de la personne, sous couvert de gratuité apparente. »}
\end{enumerate}
\end{methode}

\begin{exoresolu}[Entraînement guidé : Introduction et Conclusion sur « La personne humaine a-t-elle un prix ? »]
\textbf{Sujet :} \emph{La personne humaine a-t-elle un prix ?}

\textbf{Introduction rédigée :}
« "Tout a un prix ou une dignité. Ce qui a un prix peut être remplacé par quelque chose d'équivalent ; ce qui est au-dessus de tout prix, n'a point d'équivalent, a une dignité." (Kant, \emph{Fondements de la métaphysique des mœurs}). Cette distinction tranche net : le prix suppose l'échange, l'équivalence, le marché ; la dignité suppose l'incomparable, l'irremplaçable, le respect absolu. Pourtant, la réalité contemporaine — salaires, assurances-vie, mères porteuses, trafic d'organes, données personnelles vendues aux annonceurs — semble mettre un prix sur l'humain. \textbf{Problématique :} La valeur de la personne est-elle réductible à une valeur marchande (prix), ou possède-t-elle une valeur intrinsèque (dignité) qui la met radicalement hors commerce ? \textbf{Plan :} Nous analyserons d'abord les mécanismes par lesquels la société assigne un prix à l'humain (I), puis nous montrerons que la personne, comme fin en soi (Kant) et sujet de communion (Mounier), possède une dignité inaliénable (II), pour enfin distinguer les conditions matérielles de l'existence (qui coûtent) de la valeur de l'existant (qui ne s'achète pas) (III). »

\textbf{Conclusion rédigée :}
« Au terme de cette dissertation, \textbf{la personne humaine n'a pas de prix : elle a une dignité}. Le prix est la valeur d'un \emph{objet} échangeable ; la dignité est la valeur d'un \emph{sujet} moral. Kant a prouvé que l'autonomie — capacité de se donner sa loi — fonde ce statut de "fin en soi" ; Mounier a montré que la personne s'accomplit dans le don gratuit et la communion, non dans le contrat marchand. \textbf{Toutefois}, reconnaître la dignité ne suffit pas : encore faut-il que la société garantisse les \emph{conditions concrètes} de son exercice (santé, éducation, revenu, justice numérique). \textbf{Ouverture :} À l'ère de l'économie de l'attention et de l'IA générative, la vigilance philosophique doit porter sur les nouvelles formes d'aliénation "gratuite" où l'utilisateur devient le produit : la personne a-t-elle encore prise sur sa propre subjectivité ? »
\tcblower
\textbf{Analyse de la copie (grille Bac) :}
\begin{itemize}
    \item \textbf{Problématique} : Claire, tensionnelle, termes définis (prix vs dignité). ✓
    \item \textbf{Plan} : Dialectique (Réalité marchande / Fondement moral / Articulation). ✓
    \item \textbf{Auteurs} : Kant (autonomie, fin en soi) et Mounier (don, communion) \emph{expliqués et mobilisés}. ✓
    \item \textbf{Exemples} : Concrets, variés (salaire, GPA, données, IA). ✓
    \item \textbf{Conclusion} : Rappel probléma, tranché (Kant/Mounier), ouverture (IA/attention). ✓
    \item \textbf{Expression} : Correcte, paragraphée, vocabulaire précis (inaliénable, réification, hétéronomie). ✓
\end{itemize}
\textbf{Note estimée :} 17–19/20.
\end{exoresolu}

\courssub{Méthodologie de l'explication de texte}

L'explication de texte (souvent en devoir surveillé, parfois au Bac) suit une logique rigoureuse.

\begin{methode}[Explication de texte : 5 étapes obligatoires]
\begin{enumerate}
    \item \textbf{Situer l'auteur et l'œuvre :} Nom, époque, courant (Lumières, personnalisme...), œuvre, contexte. \emph{Ex. : « Extrait de \emph{Fondements de la métaphysique des mœurs} (1785), Kant, philosophie critique, morale déontologique. »}
    \item \textbf{Dégager la thèse principale :} En une phrase : « Le texte soutient que... » (la thèse = réponse de l'auteur à sa question).
    \item \textbf{Structurer le texte :} Découper en 2–3 \emph{moments logiques} (pas des paragraphes !) : \emph{Problème / Argument 1 / Argument 2 / Conclusion de l'auteur}. Donnez un titre à chaque moment.
    \item \textbf{Expliquer les arguments (cœur de l'exercice) :} Pour chaque moment :
        \begin{itemize}
            \item Reformulez l'argument \emph{avec vos mots}.
            \item Définissez les concepts techniques du texte (\emph{impératif catégorique, fin en soi, autonomie, hétéronomie, prix, dignité, valeur, personne, individu, communion...}).
            \item Montrez \emph{comment} l'argument prouve la thèse (enchainement logique).
            \item \textbf{Illustrez par un exemple personnel} (camerounais si possible) pour montrer que vous avez compris.
        \end{itemize}
    \item \textbf{Discussion critique \& Conclusion :}
        \begin{itemize}
            \item \textbf{Portée :} Pourquoi ce texte est-il important ? (Fondement des droits de l'homme, bioéthique...).
            \item \textbf{Limites / Objections :} Une objection possible (ex. : « Mais comment respecter la dignité d'un criminel ? » / « La dignité abstraite ne nourrit pas son homme »).
            \item \textbf{Réponse de l'auteur / Votre avis argumenté.}
            \item \textbf{Conclusion brève :} Résumé de la thèse + actualisation.
        \end{itemize}
\end{enumerate}
\end{methode}

\begin{attention}[Erreurs fatales en explication de texte]
\begin{itemize}
    \item \textbf{Le paraphrase :} Recopier le texte en changeant deux mots. $\rightarrow$ Note éliminatoire. \textbf{Il faut expliquer, pas répéter.}
    \item \textbf{Le hors-texte :} Faire sa dissertation personnelle sans coller au texte. $\rightarrow$ Hors sujet.
    \item \textbf{L'oubli des concepts :} Ne pas définir \emph{impératif catégorique}, \emph{fin en soi}, \emph{personnalisme}, \emph{aliénation} quand ils sont dans le texte.
    \item \textbf{L'absence d'exemple :} Un argument philosophique sans illustration concrète reste abstrait = note moyenne.
\end{itemize}
\end{attention}

\courssub{Critères d'évaluation au Baccalauréat technique}

\begin{exemplebox}[Grille indicative de notation (sur 20 points) — Dissertation]
\begin{center}
\begin{tabularx}{\linewidth}{|l|X|c|}\hline
\rowcolor{popblueL}
\textbf{Critère} & \textbf{Descriptif d'excellence} & \textbf{Poids} \\ \hline
\textbf{Problématique} & Question philosophique claire, tensionnelle, termes définis, enjeu explicite. & 3 pts \\ \hline
\textbf{Structure / Plan} & Dialectique (Thèse/Antithèse/Synthèse) ou analytique rigoureuse, transitions logiques, paragraphes équilibrés. & 3 pts \\ \hline
\textbf{Argumentation} & Arguments philosophiques (pas d'opinions), concepts maniés avec justesse, enchainement logique. & 5 pts \\ \hline
\textbf{Mobilisation auteurs} & Kant \& Mounier (au minimum) \emph{expliqués}, cités à propos, reliés au sujet. Pas de name-dropping. & 3 pts \\ \hline
\textbf{Exemples concrets} & Au moins 2–3 exemples précis (dont 1 camerounais), analysés comme arguments, pas décoratifs. & 2 pts \\ \hline
\textbf{Conclusion} & Rappel probléma, réponse tranchée + auteur, ouverture pertinente. & 2 pts \\ \hline
\textbf{Expression / Orthographe} & Français correct, vocabulaire philosophique précis, paragraphes aérés, lisibilité. & 2 pts \\ \hline
\end{tabularx}
\end{center}
\legende{La pertinence de la problématique et la solidité de l'argumentation pèsent plus que la longueur. Une copie de 4 pages dense et rigoureuse vaut mieux que 8 pages bavardes.}
\end{exemplebox}

\begin{savaistu}[Les classiques indémodables du Bac camerounais]
Depuis 20 ans, les sujets sur la personne reviennent sous formes variées :
\begin{itemize}
    \item \textbf{« La personne a-t-elle un prix ? »} (2008, 2014, 2021 — dissertation)
    \item \textbf{« Peut-on sacrifier une personne pour sauver la société ? »} (2011, 2017 — dissertation / explication Kant \emph{Fondements} § sur le mensonge / Mounier \emph{Traité du caractère})
    \item \textbf{« La personnalité est-elle innée ou acquise ? »} (2005, 2013 — dissertation)
    \item \textbf{« Respecter la personne, est-ce respecter ses choix ? »} (2019 — dissertation)
    \item \textbf{Textes :} Kant (\emph{Fondements}, \emph{Anthropologie}), Mounier (\emph{Traité du caractère}, \emph{Le personnalisme}), Ricoeur (\emph{Soi-même comme un autre}), extraits de la \emph{Déclaration universelle des droits de l'homme}.
\end{itemize}
\textbf{Conseil :} Entraînez-vous sur ces sujets annales (disponibles sur OnBuch+). La maîtrise de la distinction \textbf{Prix / Dignité} (Kant) et \textbf{Individu / Personne} (Mounier) est la clé de 80\% des sujets.
\end{savaistu}

\begin{popfigure}
\centering
\begin{tikzpicture}[
    node distance=1.1cm and 0.5cm,
    stepbox/.style={draw=popblue, thick, rounded corners, fill=popblue!8, text width=2.8cm, align=center, minimum height=1.6cm, font=\small\bfseries, text=popdark},
    arrow/.style={-Stealth, thick, popdark},
    subbox/.style={font=\footnotesize, text=popmuted, align=center, text width=2.6cm}
]
% Main steps
\node[stepbox, align=center] (s1){1. COMPRENDRE\\ le sujet\\ (Termes, sens commun)};
\node[stepbox, right=of s1, align=center] (s2){2. PROBLÉMATISER\\ (Question philosophique,\\ tension)};
\node[stepbox, right=of s2, align=center] (s3){3. MOBILISER\\ Kant \& Mounier\\ (Thèses, concepts)};
\node[stepbox, right=of s3, align=center] (s4){4. ILLUSTRER\\ par un cas\\ camerounais};
\node[stepbox, right=of s4, align=center] (s5){5. CONCLURE\\ (Trancher, ouvrir)};
\draw[arrow] (s1) -- (s2);
\draw[arrow] (s2) -- (s3);
\draw[arrow] (s3) -- (s4);
\draw[arrow] (s4) -- (s5);
% Sub-details
\node[subbox, below=0.2cm of s1] {Définir : personne, prix, dignité, valeur};
\node[subbox, below=0.2cm of s2] {Prix (marché) vs Dignité (morale) ?};
\node[subbox, below=0.2cm of s3, align=center]{Kant : Fin en soi, Autonomie\\Mounier : Communion, Don};
\node[subbox, below=0.2cm of s4] {Village, lycée, traite, TikTok};
\node[subbox, below=0.2cm of s5] {Rappel probléma + Auteur + Ouverture};
\end{tikzpicture}
\legende{Schéma en escalier : la démarche argumentée type pour réussir la dissertation au Bac technique. Chaque marche est indispensable ; en sauter une fait chuter la note.}
\end{popfigure}

\courssub{Tableau récapitulatif général de la leçon}

\begin{popfigure}
\centering
\begin{tabularx}{\linewidth}{|l|X|X|X|}\hline
\rowcolor{popblueL}
\textbf{Notion clé} & \textbf{Définition essentielle} & \textbf{Auteur(s) de référence} & \textbf{Exemple concret (Cameroun / Bac)} \\ \hline
\cle{Personnalité} & Ensemble des traits psychiques, affectifs, moraux qui singularisent un individu ; \emph{construction} dynamique. & — (Psychologie, Sociologie) & Un élève timide devient leader associatif (socialisation, choix). \\ \hline
\cle{Facteurs biologiques} & Hérédité, tempérament, base neurophysiologique : le « donné » de départ. & \emph{Innisme} (descendance) & Jumeaux séparés : même tempérament, personnalités divergentes. \\ \hline
\cle{Facteurs sociaux} & Famille, école, pairs, classe : socialisation primaire / secondaire. & Durkheim, Bourdieu (habitus) & Éducation au village vs ville ; poids du groupe d'âge (tontine, associations). \\ \hline
\cle{Facteurs culturels} & Langue, valeurs, rites, médias, réseaux : héritage symbolique qui forme l'esprit. & Geertz, Mounier (milieu) & Respect des aînés, rites d'initiation, influence TikTok/Instagram. \\ \hline
\cle{Liberté / Choix} & Part de responsabilité : l'individu \emph{s'approprie} ses déterminismes, s'engage. & Sartre (existence précède essence), Mounier & Refus du mariage forcé, engagement associatif, conversion religieuse. \\ \hline
\cle{Personne} & Sujet moral, être spirituel, \emph{fin en soi}, doté de \textbf{dignité} (valeur absolue, incomparable). & \textbf{Kant}, \textbf{Mounier} & Tout être humain (même handicapé, criminel, fœtus, vieillard) = personne. \\ \hline
\cle{Dignité (Kant)} & Valeur \emph{intrinsèque}, \emph{inconditionnelle}, \emph{irremplaçable}. Fondée sur l'\textbf{autonomie} (se donner sa loi morale). & \textbf{Kant} (\emph{Fondements}) & Interdiction torture, esclavage, commerce organes, mensonge instrumentalisant. \\ \hline
\cle{Prix (Kant)} & Valeur \emph{conditionnelle}, \emph{comparable}, \emph{échangeable} (marchandise, travail, service). & \textbf{Kant} & Salaire, assurance, objet, animal — tout ce qui a un équivalent. \\ \hline
\cle{Personnalisme (Mounier)} & Philosophie de la \textbf{personne} contre l'individualisme et le collectivisme. Personne = \emph{communion}, \emph{don de soi}, \emph{ouverture au transcendant}. & \textbf{Mounier} (\emph{Traité du caractère}) & Engagement humanitaire, vie monastique, militantisme désintéressé, pardon. \\ \hline
\cle{Individu (Mounier)} & Atome social, être clos sur lui-même, mesurable, interchangeable, « persona » (masque). & \textbf{Mounier} & Consommateur passif, suiveur de masse, harceleur anonyme sur réseaux. \\ \hline
\cle{Aliénation (Mounier)} & Perte de la personne : réduction à l'état de chose, d'objet, de fonction ; dépossession de sa subjectivité. & \textbf{Mounier}, Marx & Traite des personnes, travail forcé, addiction algorithmique, propagande. \\ \hline
\cle{Autonomie / Hétéronomie} & Autonomie = se donner sa loi (dignité) ; Hétéronomie = subir une loi extérieure (prix, désir, aléa). & \textbf{Kant} & Élève qui triche (hétéronomie : loi de la note) vs élève qui étudie par devoir (autonomie). \\ \hline
\end{tabularx}
\legende{Tableau de révision complet : notions, définitions, auteurs, exemples. À connaître par cœur pour le Bac.}
\end{popfigure}

\begin{aretenir}[Check-list finale avant l'épreuve]
\begin{itemize}
    \item[$\square$] Je sais \textbf{distinguer} personnalité (construction, relative) et personne (valeur, absolue).
    \item[$\square$] Je peux \textbf{lister et illustrer} les 4 facteurs de la personnalité (bio, socio, culturel, libre).
    \item[$\square$] J'explique \textbf{Kant} : Dignité vs Prix, Fin en soi, Autonomie, Impératif catégorique (2 formulations).
    \item[$\square$] J'explique \textbf{Mounier} : Personne vs Individu, Communion vs Collectif, Aliénation, Personnalisme communautaire.
    \item[$\square$] Je sais \textbf{appliquer} ces notions à 4 situations types (âgés, harcèlement, traite, numérique).
    \item[$\square$] Je maîtrise la \textbf{méthode dissertation} (Intro probléma, Plan dialectique, Conclusion justifiée + ouverture).
    \item[$\square$] Je maîtrise la \textbf{méthode explication} (Situation, Thèse, Structure, Explication concepts+exemples, Discussion).
    \item[$\square$] Je connais les \textbf{critères de notation} et les \textbf{sujets classiques} annales.
\end{itemize}
\end{aretenir}

\textbf{Travail personnel :} Rédigez \emph{entièrement} une dissertation sur « La personne humaine a-t-elle un prix ? » en 3h30 chrono, puis auto-évaluez-vous avec la grille ci-dessus. Corrigez en vous reportant au \texttt{exoresolu} modèle. C'est le meilleur entraînement.

\newpage

\coursec{Activité d'intégration}

\coursec{Personnalité et personne}

\courssub{Objectifs de l'activité}
Cette activité d'intégration vise à mettre en œuvre les savoirs et savoir-faire de la leçon sur la personnalité et la personne à travers deux temps forts :
\begin{enumerate}
    \item \textbf{Analyser concrètement la genèse d'une personnalité} : identifier et hiérarchiser les facteurs biologiques, sociaux, culturels et les choix personnels à partir de portraits de jeunes Camerounais, et représenter cette analyse sous forme de diagramme circulaire.
    \item \textbf{Argumenter philosophiquement sur la valeur de la personne} : confronter les perspectives de Kant (dignité, fin en soi) et de Mounier (personnalisme, engagement, communauté) face à une situation d'exploitation du travail, et rédiger une conclusion justifiée évaluée sur des critères précis (notions, arguments, exemples, expression).
\end{enumerate}
L'objectif final est de permettre à l'élève de \cle{passer de la compréhension théorique à l'application critique} dans des contextes vécus, conformément aux attendus du Probatoire.

\courssub{Situation}
\textbf{Contexte :} Le Cameroun compte une population jeune (plus de 60\% ont moins de 25 ans). L'insertion professionnelle est un défi majeur : selon l'INS (Enquête ECAM 4), le taux de sous-emploi visible touche environ 70\% des jeunes actifs en milieu urbain. Cette précarité expose à des situations où la \cle{dignité} de la personne est mise à l'épreuve.

\textbf{Temps 1 -- Enquête guidée : « Trois visages de la jeunesse camerounaise »}
Chaque groupe de 4 à 5 élèves reçoit une \textbf{fiche portrait} documentée. Trois profils types sont distribués dans la classe :

\begin{center}
\begin{tabularx}{\linewidth}{|l|X|X|X|}\hline
\textbf{Critères} & \textbf{Profil A : Aliou, 18 ans, Apprenti mécanicien à Douala (New-Bell)} & \textbf{Profil B : Sylvie, 17 ans, Élève interne au Lycée de Ngaoundéré} & \textbf{Profil C : Marie, 19 ans, Vendeuse au marché Mokolo (Yaoundé)} \\ \hline
\textbf{Biologique} & Force physique, endurance, asthme léger (hérité du père). & Santé fragile (anémie), mémoire visuelle forte, grande taille (1,78 m). & Résistance à la fatigue, petite taille, vue perçante. \\ \hline
\textbf{Social (Famille/École/Entourage)} & Père chauffeur, mère couturière ; 4 frères/sœurs. Quitté l'école en 3e (échec BEPC). Maître d'apprentissage autoritaire mais protecteur. & Parents fonctionnaires (père instituteur, mère infirmière). Frère aîné médecin. Encadrement strict, valeurs d'excellence scolaire. & Orpheline de père, mère revendeuse de légumes. Aînée de 3 enfants. Réseau de « bayam-sellam » (revendeuses) solidaires mais concurrentes. \\ \hline
\textbf{Culturel (Croyances, normes, médias)} & Culture « rue » : respect de la hiérarchie d'atelier, codes vestimentaires (bleu de travail), football, musique urbaine (Coupé-décalé). Croyance en la « chance » et les « marabouts » pour le succès. & Culture lettrée : lecture, débats, église protestante (chorale). Normes : ponctualité, probité, ambition universitaire. Accès internet limité mais ciblé (cours en ligne). & Culture du « débrouillard » (Article 15), solidarité féminine (tontines), prières quotidiennes (catholique + traditions Bamileke). Influence des influenceurs TikTok (mode, entrepreneuriat). \\ \hline
\textbf{Choix personnels (Liberté/Responsabilité)} & Refus du vol de pièces détachées proposé par un camarade ; décision d'épargner 10\% du salaire hebdo pour ouvrir son propre garage plus tard. & Choix de la série D (Sciences) malgré la pression pour la série A ; engagement comme déléguée de classe ; bénévolat à l'infirmerie scolaire. & Refus de la prostitution occasionnelle proposée par une « grande sœur » ; choix de suivre une formation en ligne (marketing digital) le soir sur smartphone. \\ \hline
\end{tabularx}
\end{center}

\textbf{Temps 2 -- Débat argumenté : « Le piège de l'embauche »}
\textbf{Situation-problème commune :}
\begin{quote}
\textit{« Jean, 22 ans, titulaire d'un BTS en Maintenance Industrielle, cherche du travail depuis 14 mois. Une entreprise de BTP (sous-traitante d'un grand groupe) lui propose un poste de manœuvre sur un chantier de démolition à Douala-Bonabéri. Conditions : 12h/jour, 6 jours/7, 45 000 FCFA/mois (inférieur au SMIG de 60 000 FCFA), pas de contrat écrit, pas d'équipement de protection (casque, bottes, gants), manipulation d'amiante sans formation. Le chef d'équipe lui dit : \og{} Tu as de la chance, ici y'a la queue. Prends-le ou laisse, y'en a d'autres qui attendent. \fg{} Jean a une mère malade (diabète, dialyse 3 fois/semaine, 120 000 FCFA/mois non remboursés) et un petit frère en Terminale. Il signe le contrat oral le soir même. »}
\end{quote}

\courssub{Consignes}
\textbf{Partie A -- Enquête et Modélisation (Travail de groupe, 45 min)}
\begin{enumerate}
    \item \textbf{Compléter le tableau d'analyse} (fourni sur feuille A3) pour le profil attribué : recenser au minimum 3 éléments par colonne (Biologique, Social, Culturel, Choix personnels) en vous appuyant sur la fiche portrait.
    \item \textbf{Estimer le poids relatif} de chaque facteur dans l'édification de la personnalité du personnage (en \%). Justifier chaque estimation par une phrase (ex: « Social 40\% : l'autoritarisme du maître d'apprentissage a forgé sa discipline »).
    \item \textbf{Construire le diagramme circulaire} (camembert) sur papier millimétré ou via tableur : titre, légende, pourcentages, couleurs distinctes par facteur.
    \item \textbf{Synthèse orale (2 min/groupe)} : « Ce qui fait la singularité de [Prénom], c'est l'interaction entre... »
\end{enumerate}

\textbf{Partie B -- Débat Philosophique et Rédaction (Travail de groupe puis individuel, 55 min)}
\begin{enumerate}
    \setcounter{enumi}{4}
    \item \textbf{Préparation des camps (10 min)} : La classe est divisée en deux grands camps.
        \begin{itemize}
            \item \textbf{Camp Kant} : Défend l'idée que Jean \textbf{ne devait pas signer} (ou doit démissionner immédiatement). Argument : la dignité n'a pas de prix ; traiter l'homme comme un moyen (main-d'œuvre jetable) est immoral ; le devoir moral prime sur les conséquences utilitaires (soigner sa mère).
            \item \textbf{Camp Mounier} : Défend l'idée que le choix de Jean est \textbf{compréhensible voire justifié} comme engagement personnel. Argument : la personne se réalise dans la relation et l'histoire concrète ; le sacrifice pour sa mère (communauté de destin) est un acte d'amour qui \textit{personnalise} ; il faut transformer la situation de l'intérieur (syndicat, solidarité) plutôt que de fuir dans l'abstraction.
        \end{itemize}
    \item \textbf{Confrontation structurée (15 min)} : Chaque camp désigne 2 orateurs. Échanges dirigés par le professeur (modérateur). Interdiction les arguments d'autorité sans référence textuelle.
    \item \textbf{Rédaction individuelle finale (30 min)} : Rédiger une \textbf{conclusion justifiée de 10 lignes maximum} (environ 100 mots) répondant à la question : \og{} \textbf{Dans cette situation, la dignité de Jean est-elle bafouée ou assumée ?} \fg{} La réponse doit mobiliser \textbf{au moins une thèse de Kant et une de Mounier}, citer un \textbf{exemple précis du cas Jean}, et utiliser les concepts : \cle{dignité}, \cle{moyen/fin}, \cle{personnalisme}, \cle{engagement}.
\end{enumerate}

\courssub{Ressources à mobiliser}
\begin{definition}[Personnalité]
Ensemble des traits psychologiques, affectifs, intellectuels et comportementaux qui caractérisent un individu et le rendent unique. Elle résulte de l'interaction entre l'inné (biologique) et l'acquis (social, culturel) médiatisée par la liberté.
\end{definition}

\begin{definition}[Personne (Kant)]
Être raisonnable, doué de liberté et d'autonomie, qui possède une \cle{dignité} (valeur intrinsèque, incomparable) et non un \cle{prix} (valeur marchande, comparable). \textbf{Impératif catégorique (2e formulation)} : \og{} Agis de telle sorte que tu traites l'humanité, aussi bien dans ta personne que dans la personne de tout autre, toujours en même temps comme une fin, et jamais simplement comme un moyen. \fg{}
\end{definition}

\begin{definition}[Personnalisme (Mounier -- \textit{Traité du caractère}, \textit{Personnalisme})]
La personne n'est pas une substance close (l'individu), mais un \cle{être en relation}, ouverte sur le monde et les autres. Elle s'édifie par l'\cle{engagement} libre et responsable dans une \cle{communauté} (famille, nation, humanité). Refus du personnalisme individualiste (bourgeois) et du collectivisme (étatiste). \og{} La personne est le plus de l'individu. \fg{}
\end{definition}

\begin{propriete}[Facteurs de la personnalité]
\begin{itemize}
    \item \textbf{Biologiques (Inné) :} Tempérament, hérédité génétique, sexe, santé, aptitudes natives. Condition nécessaire mais non suffisante.
    \item \textbf{Sociaux (Acquis) :} Famille (style éducatif), école, groupe de pairs, classe sociale, monde professionnel. Fournissent les normes, valeurs, langage.
    \item \textbf{Culturels (Acquis) :} Croyances, religions, traditions, médias, arts, idéologies. Donnent du sens, des modèles, une vision du monde.
    \item \textbf{Choix personnels (Liberté) :} Décisions conscientes, engagements, refus, projets de vie. La personnalité n'est pas \textit{donnée}, elle se \textit{fait} (Sartre : \og{} L'existence précède l'essence \fg{} -- hors programme mais utile comme repère).
\end{itemize}
\end{propriete}

\courssub{Démarche et corrigé commenté}
\begin{exoresolu}[Corrigé type et commentaire pédagogique]
\textbf{Énoncé de l'activité (résumé) :} Analyser les facteurs de la personnalité pour trois profils de jeunes Camerounais (Partie A) et trancher le dilemme éthique du cas Jean en confrontant Kant et Mounier (Partie B).

\tcblower

\textbf{PARTIE A -- Exemple de traitement pour le Profil A (Aliou, Apprenti mécanicien)}

\textbf{1. Tableau d'analyse complété (extrait) :}
\begin{center}
\begin{tabularx}{\linewidth}{|l|X|}\hline
\textbf{Facteur} & \textbf{Éléments identifiés (min 3)} \\ \hline
Biologique & Force physique (manutention moteurs), Endurance (debout 10h/jour), Asthme léger (limite peinture/solvants). \\ \hline
Social & Modèle père (chauffeur = mobilité, effort), Fratrie nombreuse (partage, responsabilité), Maître d'apprentissage (autorité technique, transmission orale), Échec BEPC (rupture scolaire, stigmate). \\ \hline
Culturel & Code honneur « rue » (fidélité, non-dénonciation), Croyance chance/marabout (extériorisation réussite), Football (cohésion groupe), Esthétique bleu de travail (identité pro). \\ \hline
Choix personnels & Refus vol (intégrité), Épargne 10\% (projet autonomie), Acceptation discipline dure (stratégie apprentissage). \\ \hline
\end{tabularx}
\end{center}

\textbf{2. Estimation des poids (exemple d'élève) :}
\begin{itemize}
    \item Biologique : 15\% (Base physique nécessaire mais commune).
    \item Social : 45\% (L'atelier EST son univers formateur ; le maître remplace l'école ; la famille pousse à l'apprentissage).
    \item Culturel : 20\% (Les codes de la rue et la foi en la chance structurent sa vision du succès).
    \item Choix personnels : 20\% (C'est le facteur discriminant : deux apprentis même atelier $\neq$ même personnalité grâce aux choix éthiques/projet).
\end{itemize}
\textit{Justification pédagogique :} Le total fait 100\%. L'élève montre que le social domine (milieu technique oral), mais que la liberté (choix) n'est pas nulle (20\%), ce qui valide la thèse : \og{} La personnalité se fait dans la situation par la liberté \fg{}.

\textbf{3. Diagramme circulaire :} Secteurs proportionnels aux pourcentages. Couleurs : Vert (Bio), Bleu (Social), Orange (Culturel), Rouge (Choix). Légende claire. Titre : \og{} Facteurs édifiant la personnalité d'Aliou (Apprenti, Douala) \fg{}.

\vspace{0.5cm}
\textbf{PARTIE B -- Corrigé type de la conclusion justifiée (10 lignes / $\approx$ 100 mots)}

\begin{quote}
« La dignité de Jean est \textbf{bafouée objectivement} mais \textbf{assumée subjectivement} comme engagement. 

D'un côté, \textbf{Kant} condamne l'employeur qui traite Jean \textbf{uniquement comme un moyen} (main-d'œuvre précaire, amiante, salaire $<$ SMIG) et non comme une \textbf{fin en soi} : la dignité « n'a pas de prix » et ne saurait être monnayée contre 45 000 FCFA, fut-ce pour sauver sa mère. Jean perd son autonomie morale en acceptant l'inacceptable.

De l'autre, \textbf{Mounier} éclaire le \textbf{personnalisme} de Jean : son « oui » n'est pas soumission mais \textbf{engagement lucide} par amour pour sa mère (communauté de destin). Il personnalise sa condition en transformant une contrainte en don libre. 

\textbf{Synthèse} : La situation \textit{structurellement} nie sa dignité (Kant) ; l'acte de Jean \textit{existentiellement} la réaffirme par le sacrifice relationnel (Mounier). L'urgence éthique reste de transformer les structures pour que le choix ne soit plus entre dignité et survie. »
\end{quote}

\textbf{Commentaire sur la grille d'évaluation (sur 10 pts) :}
\begin{center}
\begin{tabularx}{\linewidth}{|l|X|c|}\hline
\textbf{Critères} & \textbf{Indicateurs de réussite} & \textbf{Pts} \\ \hline
Notions (3 pts) & Dignité, Moyen/Fin (Kant), Personnalisme/Engagement/Communauté (Mounier) définis et utilisés à bon escient. & 3 \\ \hline
Arguments (3 pts) & Articulation claire : Kant = refus structurel / Mounier = sens subjectif. Pas de confusion individu/personne. & 3 \\ \hline
Exemples (2 pts) & Citation précise : « 45 000 FCFA », « amiante », « mère dialysée », « engagement lucide ». Ancrage camerounais réel. & 2 \\ \hline
Expression (2 pts) & 10 lignes $\pm$ 1, connecteurs logiques (D'un côté/De l'autre/Synthèse), vocabulaire précis, pas de HS. & 2 \\ \hline
\end{tabularx}
\end{center}
\end{exoresolu}

\courssub{Bilan de l'activité}
Cette activité a permis de mettre en évidence trois acquis majeurs pour le Probatoire :
\begin{enumerate}
    \item \textbf{La distinction opératoire Personnalité/Personne :} Les élèves ont constaté que la \cle{personnalité} (Aliou, Sylvie, Marie) est un \textbf{fait empirique, multiple, mesurable} (diagramme circulaire), fruit de déterminismes (bio, socio, culturel) et de liberté. À l'inverse, la \cle{personne} (Jean) est une \textbf{valeur normative, une, infinie}, qui commande le respect (Kant) ou appelle l'engagement (Mounier), indépendamment des facteurs sociologiques.
    \item \textbf{L'articulation description / prescription :} Le passage du tableau descriptif (Partie A) au débat normatif (Partie B) entraîne l'élève à ne pas confondre \og{} expliquer pourquoi Jean a signé \fg{} (sociologie/psychologie) et \og{} juger s'il a eu raison de signer \fg{} (philosophie morale).
    \item \textbf{La fécondité du dialogue des pensées :} Kant et Mounier ne s'annulent pas ; ils opèrent à deux niveaux. Kant pose la \textbf{ligne rouge infranchissable} (dignité intangible = critique sociale du travail précaire). Mounier donne du \textbf{sens à l'épreuve vécue} (personnalisation par l'engagement = résilience éthique). La synthèse finale (« bafouée objectivement / assumée subjectivement ») modélise la réponse attendue en dissertation : \og{} Thèse / Antithèse / Synthèse dépassante \fg{}.
\end{enumerate}
\textbf{Transférabilité :} Cette démarche (Enquête factorielle $\to$ Modélisation $\to$ Cas éthique $\to$ Confrontation auteurs $\to$ Rédaction critériée) est réutilisable pour toute leçon du programme (ex: \og{} Sujet et culture \fg{}, \og{} Science et vérité \fg{}, \og{} Politique et justice \fg{}).

\newpage

\coursec{Méthodes et exercices résolus}

\courssub{Méthodes et exercices résolus}

\begin{methode}[Analyser une situation concrète pour identifier les facteurs de la personnalité]
\textbf{Objectif :} Appliquer la distinction entre facteurs biologiques, sociaux et culturels à un cas réel (savoir-faire 1).
\begin{enumerate}
    \item \textbf{Lire la situation} en soulignant les éléments descriptifs (âge, milieu, événements, comportements).
    \item \textbf{Classifier chaque élément} dans l'un des trois registres :
    \begin{itemize}
        \item \textbf{Biologique (inné) :} hérédité, tempérament, sexe, constitution physique, santé, maturation nerveuse.
        \item \textbf{Social (acquis par interaction) :} famille, école, pairs, travail, classe sociale, normes, rôles sociaux.
        \item \textbf{Culturel (acquis par transmission symbolique) :} langue, religion, valeurs, traditions, médias, idéologies, systèmes éducatifs formels.
    \end{itemize}
    \item \textbf{Montrer l'interaction :} La personnalité n'est pas une somme, mais une \cle{synthèse originale} (Le Moigne). Expliquer comment le biologique est modelé par le social/culturel (ex: un tempérament colérique canalisé par l'éducation).
    \item \textbf{Conclure} en caractérisant la personnalité résultante (unique, historique, en devenir).
\end{enumerate}
\cle{Interactionnisme} : La personnalité émerge de la dialectique entre le « donné » (biologique) et le « construit » (socio-culturel).
\end{methode}

\begin{exoresolu}[Analyse des facteurs de personnalité : le cas de « Amadou »]
\textbf{Énoncé :} Amadou, 17 ans, élève en Première F3 au Lycée Technique de Douala, est issu d'un père charpentier et d'une mère commerçante au marché Sandaga. Grand, costaud, il a un tempérament naturellement calme et réfléchi (on dit de lui qu'il « ne s'énerve jamais »). Élevé dans le respect des traditions Bassa (rites d'initiation, respect des aînés), il pratique le football en club. Depuis l'obtention de son BEPC, il manifeste un fort désir de devenir ingénieur en génie civil pour « construire des ponts solides comme mon père fait des meubles solides ». Il refuse cependant de tricher aux examens, contrairement à certains camarades, arguant que « l'honneur de la famille est en jeu ».
\textbf{Question :} Identifiez et classez les facteurs biologiques, sociaux et culturels qui ont contribué à édifier la personnalité d'Amadou. Montrez comment ils interagissent.
\tcblower
\textbf{Solution détaillée :}

\textbf{1. Facteurs biologiques (le « donné » naturel) :}
\begin{itemize}
    \item \textbf{Constitution physique :} « Grand, costaud » — prédispositions génétiques (hérédité polygène) favorisant l'aptitude au sport et une image corporelle valorisante.
    \item \textbf{Tempérament :} « Naturellement calme et réfléchi » — base neurobiologique (système nerveux équilibré, faible réactivité émotionnelle innée). C'est le \cle{socle} sur lequel l'éducation va s'appuyer.
\end{itemize}

\textbf{2. Facteurs sociaux (l'acquis par interaction relationnelle) :}
\begin{itemize}
    \item \textbf{Famille (socialisation primaire) :} Parents charpentier/commerçante — transmission d'une \cle{éthique du travail} manuel, de la rigueur, de la valeur de l'effort concret. Modèle parental (père artisan) source d'identification professionnelle.
    \item \textbf{École/Lycée (socialisation secondaire) :} Milieu technique (F3) — acquisition de savoirs scientifiques, normes de la vie scolaire, compétition pour les notes, pression du Probatoire.
    \item \textbf{Groupe des pairs (club de football) :} Apprentissage de la coopération, de la compétition réglée, gestion de l'échec/réussite collective. Contraste avec les « camarades qui trichent » : groupe de référence négatif qui renforce sa propre norme.
    \item \textbf{Classe sociale :} Milieu populaire/artisanal — conscience de la nécessité de la réussite scolaire comme ascension sociale (ingénieur vs ouvrier).
\end{itemize}

\textbf{3. Facteurs culturels (l'acquis par transmission symbolique) :}
\begin{itemize}
    \item \textbf{Tradition Bassa :} Rites d'initiation, respect des aînés — intériorisation de valeurs \cle{communautaires} : honneur de la famille, solidarité, responsabilité collective. Le refus de tricher n'est pas seulement moral individuel, il est \emph{honorabilité lignagère}.
    \item \textbf{Langue et représentations :} Métaphore du « pont solide comme les meubles du père » — réappropriation culturelle du métier paternel sublimé par la technique moderne.
    \item \textbf{Valeurs républicaines/citoyennes (École) :} Mérite, honnêteté intellectuelle — renforcement de la norme anti-triche.
\end{itemize}

\textbf{4. Interaction et Synthèse (Le devenir personne) :}
\begin{itemize}
    \item Le \textbf{tempérament calme} (biologique) a facilité l'intériorisation des \textbf{exigences de réflexion} de l'initiation Bassa et de la rigueur technique (social/culturel). Sans ce tempérament, l'impulsivité aurait pu rendre l'apprentissage du métier d'ingénieur plus difficile.
    \item La \textbf{valeur « honneur familial »} (culturel) transforme une contrainte scolaire (ne pas tricher) en \textbf{impératif personnel existentiel}. Amadou ne subit pas la norme, il s'en \cle{approprie} le sens.
    \item Le \textbf{projet professionnel} (ingénieur) réalise une \textbf{synthèse originale} : il ne reproduit pas le père (charpentier), il \emph{sublime} le geste technique (bois $\to$ béton, artisanat $\to$ science) tout en gardant la \emph{finalité} (solidité, utilité sociale).
\end{itemize}
\textbf{Conclusion :} La personnalité d'Amadou n'est pas la somme Bio + Socio + Culturel. C'est une \textbf{structure dynamique} où le biologique offre des potentialités, le social offre des occasions et des modèles, le culturel offre du sens et des valeurs. Amadou \textbf{se fait personne} en unifiant ces apports dans un projet libre (devenir ingénieur honnête).
\end{exoresolu}

\begin{methode}[Distinguer rigoureusement « Personnalité » et « Personne »]
\textbf{Objectif :} Maîtriser la distinction conceptuelle fondamentale pour toute dissertation ou commentaire (savoir-faire 2).
\begin{enumerate}
    \item \textbf{Définir la Personnalité (l'ordre de l'avoir/être psychologique) :}
    \begin{itemize}
        \item Ensemble de traits stables (caractère, tempérament, intelligence, habitudes).
        \item Résultat de l'édification (facteurs bio/socio/culturels).
        \item Se \textbf{mesure}, se \textbf{compare}, se \textbf{détermine} par des causes (psychologie scientifique).
        \item C'est un \textbf{« quoi »} (un objet d'étude).
    \end{itemize}
    \item \textbf{Définir la Personne (l'ordre de la valeur/être moral/juridique) :}
    \begin{itemize}
        \item Sujet de droit, de devoirs, de responsabilité.
        \item \textbf{Fin en soi} (Kant), être d'\textbf{existence} et de \textbf{communion} (Mounier).
        \item Ne se mesure pas, ne se compare pas : \cle{dignité} (incomparable) vs \cle{prix} (comparable).
        \item C'est un \textbf{« qui »} (un sujet moral).
    \end{itemize}
    \item \textbf{Articuler :} La personnalité est le « véhicule » ou l'« instrument » de la personne. On \emph{a} une personnalité (plus ou moins riche, équilibrée) ; on \emph{est} une personne (absolument, inconditionnellement).
    \item \textbf{Reflexe Bac :} Ne jamais écrire « Sa personne est timide » (confusion). Écrire : « Sa personnalité présente un trait d'introversion ; en tant que personne, il mérite le respect. »
\end{enumerate}
\cle{Personnalité = Fait (psychologique, déterminé). Personne = Valeur (morale/juridique, libre).}
\end{methode}

\begin{exoresolu}[Distinction conceptuelle et pièges de langage]
\textbf{Énoncé :} Corrigez et justifiez les affirmations suivantes :
\begin{enumerate}
    \item « Le juge a tenu compte de la personne de l'accusé, qui est un grand colérique, pour atténuer la peine. »
    \item « Chaque homme a une personne unique, c'est pourquoi il faut la respecter. »
    \item « La personnalité de ce dirigeant est impressionnante : il a une dignité naturelle. »
\end{enumerate}
\tcblower
\textbf{Solution détaillée :}

\textbf{1. Correction :} « Le juge a tenu compte de la \textbf{personnalité} de l'accusé (son tempérament colérique, trait psychologique déterminé), non de sa \textbf{personne} (sa qualité de sujet de droit). »
\textbf{Justification :} La colère est un trait de \textbf{personnalité} (fait psychologique, relevant de la causalité naturelle). La \textbf{personne} est le fondement de la responsabilité pénale (liberté). Tenir compte de la personnalité pour atténuer la peine (circonstance atténuante : trouble psychique) ne nie pas la personne (qui reste responsable), mais module la sanction en fonction de la capacité de discernement.

\textbf{2. Correction :} « Chaque homme a une \textbf{personnalité} unique (ensemble de traits singuliers), c'est pourquoi, en tant que \textbf{personne}, il faut le respecter. »
\textbf{Justification :} L'unicité descriptive (empreinte digitale, histoire, caractère) relève de la \textbf{personnalité}. Le fondement du respect (dignité) relève de la \textbf{personne}. \textbf{Attention :} Ce n'est pas \emph{parce que} la personnalité est unique qu'on respecte la personne (sinon, un clone ou un homme « banal » ne mériterait pas le respect). On respecte la personne \emph{en raison de sa nature d'être raisonnable/liberté} (Kant) ou d'\emph{être spirituel} (Mounier), indépendamment de sa singularité psychologique.

\textbf{3. Correction :} « La \textbf{personnalité} de ce dirigeant est impressionnante (charisme, compétences, traits de caractère) ; il \textbf{incarne la dignité} de la personne (ou : il agit en personne responsable). »
\textbf{Justification :} « Impressionnante » qualifie un ensemble de qualités psychologiques/sociales (\textbf{personnalité}). La \textbf{dignité} n'est pas un trait de personnalité (on ne peut pas être « un peu digne » ou « très digne » comme on est « un peu timide » ou « très extraverti »). La dignité est l'attribut de la \textbf{personne} : elle est \cle{inconditionnelle} (on l'a ou on ne l'a pas, tout être humain la possède). Dire « il a une dignité naturelle » risque de confondre dignité (valeur morale intrinsèque) et prestance naturelle (trait de personnalité).
\end{exoresolu}

\begin{methode}[Expliquer et argumenter la thèse de Kant : La personne comme fin en soi (Dignité)]
\textbf{Objectif :} Restituer l'argumentation kantienne et l'appliquer à un cas d'éthique appliquée (savoir-faire 1 \& 2).
\begin{enumerate}
    \item \textbf{Rappeler le principe suprême (Impératif Catégorique, 2\textsuperscript{e} formulation) :} « Agis de telle sorte que tu traites l'humanité, aussi bien dans ta personne que dans la personne de tout autre, toujours en même temps comme une fin, et jamais simplement comme un moyen. »
    \item \textbf{Distinguer \cle{Prix} et \cle{Dignité} :}
    \begin{itemize}
        \item \textbf{Prix (Valeur marchande/conditionnelle) :} Ce qui a un équivalent, peut être remplacé, échangé, utilisé comme \textbf{moyen} (objets, compétences, travail, personnalité).
        \item \textbf{Dignité (Valeur intrinsèque/inconditionnelle) :} Ce qui n'a pas d'équivalent, ne peut être remplacé, ne doit \textbf{jamais} être utilisé comme simple moyen. C'est l'attribut de la \textbf{Personne}.
    \end{itemize}
    \item \textbf{Fonder la dignité :} La personne est \textbf{être raisonnable} (capacité d'autonomie, de se donner sa propre loi morale). Elle est \textbf{fin en soi} car elle est la source de la valeur (sans sujet rationnel, rien n'a de valeur).
    \item \textbf{Appliquer (Cas concret) :} Identifier si une action traite l'humain \textbf{aussi} comme fin (respect de son consentement, de ses fins propres, information, autonomie) ou \textbf{uniquement} comme moyen (manipulation, contrainte, tromperie, instrumentalisation totale).
    \item \textbf{Formuler la conclusion :} « L'action X respecte/violente la dignité car elle [traite/ne traite pas] l'agent comme une fin en soi, en [respectant/violant] son autonomie. »
\end{enumerate}
\cle{Autonomie = Capacité d'être sa propre fin. Hétéronomie = Être moyen pour une fin extérieure.}
\end{methode}

\begin{exoresolu}[Application de l'éthique kantienne : Le stage obligatoire non rémunéré]
\textbf{Énoncé :} Une entreprise de BTP à Yaoundé impose aux élèves de Terminale Génie Civil un stage de 3 mois non rémunéré, sous peine de ne pas valider leur module « Stage Professionnel ». L'entreprise utilise ces élèves comme main-d'œuvre pour des tâches répétitives (port de charges, nettoyage) sans encadrement pédagogique. Le directeur argue : « Cela fait partie de leur formation, ils acquièrent de l'expérience, c'est un échange de bon procédé. »
\textbf{Question :} En vous appuyant sur la distinction kantienne Prix/Dignité, analysez si cette pratique respecte la personne des élèves.
\tcblower
\textbf{Solution détaillée :}

\textbf{1. Identification du statut des élèves :}
Les élèves sont des \textbf{personnes} (êtres raisonnables, fins en soi, doués de dignité). Ils sont aussi des \textbf{personnalités} en formation (acquisition de compétences, prix sur le marché du travail).

\textbf{2. Analyse de la proposition du directeur (Thèse de l'échange) :}
Le directeur présente la relation comme un \textbf{contrat d'échange} : Travail (moyen) $\leftrightarrow$ Expérience/Validation (fin).
\begin{itemize}
    \item \textbf{Niveau « Prix » :} L'élève vend sa \textbf{force de travail} (compétence, temps, énergie) — cela a un \textbf{prix} (salaire normalement). L'entreprise « achète » cette force de travail. L'élève « achète » l'expérience. C'est une logique marchande, utilitariste.
    \item \textbf{Problème :} L'échange est \textbf{inégal} (contrainte de validation du diplôme $\to$ absence de liberté contractuelle réelle). L'élève \emph{doit} accepter pour obtenir son diplôme. Il n'y a pas \textbf{consentement libre} (condition de l'autonomie).
\end{itemize}

\textbf{3. Test de l'Impératif Catégorique (Formule de l'Humanité) :}
\begin{itemize}
    \item \textbf{Moyen :} L'entreprise utilise la force physique des élèves (port de charges, nettoyage) comme \textbf{moyen} pour son profit (économie de main-d'œuvre qualifiée).
    \item \textbf{Fin ?} L'élève est-il traité \textbf{aussi} comme une fin ?
    \begin{itemize}
        \item \textbf{Non} : Absence d'encadrement pédagogique $\to$ pas de transmission de savoir (non-respect de sa fin propre : \emph{apprendre}).
        \item \textbf{Non} : Non-rémunération de travail productif $\to$ exploitation (l'élève est \emph{simple} moyen de production).
        \item \textbf{Non} : Contrainte (chantage au diplôme) $\to$ violation de l'\textbf{autonomie} (l'élève ne se donne pas sa propre fin, il subit une fin étrangère).
    \end{itemize}
\end{itemize}

\textbf{4. Argumentation kantienne rigoureuse :}
\begin{itemize}
    \item \textbf{Violation de la Dignité :} L'élève est traité \textbf{uniquement comme un moyen} (main-d'œuvre corvéable à merci).
    \item \textbf{Instrumentalisation :} La « formation » est un \textbf{prétexte} (masque) pour masquer l'exploitation. Kant dit : « Celui qui veut se servir de l'homme comme d'un moyen doit en même temps contenir en lui la fin. » Ici, la fin de l'élève (apprendre, valider son diplôme par la compétence) est niée par la réalité des tâches (exécution sans apprentissage).
    \item \textbf{Respect de la personnalité $\neq$ Respect de la personne :} L'entreprise « respecte » la personnalité de l'élève (sa force physique, sa docilité, son besoin de diplôme) pour s'en servir. Elle \textbf{nie} sa personne (sa liberté, sa rationalité, son droit à un contrat équitable).
\end{itemize}

\textbf{Conclusion :} Cette pratique \textbf{viole la dignité} des élèves. Elle les traite comme des \textbf{objets} (prix, main-d'œuvre) et non comme des \textbf{sujets} (fins en soi, autonomes). Une pratique respectueuse impliquerait : convention de stage claire, encadrement pédagogique effectif, tâches formatrices, et idéalement gratification (reconnaissance de la valeur du travail fourni), le tout accepté \textbf{librement}.
\end{exoresolu}

\begin{methode}[Expliquer et argumenter la thèse de Mounier : Le Personnalisme (Existence et Communion)]
\textbf{Objectif :} Restituer la spécificité du personnalisme mouniérien face au kantisme et l'appliquer (savoir-faire 1 \& 2).
\begin{enumerate}
    \item \textbf{Situer l'opposition :} Contre l'\textbf{individualisme} (atome social, droits subjectifs, égoïsme) et le \textbf{collectivisme} (individu noyé dans la masse, État totalitaire). Pour la \textbf{Personne}.
    \item \textbf{Définir la Personne (Mounier) :} « Être en situation », « Être d'existence », « Être de communion ».
    \begin{itemize}
        \item \textbf{Existence (vs Essence) :} La personne ne se \emph{définit} pas (essence close), elle \emph{s'invente} par des choix libres, des engagements, des réponses à l'appel de la situation. Elle est \textbf{devenir}.
        \item \textbf{Communion (vs Communication/Collection) :} Relation \textbf{sujet-sujet} (I-Thou / Je-Tu). Don de soi, réciprocité, amour, amitié, engagement militant. \textbf{Pas} relation sujet-objet (utilisation) ni fusion (perte de soi).
        \item \textbf{Intériorité spirituelle :} Cœur de la personne, lieu de la liberté, de la réflexion, de la prière/meditation. Source de l'unité.
    \end{itemize}
    \item \textbf{Distinction cruciale Personnalité / Personne (Mounier) :}
    \begin{itemize}
        \item \textbf{Personnalité} = « Masque », « Enveloppe », ensemble des traits psychologiques, sociaux, héréditaires. C'est le \textbf{côté objectif, mesurable, déterminé}. On \emph{a} une personnalité.
        \item \textbf{Personne} = « Noyau spirituel », « Centre unifiant », le \emph{Je} profond qui dit « Je » et s'engage. C'est le \textbf{côté subjectif, libre, infini}. On \emph{est} une personne.
    \end{itemize}
    \item \textbf{Appliquer (Cas concret) :} Repérer si une situation favorise l'\textbf{éveil de la personne} (prise de responsabilité, engagement, créativité, communion) ou si elle enferme dans la \textbf{personnalité} (conformisme, conditionnement, rôle social figé, « masque »).
\end{enumerate}
\cle{Personnalité = Ce que je \emph{suis} (donné). Personne = Ce que je \emph{fais} de ce que je suis (liberté).}
\end{methode}

\begin{exoresolu}[Personnalisme et engagement associatif : Le cas de « Aïcha »]
\textbf{Énoncé :} Aïcha, 18 ans, élève en Première STT (Secrétariat), est très timide (personnalité introvertie). Son père souhaite qu'elle épouse un cousin au village après le Probatoire, « comme la tradition l'exige ». Aïcha refuse. Elle s'engage dans une association de jeunes « Jeunes Leaders Cameroun » où elle anime des ateliers de prise de parole en public. Elle dit : « Je découvre que je peux agir, que ma voix compte. Je ne suis plus seulement la fille de mon père ou la future épouse, je suis Aïcha. »
\textbf{Question :} Analysez cette situation à la lumière de la distinction mouniérienne entre Personnalité et Personne, et des notions d'Existence et de Communion.
\tcblower
\textbf{Solution détaillée :}

\textbf{1. Le pôle « Personnalité » (Le donné, le masque, le conditionnement) :}
\begin{itemize}
    \item \textbf{Traits psychologiques :} Timidité, introversion (tempérament biologique).
    \item \textbf{Conditionnement social/culturel :} Rôle assigné : « Fille de... », « Future épouse de... ». Attente familiale (père), norme traditionnelle (mariage forcé/arrangé). C'est la \textbf{personnalité « close »}, déterminée par l'extérieur. Aïcha \emph{subit} cette personnalité comme un destin.
    \item \textbf{Masque social :} L'identité « fille soumise » est un rôle (persona) que la société/la famille lui colle sur le visage.
\end{itemize}

\textbf{2. Le pôle « Personne » (L'éveil, la liberté, l'existence) :}
\begin{itemize}
    \item \textbf{Existence (S'inventer) :} Le refus du mariage imposé est un \textbf{acte d'existence}. Aïcha sort de l'« essence » (fille soumise) pour entrer dans l'\textbf{existence} (sujet libre qui dit « Non »). Elle \textbf{se choisit}.
    \item \textbf{Engagement (Communion) :} L'association « Jeunes Leaders » n'est pas un simple club (collection d'individus). C'est un lieu de \textbf{communion} :
    \begin{itemize}
        \item Relation \textbf{Sujet-Sujet} : Animateurs/animés s'écoutent, se respectent, grandissent ensemble.
        \item \textbf{Don de soi :} Aïcha donne son temps, son énergie, sa parole pour les autres.
        \item \textbf{Réciprocité :} Elle reçoit confiance, compétences, reconnaissance. « Ma voix compte. »
    \end{itemize}
    \item \textbf{Intériorité :} « Je découvre que je peux agir » $\to$ Éveil du \textbf{noyau spirituel}. La timidité (personnalité) n'a pas disparu, mais la \textbf{personne} (le Je profond) ne s'y \emph{identifie} plus. Elle \emph{a} une personnalité timide, mais elle \emph{est} une personne audacieuse.
\end{itemize}

\textbf{3. Synthèse Mouniérienne :}
\begin{itemize}
    \item \textbf{Personnalité} = « Fille timide promise au cousin » (Fait, déterminisme, passé).
    \item \textbf{Personne} = « Aïcha, animatrice, citoyenne engagée » (Valeur, liberté, avenir).
    \item La \textbf{personnalité} est le \textbf{matériau} (l'argile) ; la \textbf{personne} est le \textbf{sculpteur} (l'esprit). Aïcha \textbf{personnalise} sa personnalité : elle assume sa timidité (elle ne devient pas extravertie magiquement) mais elle \textbf{la dépasse} par l'engagement.
    \item \cle{Personnalisation} = Processus par lequel la personne s'approprie sa personnalité et l'oriente vers des valeurs (liberté, justice, vérité).
\end{itemize}
\textbf{Conclusion :} Aïcha illustre parfaitement la thèse de Mounier : la personne n'est pas une substance figée (personnalité), mais un \textbf{acte permanent de liberté} (existence) qui se réalise dans la \textbf{communion} (engagement associatif). Elle refuse d'être un « objet » de tradition pour devenir « sujet » de sa vie.
\end{exoresolu}

\begin{methode}[Comparer Kant et Mounier sur la valeur de la personne (Dissertation/Commentaire)]
\textbf{Objectif :} Structurer une comparaison philosophique rigoureuse (savoir-faire 2).
\begin{enumerate}
    \item \textbf{Identifier le point commun (Le socle) :} Tous deux affirment la \textbf{valeur absolue, inconditionnelle, incomparable} de la personne (Dignité chez Kant, Sacralité/Invulnérabilité chez Mounier). Contre l'instrumentalisation (moyen), le relativisme, le totalitarisme.
    \item \textbf{Structurer les différences selon 3 axes :}
    \begin{center}
    \begin{tabularx}{\linewidth}{|l|X|X|}\hline
    \rowcolor{popblueL} \textbf{Axe} & \textbf{Kant (Rationalisme critique)} & \textbf{Mounier (Existentialisme personnaliste)} \\ \hline
    \textbf{Fondement} & \textbf{Raison pure / Autonomie législative.} La personne est \textbf{sujet de droit moral} (nouménal). & \textbf{Existence concrète / Esprit incarné.} La personne est \textbf{être en situation}, corps et histoire. \\ \hline
    \textbf{Structure} & \textbf{Dualisme} : Personne (noumène/liberté) \textbf{vs} Personnalité (phénomène/nature). Séparation nette. & \textbf{Unité dialectique} : Personne (noyau) \textbf{personnalise} la Personnalité (masque). Continuité. \\ \hline
    \textbf{Relation à l'Autre} & \textbf{Respect (Devoir) :} Reconnaissance mutuelle de la dignité (Droit, Justice). Relation \textbf{universelle}, formelle. & \textbf{Communion (Amour/Amitié/Engagement) :} Don de soi, réciprocité concrète. Relation \textbf{singulière}, existentielle. \\ \hline
    \textbf{Réalisation} & \textbf{Devoir être} (Idéal régulateur). Agir par devoir, selon la maxime universalisable. & \textbf{Devenir} (Processus historique). S'engager, créer, répondre à l'appel de la situation. \\ \hline
    \textbf{Rapport au corps/histoire} & \textbf{Indifférent / Obstacle :} Le corps/appartenances sont du côté de la nature (phénomène), pas de la personne (noumène). & \textbf{Constituant :} Le corps, l'histoire, la culture, le milieu sont la \textbf{matière} même de la personnalisation. \\ \hline
    \end{tabularx}
    \end{center}
    \item \textbf{Synthèse dialectique (Niveau Bac « Excellent ») :} Kant garantit le \textbf{droit} (protection contre l'arbitraire, universalité). Mounier garantit la \textbf{vie} (épanouissement, singularité, concrétude). Une société juste a besoin des deux : \textbf{Kant pour les lois (Justice), Mounier pour la culture/éducation (Vie bonne)}.
    \item \textbf{Rédaction :} Utiliser des connecteurs logiques : \emph{En revanche, Par conséquent, Dès lors, Certes... mais.}
\end{enumerate}
\cle{Kant = Personne = Sujet de droit (Universalisme formel). Mounier = Personne = Sujet d'existence (Singularité concrète).}
\end{methode}

\begin{exoresolu}[Sujet de dissertation type : « La personne est-elle réductible à sa personnalité ? »]
\textbf{Énoncé :} Sujet de dissertation : \textbf{« La personne est-elle réductible à sa personnalité ? »}
\textbf{Consigne :} Proposez un plan détaillé (Introduction, 2 ou 3 parties avec arguments/auteurs/exemples, Conclusion) en mobilisant Kant et Mounier.
\tcblower
\textbf{Solution détaillée : Plan de dissertation structuré}

\textbf{INTRODUCTION}
\begin{itemize}
    \item \textbf{Accroche :} Cit. Sénèque : « L'homme est une chose sacrée pour l'homme. » Ou distinction courante : « Il a une forte personnalité » vs « C'est une personne bien ».
    \item \textbf{Définitions :} \textbf{Personnalité} = Ensemble de traits psychologiques, sociaux, culturels (le « quoi », le masque, le déterminé). \textbf{Personne} = Sujet moral, être de valeur, de liberté, de responsabilité (le « qui », le noyau, l'absolu).
    \item \textbf{Problématique :} Si la personnalité est l'enveloppe observable, la personne en est-elle le simple contenu ? La réduction de la personne à la personnalité ne risque-t-elle pas de nier la liberté et la dignité (Kant) ou l'existence et la communion (Mounier) ?
    \item \textbf{Annonce du plan :} Nous verrons d'abord que la personnalité est le support nécessaire mais insuffisant de la personne (I), puis que la personne excède radicalement la personnalité par sa dimension morale (Kant) et existentielle (Mounier) (II), pour conclure sur leur articulation dialectique (III).
\end{itemize}

\textbf{DÉVELOPPEMENT}

\textbf{I. La personnalité : condition nécessaire, support empirique de la personne (Le « matériel »)}
\begin{itemize}
    \item \textbf{A. Facteurs d'édification (Programme) :} Biologiques (tempérament), Sociaux (famille, école), Culturels (valeurs, langue). La personnalité est le \textbf{produit} de ces déterminismes. Ex: Amadou (exercice 1).
    \item \textbf{B. Fonction opératoire :} La personnalité fournit les \textbf{moyens} d'agir (intelligence, force, langage, compétences, habitudes). Sans personnalité (ex: coma profond, handicap mental lourd), la personne peine à s'exprimer (problème du « végétatif »).
    \item \textbf{C. Risque de confusion (Réductionnisme) :} Psychologisme (Freud : « Le moi n'est pas maître en sa maison »), Sociologisme (Durkheim : individu déterminé par faits sociaux), Biologisme (neurosciences : « Nous sommes notre cerveau »). \textbf{Thèse adverse :} La personne \emph{est} sa personnalité (matérialisme).
\end{itemize}

\textbf{II. La personne : excédent irréductible à la personnalité (Le « formel » / Le « spirituel »)}
\begin{itemize}
    \item \textbf{A. Chez Kant : La personne comme Fin en soi (Nouménal).}
    \begin{itemize}
        \item Distinction radicale \textbf{Phénomène / Noumène}. Personnalité = Phénomène (soumis aux lois naturelles, causalité, prix). Personne = Noumène (liberté, autonomie, dignité).
        \item \textbf{Argument :} Je ne peux pas \emph{choisir} ma personnalité (donnée), mais je \emph{dois} me comporter en personne (devoir moral). La dignité n'a pas de prix (incomparable), la personnalité a un prix (échangeable).
        \item \textbf{Exemple :} Un criminel a une personnalité « perverse » (fait), mais il reste une personne (droit au procès équitable, interdiction de la torture). On juge l'acte (personnalité), on respecte l'être (personne).
    \end{itemize}
    \item \textbf{B. Chez Mounier : La personne comme Existence et Communion (Spirituel incarné).}
    \begin{itemize}
        \item Distinction \textbf{Personnalité (Masque) / Personne (Noyau)}. La personnalité est « ce que je \emph{suis} » (donné) ; la personne est « ce que je \emph{fais} de ce que je suis » (liberté).
        \item \textbf{Argument :} \textbf{Personnalisation} = Mouvement par lequel le noyau spirituel s'approprie le masque. La personne \textbf{n'est pas}, elle \textbf{devient} (Existence). Elle se réalise dans la \textbf{Communion} (don de soi), impossible pour une simple personnalité (objet).
        \item \textbf{Exemple :} Aïcha (exercice 4). Timide (personnalité) $\to$ Engagée (personne). La personnalité ne \emph{détermine} pas la personne ; la personne \emph{transfigure} la personnalité.
    \end{itemize}
\end{itemize}

\textbf{III. Articulation dialectique : La personne \emph{personnalise} sa personnalité (Synthèse)}
\begin{itemize}
    \item \textbf{A. Pas de personne sans personnalité (Incarnation) :} Mounier contre Kant : La personne n'est pas un « fantôme » nouménal pur. Elle a \textbf{besoin} de la personnalité (corps, histoire, culture) pour s'exprimer. « L'esprit a besoin de la matière pour se manifester. »
    \item \textbf{B. Pas de personnalité humaine sans personne (Unification) :} Une personnalité sans personne (schizophrénie, aliénation totale, conformisme moutonnier) est une \textbf{personnalité « close », fragmentée, non unifiée}. C'est la personne (le « Je » unifiant) qui donne du \textbf{sens} et de l'\textbf{unité} aux traits dispersés.
    \item \textbf{C. Éthique de la personnalisation :} L'idéal n'est pas de \emph{nier} sa personnalité (ascétisme kantien mal compris) ni de s'y \emph{complaire} (narcissisme), mais de la \textbf{personnaliser} : assumer ses traits (timidité, force, culture) pour les mettre au service de valeurs librement choisies (justice, vérité, amour).
\end{itemize}

\textbf{CONCLUSION}
\begin{itemize}
    \item \textbf{Bilan :} Non, la personne n'est pas réductible à sa personnalité. La personnalité est le \textbf{nécessaire} (condition), la personne est la \textbf{fin} (le supérieur en dignité).
    \item \textbf{Ouverture :} Cette distinction fonde nos sociétés : \textbf{Droits de l'homme} (Kant : protection de la personne contre l'État/marché qui ne voient que des « personnalités » productives) et \textbf{Éducation/Santé mentale} (Mounier : soigner/éduquer pour libérer la personne derrière le symptôme ou le rôle social).
\end{itemize}
\end{exoresolu}

\begin{methode}[Rédiger une introduction et une conclusion « type Bac » en Philosophie]
\textbf{Objectif :} Maîtriser le formalisme exigé à l'écrit du Probatoire/Baccalauréat (savoir-faire 2).
\begin{enumerate}
    \item \textbf{Introduction (Entonnoir : 4 étapes obligatoires) :}
    \begin{enumerate}
        \item \textbf{Accroche (A) :} Citation d'auteur \textbf{au programme} (Kant, Mounier, Descartes, Aristote...), proverbe, fait social précis, définition étymologique. \textit{Interdit :} « Depuis la nuit des temps... », « De tout temps les hommes... ».
        \item \textbf{Définition des termes (D) :} Définir \textbf{tous} les termes clés du sujet (sens commun + sens philosophique). Distinguer les concepts (ex: Personnalité vs Personne).
        \item \textbf{Problématisation (P) :} Transformer le sujet en \textbf{question philosophique} tensionnelle (opposition de thèses). « Si [thèse 1], alors [conséquence], mais [thèse 2] suggère [autre chose]. \textbf{Dès lors, [Question centrale] ?} »
        \item \textbf{Annonce du plan (A) :} « Pour répondre, nous examinerons d'abord... (I), puis... (II), avant de montrer... (III). » \textit{Verbes d'action :} Analyser, Montrer, Démontrer, Étudier, Interroger.
    \end{enumerate}
    \item \textbf{Conclusion (Entonnoir inverse : 3 étapes) :}
    \begin{enumerate}
        \item \textbf{Bilan / Synthèse (B) :} Reprendre la \textbf{réponse directe} à la problématique en résumant l'apport de chaque partie (I a montré que... ; II a révélé que... ; III a synthétisé...). \textit{Pas de copier-coller de l'intro.}
        \item \textbf{Réponse définitive (R) :} Une phrase claire, tranchée, qui clôt le débat. « En définitive, la personne n'est pas réductible à sa personnalité car... »
        \item \textbf{Ouverture (O) :} Élargir le horizon : autre auteur (pas au programme), autre discipline (Droit, Sciences politiques, Littérature), actualité brûlante, question inverse. \textit{Interdit :} « Pour conclure, on peut se demander... » (relance le devoir). L'ouverture \textbf{ferme} la porte en l'entrouvrant sur un ailleurs.
    \end{enumerate}
    \item \textbf{Reflexes de rédaction :} \textbf{Un paragraphe = Une idée = Un argument.} Connecteurs logiques (\emph{Premièrement, En outre, Cependant, Par conséquent, En somme}). \textbf{Zéro faute d'orthographe sur les noms d'auteurs et concepts.} Majuscule à \textbf{Personne, Dignité, Existence, Communion, Autonomie, Impératif Catégorique}.
\end{enumerate}
\cle{ADPA (Accroche, Définitions, Problématique, Annonce) / BRO (Bilan, Réponse, Ouverture).}
\end{methode}

\begin{exoresolu}[Rédaction guidée : Introduction et Conclusion sur un sujet inédit]
\textbf{Énoncé :} Sujet : \textbf{« Le respect de la personne exige-t-il l'effacement de la personnalité ? »}
\textbf{Consigne :} Rédigez \textbf{uniquement} l'Introduction et la Conclusion complètes, prêtes à être copiées sur la copie.
\tcblower
\textbf{Solution détaillée :}

\textbf{INTRODUCTION}
\begin{quote}
\textbf{[Accroche]} « La personnalité est le masque que nous portons ; la personne est le visage que nous sommes », pourrait-on paraphraser en écho à \textbf{Emmanuel Mounier} (\textit{Traité du caractère}, \textit{Le Personnalisme}). 

\textbf{[Définitions]} La \textbf{personnalité} désigne l'ensemble des traits psychologiques (tempérament, caractère), sociaux (rôles, statut) et culturels (valeurs, langue) qui constituent l'identité empirique, mesurable et déterminée d'un individu : c'est le « quoi », le masque (du latin \emph{persona}), l'enveloppe façonnée par l'hérédité et le milieu. La \textbf{personne}, en revanche, désigne le sujet moral et existentiel, porteur d'une valeur absolue — \textbf{dignité} chez \textbf{Kant} (\textit{Fondements de la métaphysique des mœurs}), \textbf{existence} et \textbf{communion} chez \textbf{Mounier} — qui ne se mesure pas, ne s'échange pas, et fonde les droits et la responsabilité : c'est le « qui », le noyau spirituel libre.

\textbf{[Problématisation]} Le sujet pose une tension paradoxale. D'un côté, le respect de la personne (Kant) semble exiger de ne pas juger sur les apparences (personnalité) : « Traite l'humanité... comme une fin », abstraction faite des qualités contingentes. L'effacement des différences de personnalité (richesse, intelligence, origine) serait la condition de l'égale dignité. De l'autre côté (Mounier), la personne ne s'incarne que \emph{dans} et \emph{par} sa personnalité : l'effacement de la personnalité (uniformisation, totalitarisme, aliénation) serait la \emph{négation} même de la personne qui a besoin de s'exprimer singulièrement. \textbf{Dès lors, le respect de la personne passe-t-il par l'abstraction des particularités de la personnalité (universalisme kantien) ou par la reconnaissance et la promotion de sa singularité concrète (personnalisme mouniérien) ?}

\textbf{[Annonce du plan]} Nous analyserons d'abord comment le respect kantien de la personne impose une mise entre parenthèses de la personnalité (I), avant de montrer que pour Mounier, au contraire, la personne ne se réalise qu'en \emph{personnalisant} sa personnalité (II), pour enfin dégager une articulation : le respect juridique (Kant) protège l'espace où s'exerce la liberté de personnalisation (Mounier) (III).
\end{quote}

\textbf{CONCLUSION}
\begin{quote}
\textbf{[Bilan]} Au terme de cette analyse, il apparaît que le respect de la personne n'exige pas l'effacement pur et simple de la personnalité, mais en \textbf{transforme le statut}. La première partie a montré que, pour \textbf{Kant}, le respect de la dignité impose une \textbf{abstraction égalitaire} : face à la loi morale et au droit, la personnalité (talents, origine, caractère) doit s'effacer pour ne laisser paraître que l'humanité commune, fin en soi. La deuxième partie a démontré que, pour \textbf{Mounier}, cet effacement serait une \textbf{mort de la personne} : celle-ci n'existe qu'en \textbf{existant}, c'est-à-dire en s'appropriant sa personnalité concrète (corps, histoire, culture) pour la porter à des valeurs universelles dans la communion. La troisième partie a synthétisé ces deux pôles : le \textbf{respect formel (Kant)} crée l'espace de sécurité (droits, non-instrumentalisation) sans lequel la \textbf{personnalisation concrète (Mounier)} est impossible.

\textbf{[Réponse]} \textbf{En définitive, le respect de la personne n'exige pas l'effacement de la personnalité, mais sa \emph{subordination} : la personnalité doit servir la personne, non l'asservir.}

\textbf{[Ouverture]} Cette articulation éclaire les débats bioéthiques actuels : le \textbf{consentement libre et éclairé} (héritage kantien : respect de la personne comme autonomie) ne suffit-il pas ? Ne faut-il pas ajouter, avec Mounier, une attention à la \textbf{vulnérabilité concrète} de la personnalité (handicap, vieillesse, détresse psychique) pour que le respect ne reste pas une formalité vide mais devienne une \emph{communion solidaire} ?
\end{quote}
\end{exoresolu}

\begin{methode}[Commentaire de texte philosophique : La méthode en 4 temps (Mounier/Kant)]
\textbf{Objectif :} Réussir l'exercice spécifique du commentaire de texte au Bac (savoir-faire 2).
\begin{enumerate}
    \item \textbf{Temps 1 : L'Introduction (Identification + Problème + Plan) :}
    \begin{itemize}
        \item \textbf{Référence :} Auteur, Œuvre, Date, Contexte (mouvement : Personnalisme, Critique...).
        \item \textbf{Thèse principale :} Une phrase : « Dans ce texte, X soutient que... »
        \item \textbf{Problématique du texte :} Quelle question l'auteur résout-il ? (Souvent : Comment penser la personne sans la réduire à la personnalité/à l'individu ?).
        \item \textbf{Plan du texte (suivre la structure logique de l'extrait) :} « Le texte s'articule en trois temps : d'abord... (lignes 1 à X), ensuite... (lignes X à Y), enfin... (lignes Y à la fin). »
    \end{itemize}
    \item \textbf{Temps 2 : L'Explication linéaire (Suivre le fil du texte) :}
    \begin{itemize}
        \item Parcourir le texte \textbf{paragraphe par paragraphe} (ou mouvement par mouvement).
        \item \textbf{Pour chaque étape :} Reformuler l'idée \textbf{en propres mots} + \textbf{Analyser les concepts clés} (définir, distinguer) + \textbf{Justifier l'argumentation} (pourquoi dit-il ça ? exemple, référence implicite, opposition) + \textbf{Montrer l'enchaînement} (connecteurs : « C'est pourquoi », « Cependant », « En ce sens »).
        \item \textit{Ne pas faire de hors-texte} (pas de dissertation déguisée), mais \textbf{éclairer} le texte par le cours (Kant/Mounier).
    \end{itemize}
    \item \textbf{Temps 3 : La Discussion Critique (Évaluation portée du texte) :}
    \begin{itemize}
        \item \textbf{Portée / Vérité :} Pourquoi cette thèse est-elle forte ? (Apport conceptuel, réponse à un problème réel, ex: Mounier contre individualisme/collectivisme).
        \item \textbf{Limites / Objections :} Qu'est-ce que le texte ne dit pas ? Quelles difficultés ? (Ex: Comment garantir la communion sans État de droit kantien ? Risque de subjectivisme ? Difficulté de définir le « noyau »).
        \item \textbf{Dépassement / Complément :} Autre auteur (Kant si texte Mounier, ou Ricœur, Levinas, Arendt) qui complète ou corrige.
    \end{itemize}
    \item \textbf{Temps 4 : La Conclusion (Bilan + Ouverture) :} Même structure BRO que dissertation.
\end{enumerate}
\cle{Linéaire $\neq$ Paraphrase. Il faut \textbf{expliciter} (rendre explicite l'implicite) et \textbf{articuler} (montrer la logique).}
\end{methode}

\begin{exoresolu}[Commentaire guidé : Extrait de Mounier « Qu'est-ce que le personnalisme ? »]
\textbf{Énoncé :} \textbf{Texte :} « La personne n'est pas l'individu. L'individu est un atome, une monade close sur elle-même, qui ne communique qu'accidentellement. La personne, au contraire, est un être de communion. Elle ne vit que par les autres, elle ne se possède qu'en se donnant. La personnalité est l'enveloppe, le masque, l'ensemble des caractères acquis ou innés. La personne est le noyau spirituel qui anime cette enveloppe, qui la dépasse infiniment. On a une personnalité, on est une personne. »
\textit{(Inspiré de Mounier, \textit{Qu'est-ce que le personnalisme ?}, 1947)}
\textbf{Consigne :} Proposez une introduction type, puis les grandes lignes de l'explication linéaire (3 mouvements), et les axes de discussion critique.
\tcblower
\textbf{Solution détaillée :}

\textbf{INTRODUCTION TYPE}
\begin{itemize}
    \item \textbf{Référence :} \textbf{Emmanuel Mounier} (1905-1950), figure majeure du \textbf{personnalisme communautaire} français, directeur de la revue \textit{Esprit}. Texte extrait de \textit{Qu'est-ce que le personnalisme ?} (1947), écrit dans l'après-guerre, contre le nazisme (collectivisme totalitaire) et le libéralisme (individualisme atomisé).
    \item \textbf{Thèse :} La \textbf{personne} se distingue radicalement de l'\textbf{individu} (être clos) et de la \textbf{personnalité} (masque déterminé) : elle est \textbf{être de communion} (relation constitutive) et \textbf{noyau spirituel libre} (existence).
    \item \textbf{Problématique :} Comment penser l'unité de l'être humain sans le réduire ni à un atome social (individu) ni à un déterminisme psychologique (personnalité) ?
    \item \textbf{Plan du texte (3 mouvements) :}
    \begin{enumerate}
        \item \textbf{Lignes 1-3 :} Distinction ontologique \textbf{Individu vs Personne} (Clos vs Communion).
        \item \textbf{Lignes 3-5 :} Distinction structurelle \textbf{Personnalité vs Personne} (Masque/Enveloppe vs Noyau/Esprit).
        \item \textbf{Ligne 5 :} Synthèse existentielle : \textbf{Avoir vs Être} (Possession vs Don).
    \end{enumerate}
\end{itemize}

\textbf{EXPLICATION LINÉAIRE (Grandes lignes)}

\textbf{Mouvement 1 : L'Individu, figure de la clôture (lignes 1-3).}
\begin{itemize}
    \item \textbf{Analyse :} Mounier définit l'individu par la \textbf{négation} : « atome », « monade close » (réf. Leibniz détournée : la monade n'a pas de fenêtres, mais Mounier dit qu'elle communique « accidentellement »).
    \item \textbf{Concepts :} \textbf{Individualisme} (thèse adverse : société = somme d'atomes). \textbf{Accidentel} = contingent, externe, non nécessaire.
    \item   \textbf{Argumentation :} L'individu est une \textbf{fiction abstraite} (produit de la modernité libérale). La relation aux autres est extérieure (contrat, échange, conflit), pas constitutive.
\end{itemize}

\textbf{Mouvement 2 : La Personne, figure de la communion (lignes 3-5 « La personne... par les autres »).}
\begin{itemize}
    \item \textbf{Rupture :} « Au contraire » (adversatif fort). \textbf{Communion} $\neq$ Communication (échange d'infos) $\neq$ Fusion (perte de soi).
    \item \textbf{Paradoxe :} « Ne se possède qu'en se donnant. » Loi de l'\textbf{amour/amitié/engagement} : la plénitude de l'être est dans le don (théologie chrétienne implicite : Trinité, Kenose).
    \item \textbf{Concept clé :} \textbf{Existence} = Être hors de soi (ek-sistere), tourné vers l'autre.
\end{itemize}

\textbf{Mouvement 3 : Personnalité (Masque) vs Personne (Noyau) - Avoir vs Être (fin du texte).}
\begin{itemize}
    \item \textbf{Métaphore :} \textbf{Enveloppe / Masque (Persona théâtrale)} = Personnalité. Donnée (inné/acquis), objective, mesurable, finie. « Caractères acquis ou innés » $\to$ Reprise des facteurs bio/socio/culturels du cours.
    \item \textbf{Réalité profonde :} \textbf{Noyau spirituel} = Personne. Sujet, liberté, infini (dépasse l'enveloppe), source d'unité.
    \item \textbf{Formule choc :} « On \textbf{a} une personnalité, on \textbf{est} une personne. » (Verbe \emph{avoir} = possession, objet ; Verbe \emph{être} = existence, sujet).
    \item \textbf{Conséquence :} La personnalité est \textbf{instrumentalisée} par la personne pour la communion. La personne \textbf{personnalise} la personnalité.
\end{itemize}

\textbf{DISCUSSION CRITIQUE (Axes)}

\textbf{1. Portée (Force du texte) :}
\begin{itemize}
    \item \textbf{Contre le totalitarisme :} L'individu est sacrifiable pour le Tout (État, Race, Classe). La Personne est \textbf{insacrifiable} (fin en soi) car elle \emph{est} relation (communion), pas partie d'une machine.
    \item \textbf{Contre le psychologisme/réductionnisme :} Libère l'homme de son « CV » (personnalité). Espérance : je ne suis pas mon passé, mes gènes, mes névroses. Je suis un \textbf{projet}.
    \item \textbf{Fondation de l'éthique :} Le don de soi (communion) fondement de la vie bonne.
\end{itemize}

\textbf{2. Limites / Objections :}
\begin{itemize}
    \item \textbf{Risque de spiritualisme idéaliste :} Où est le corps ? La maladie, la douleur, la mort « personnelles » ? Le « noyau » est-il observable ? Risque de \textbf{dualisme} (âme/corps) que Mounier veut éviter (« esprit incarné ») mais que le texte suggère (noyau/enveloppe).
    \item \textbf{La communion est-elle toujours possible ?} Que faire de l'ennemi, de l'étranger radical, du pervers narcissique ? \textbf{Kant (Droit/Respect)} est nécessaire là où la \textbf{Communion (Amour)} fait défaut. Le respect formel protège là où le don est impossible.
    \item \textbf{Subjectivisme :} Si la personne est « noyau » inaccessible, comment juger les actes ? Responsabilité ? Mounier répond : la personne \emph{s'engage}, elle est responsable de sa personnalisation. Mais le texte ne le dit pas explicitement.
\end{itemize}

\textbf{3. Dépassement (Ricœur / Levinas) :}
\begin{itemize}
    \item \textbf{Paul Ricœur (\textit{Soi-même comme un autre}) :} Réconcilie \emph{ipseité} (moi-même, noyau, fidélité à soi = Personne Mounier) et \emph{identité narrative} (personnalité, histoire, masque). La personne \emph{raconte} sa personnalité.
    \item \textbf{Emmanuel Levinas :} Le « Visage » de l'Autre (proche de « Personne ») m'oblige \emph{avant} toute communion choisie. Responsabilité infinie $\to$ Fondement éthique plus rigoureux que la communion (qui suppose réciprocité).
\end{itemize}

\textbf{CONCLUSION TYPE (BRO)}
\begin{quote}
\textbf{Bilan :} Mounier redéfinit l'humain : ni atome individuel, ni masque psychologique, la personne est \textbf{communion vivante} et \textbf{noyau de liberté}. \textbf{Réponse :} La personnalité est le véhicule mortel ; la personne est le voyageur immortel (par le don). \textbf{Ouverture :} Cette définition fonde-t-elle une politique ? Le personnalisme peut-il éviter le communautarisme (repli sur sa communion) sans le droit universel kantien ?
\end{quote}
\end{exoresolu}

\begin{attention}[Les 5 erreurs qui coûtent des points au Bac]
\begin{enumerate}
    \item \textbf{Confondre Personnalité et Personne (Erreur conceptuelle majeure).} \\
    \textit{Exemple :} « La dignité de la personnalité... », « Il a une belle personne (au sens moral) ». \\
    \textit{Reflexe :} \textbf{Personnalité = Fait (Psychologie/Sociologie). Personne = Valeur (Éthique/Ontologie/Juridique).} Utiliser \textbf{« Avoir » vs « Être »} (Mounier) ou \textbf{« Prix » vs « Dignité »} (Kant) à chaque paragraphe.
    
    \item \textbf{Réduire Kant à « Respectez les gens » ou Mounier à « Aimez-vous les uns les autres » (Réductionnisme moralisateur).} \\
    \textit{Erreur :} Oublier la \textbf{rigueur conceptuelle} : Autonomie, Noumène, Impératif Catégorique, Universalisation (Kant) ; Existence, Communion, Personnalisation, Spiritualité incarnée (Mounier). \\
    \textit{Reflexe :} Citer les \textbf{concepts techniques} exacts. Expliquer \emph{pourquoi} c'est une fin en soi (autonomie législative), pas juste « c'est bien ».
    
    \item \textbf{Traiter les facteurs de personnalité comme une liste de courses (Sans interaction).} \\
    \textit{Erreur :} « Il y a le biologique : ... Le social : ... Le culturel : ... » (Juxtaposition). \\
    \textit{Reflexe :} \textbf{Montrer la dialectique :} Le biologique est \emph{modelé} par le social ; le culturel donne du \emph{sens} au biologique. La personnalité est une \textbf{synthèse originale} (Le Moigne), pas une addition. Citer l'exemple d'Amadou ou Aïcha.
    
    \item \textbf{Introduction/Conclusion « hors-sujet » ou « bac à sable ».} \\
    \textit{Erreurs :} Accroche vague (« Depuis l'Antiquité... »), Définitions manquantes ou fausses, Problématique absente ou fermée (Oui/Non), Plan flou (« Nous verrons les pour et les contre »). Conclusion : « En conclusion, on a vu que... » (Répétition), Ouverture qui relance le sujet (« Peut-on vraiment savoir ? »). \\
    \textit{Reflexe :} Appliquer \textbf{ADPA / BRO} à la lettre. S'entraîner sur 5 sujets types avant l'examen.
    
    \item \textbf{Commentaire de texte = Dissertation déguisée ou Paraphrase.} \\
    \textit{Erreurs :} Ne pas suivre l'ordre du texte (sauter des lignes), Importer son cours sans l'accrocher aux \textbf{mots du texte} (guillemets obligatoires), Faire l'impasse sur la \textbf{Discussion critique} (1/3 de la note), Oublier la \textbf{Problématique du texte} (différente de celle du cours). \\
    \textit{Reflexe :} \textbf{Linéaire rigoureux} (citer l. 1-3, l. 4-6...), \textbf{Guillemets} pour chaque analyse, \textbf{Discussion} structurée (Portée / Limites / Autre auteur).
\end{enumerate}
\end{attention}

\newpage

\coursec{Exercices}

\courssub{Je vérifie mes connaissances}

\begin{exercice}[QCM : notions fondamentales]
Pour chaque question, une seule réponse est correcte. Recopiez le numéro et la lettre choisie.
\begin{enumerate}
    \item La personnalité, au sens philosophique, se distingue du tempérament par le fait qu'elle est :
    \begin{itemize}
        \item[A.] Innée et fixée génétiquement.
        \item[B.] Le résultat d'une construction par l'individu lui-même au cours de son existence.
        \item[C.] Identique au caractère.
        \item[D.] Une simple somme de réflexes conditionnés.
    \end{itemize}
    \item Pour Kant, dans le \emph{Fondement de la métaphysique des mœurs}, ce qui a un \textbf{prix} peut être remplacé par un équivalent ; ce qui a une \textbf{dignité} :
    \begin{itemize}
        \item[A.] A un prix très élevé.
        \item[B.] Ne peut être remplacé par aucun équivalent et est au-dessus de tout prix.
        \item[C.] Se négocie sur le marché du travail.
        \item[D.] Dépend de l'utilité sociale de la personne.
    \end{itemize}
    \item Le personnalisme de Mounier affirme que la personne ne se réalise pleinement que dans la \textbf{communion}. Cela signifie que :
    \begin{itemize}
        \item[A.] L'individu doit s'isoler pour protéger son intériorité.
        \item[B.] La personne s'accomplit dans l'ouverture à autrui et l'engagement commun.
        \item[C.] La valeur de la personne dépend de son rang social.
        \item[D.] La communion est une fusion qui efface les différences.
    \end{itemize}
    \item Parmi les facteurs d'édification de la personnalité, l'école, la famille et les médias relèvent du :
    \begin{itemize}
        \item[A.] Facteur biologique (hérédité).
        \item[B.] Facteur psychologique strict (tempérament).
        \item[C.] Facteur social et culturel (environnement).
        \item[D.] Facteur métaphysique (destin).
    \end{itemize}
\end{enumerate}
\end{exercice}

\begin{exercice}[Vrai ou Faux ? Justifiez brièvement]
Indiquez si l'affirmation est vraie (V) ou fausse (F) et justifiez votre réponse en une phrase en mobilisant le cours.
\begin{enumerate}
    \item « Les jumeaux monozygotes élevés dans le même milieu développent nécessairement la même personnalité. »
    \item « Traiter un être humain comme une simple main-d'œuvre dans une plantation de cacao, c'est lui attribuer un prix et nier sa dignité. »
    \item « Pour Mounier, la personne est une substance close sur elle-même, une \emph{monade} sans fenêtre. »
    \item « Le tempérament est acquis par l'éducation alors que le caractère est inné. »
\end{enumerate}
\end{exercice}

\begin{exercice}[Définitions attendues au Bac]
Donnez une définition précise et philosophique (2 à 3 lignes) pour chacun des termes suivants :
\begin{enumerate}
    \item \textbf{Personnalité} (distinction avec individualité).
    \item \textbf{Personne} (selon la définition juridique et morale).
    \item \textbf{Dignité} (selon Kant).
    \item \textbf{Personnalisme} (selon Mounier : notion de communion).
\end{enumerate}
\end{exercice}

\begin{exercice}[Association et complétion]
\begin{enumerate}
    \item Reliez chaque auteur à sa thèse centrale sur la valeur de la personne :
    \begin{center}
        \begin{tabular}{|c|c|}
        \hline
        \textbf{Auteur} & \textbf{Thèse} \\ \hline
        Kant & 1. La personne se réalise dans la communion et l'engagement. \\ \hline
        Mounier & 2. La personne a une dignité, fin en soi, au-dessus de tout prix. \\ \hline
        \end{tabular}
    \end{center}
    \item Complétez la citation de Kant : « Tout a soit un \underline{\hspace{3cm}}, soit une \underline{\hspace{3cm}}. »
\end{enumerate}
\end{exercice}

\courssub{Je m'entraîne}

\begin{exercice}[Analyse des facteurs biologiques : le cas des jumeaux]
Dans un quartier de Yaoundé, deux frères jumeaux monozygotes (vrais jumeaux), Paul et Pierre, ont grandi ensemble. Paul est devenu un artisan menuisier calme et méthodique ; Pierre est devenu un chauffeur de taxi-brousse extraverti et impatient.
\begin{enumerate}
    \item Identifiez le facteur biologique commun aux deux frères.
    \item Expliquez pourquoi, malgré ce facteur commun, leurs personnalités diffèrent. Citez au moins deux facteurs non biologiques en jeu.
    \item En vous appuyant sur cet exemple, formulez une thèse nuancée sur le rôle de l'hérédité dans l'édification de la personnalité.
\end{enumerate}
\end{exercice}

\begin{exercice}[Facteurs sociaux et culturels : l'école et le quartier]
Amina, 17 ans, élève en Première Technique à Garoua, provient d'une famille peule traditionnelle. Elle suit des cours de philosophie, pratique le codage informatique au club numérique de l'école et participe aux veillées traditionnelles le week-end.
\begin{enumerate}
    \item Classez les influences suivantes dans les catégories « Facteur social » ou « Facteur culturel » : la langue fulfulde, l'institution scolaire, les rites d'initiation, le groupe d'amis du club numérique.
    \item Montrez comment l'interaction de ces facteurs contribue à forge \emph{une} personnalité singulière (et non déterminée).
    \item Pourquoi dire qu'Amina « subit » son milieu serait une erreur philosophique au regard de la notion de personnalité ?
\end{enumerate}
\end{exercice}

\begin{exercice}[Kant : Prix vs Dignité dans une situation concrète]
Un patron d'une PME de Bafoussam propose à son employé, M. Kamga, une prime de 50 000 FCFA pour qu'il accepte de travailler un jour férié sans récupération, alors que le code du travail l'interdit. M. Kamga refuse en disant : « Ma dignité n'a pas de prix. »
\begin{enumerate}
    \item Expliquez la distinction kantienne entre \emph{prix} (valeur marchande) et \emph{dignité} (valeur intrinsèque).
    \item Montrez en quoi la proposition du patron traite M. Kamga comme un \emph{moyen} et non comme une \emph{fin en soi}.
    \item Le refus de M. Kamga illustre-t-il l'autonomie de la volonté ? Justifiez.
\end{enumerate}
\end{exercice}

\begin{exercice}[Mounier : Personnalisme et engagement associatif]
Un groupe d'élèves de Première Technique de Douala crée une association « Jeunes Solidaires » pour nettoyer les caniveaux de leur quartier avant la saison des pluies et aider les personnes âgées.
\begin{enumerate}
    \item Selon Mounier, en quoi cette action est-elle une \emph{communion} et non une simple \emph{collaboration} technique ?
    \item Expliquez la phrase de Mounier : « La personne est un être qui \emph{se donne} » à la lumière de cet engagement.
    \item En quoi le personnalisme s'oppose-t-il à l'individualisme égoïste et au collectivisme qui écrase l'individu ? Illustrer par l'exemple.
\end{enumerate}
\end{exercice}

\courssub{Je me prépare au Bac}

\begin{exercice}[Dissertation type Baccalauréat Technique]
\textbf{Sujet :} \og La personne humaine a-t-elle un prix ? \fg
\begin{enumerate}
    \item Analysez le sujet : définissez les termes, identifiez la tension (prix vs dignité), formulez la problématique.
    \item Proposez un plan détaillé en trois parties (Thèse / Antithèse / Synthèse) avec les arguments principaux, les références à Kant (obligatoire) et un exemple concret camerounais (travail des enfants, traite, corruption, don d'organes, etc.) pour chaque grande partie.
    \item Rédigez l'introduction complète (accroche, définition des termes, problématique, annonce du plan).
\end{enumerate}
\textit{Indication :} Vous devez mobiliser la distinction kantienne Prix/Dignité et le personnalisme de Mounier (communion) pour la synthèse.
\end{exercice}

\begin{exercice}[Explication de texte : Kant, \emph{Fondement de la métaphysique des mœurs}]
\textbf{Texte :}
\og Dans le royaume des fins, tout a soit un prix, soit une dignité. Ce qui a un prix peut être remplacé par quelque autre chose comme équivalent ; ce qui, au contraire, est au-dessus de tout prix, c'est-à-dire qui n'admet point d'équivalent, cela a une dignité. [...]
L'humanité même est une dignité ; car l'homme ne peut être utilisé par aucun homme (ni même par lui-même) comme simple moyen, mais doit en tout temps être traité comme une fin. C'est ici qu'est le principe suprême de la morale et de la législation universelle. \fg
\begin{enumerate}
    \item \textbf{Analyse du texte :}
    \begin{itemize}
        \item[a)] Dégagez la thèse principale du texte.
        \item[b)] Expliquez la distinction entre \emph{prix} et \emph{dignité} en vos mots.
        \item[c)] Que signifie la formule : « l'homme ne peut être utilisé [...] comme simple moyen, mais doit en tout temps être traité comme une fin » ?
    \end{itemize}
    \item \textbf{Discussion / Argumentation :}
    \begin{itemize}
        \item[a)] Pourquoi Kant affirme-t-il que « l'humanité même est une dignité » ? (Rôle de l'autonomie et de la raison).
        \item[b)] Donnez un exemple contemporain (Cameroun ou international) où des êtres humains sont traités comme ayant un « prix » (moyens) et non une dignité (fins). Analysez cet exemple à la lumière du texte.
        \item[c)] La reconnaissance de la dignité suffit-elle à garantir le respect des droits de l'homme dans la réalité sociale ? Argumentez.
    \end{itemize}
\end{enumerate}
\end{exercice}

\begin{exercice}[Étude de cas \& Question à argumentation construite : Travail des enfants et Hérédité]
\textbf{Document :} Dans une plantation de cacao du département de la Mbam-et-Inoubou (Centre-Cameroun), des enfants de 10 à 14 ans travaillent à la machette pour l'entretien des champs et la récolte, pendant les vacances scolaires et parfois en période de cours. Ils sont rémunérés 500 FCFA la journée. Les parents estiment que c'est une « éducation au travail » et un revenu nécessaire pour la famille. L'ONG locale « Enfance Protégée » dénonce une violation de la dignité de l'enfant et du droit à l'éducation.

\textbf{Questions :}
\begin{enumerate}
    \item \textbf{Analyse philosophique du cas (Mobilisation Kant / Mounier) :}
    \begin{itemize}
        \item[a)] En quoi la rémunération de 500 FCFA/jour institue-t-elle un \emph{prix} pour l'enfant ? Montrez que cela le réduit à l'état de \emph{moyen} (main-d'œuvre) au détriment de sa fin (éducation, développement).
        \item[b)] En vous référant à Mounier, expliquez pourquoi ce travail précoce et pénible empêche l'\emph{édification de la personnalité} (manque de communion scolaire, aliénation, absence de don de soi libre).
    \end{itemize}
    \item \textbf{Question à argumentation construite :} \og L'hérédité suffit-elle à expliquer la personnalité ? \fg
    \begin{itemize}
        \item[a)] Présentez la thèse de l'innéisme (hérédité = destin) et ses limites (exemple des jumeaux).
        \item[b)] Présentez la thèse de l'acquis (environnement, culture, éducation) et ses limites (déterminisme social).
        \item[c)] Synthétisez : La personnalité comme \emph{conquête} (intégration originale des donnés biologiques et sociaux par la liberté). Illustrer par le cas d'un enfant de la plantation qui reprend ses études et devient instituteur.
    \end{itemize}
\end{enumerate}
\end{exercice}

\newpage

\coursec{Corrigés des exercices}

\courssub{Je vérifie mes connaissances}

\begin{corrige}[Exercice 1 — QCM : notions fondamentales]
\begin{enumerate}
    \item \textbf{B.} La personnalité, au sens philosophique, n'est pas un donné biologique figé (tempérament), mais une \cle{construction} progressive. Elle résulte de l'intégration originale, par le sujet lui-même, de ses donnés innés, de son histoire, de ses choix et de ses engagements. C'est une \oe{}uvre de liberté.
    \item \textbf{B.} C'est la thèse centrale du \emph{Fondement de la métaphysique des mœurs}. Le \cle{prix} relève de l'échange marchand (valeur relative, équivalent possible). La \cle{dignité} est une \cle{valeur intrinsèque}, incomparable, qui place l'être humain \emph{au-dessus de tout prix} et interdit toute instrumentalisation.
    \item \textbf{B.} Pour Emmanuel \cle{Mounier}, la personne n'est pas une monade close. Elle ne s'accomplit que dans la \cle{communion} : une relation d'ouverture, de don de soi et d'engagement commun avec d'autres personnes. L'isolement (A) nie la dimension relationnelle ; le rang social (C) est extérieur ; la fusion (D) détruit la singularité personnelle.
    \item \textbf{C.} L'école, la famille, les médias, les pairs, les institutions sont des \cle{facteurs sociaux et culturels} (environnement). Ils fournissent les modèles, les normes, le langage, les valeurs qui structurent la personnalité. L'hérédité (A) est biologique ; le tempérament (B) est une base psycho-physiologique innée ; le destin (D) n'est pas une catégorie explicative philosophique standard.
\end{enumerate}
\end{corrige}

\begin{corrige}[Exercice 2 — Vrai ou Faux ? Justifiez brièvement]
\begin{enumerate}
    \item \textbf{Faux.} Les jumeaux monozygotes partagent le même patrimoine génétique (facteur biologique identique), mais la personnalité résulte de l'interaction entre l'inné et l'acquis. Des expériences de vie différentes (rencontres, choix, épreuves, interprétations subjectives) mènent à des personnalités distinctes. L'hérédité ne \emph{détermine} pas la personnalité.
    \item \textbf{Vrai.} Payer un travailleur pour une tâche interdite par le code du travail, sans respect du repos légal, c'est le considérer comme un \cle{moyen} (une main-d'œuvre interchangeable, avoir un \cle{prix}) au service du profit, niant sa qualité de \cle{fin en soi} (dignité, autonomie, droits).
    \item \textbf{Faux.} C'est l'exact contraire de la thèse de Mounier. La personne mouniérienne est \emph{ouverte}, \emph{relationnelle}, elle « a des fenêtres ». La \cle{communion} est sa loi d'être. La « monade sans fenêtre » est une image leibnizienne de la substance close, que le personnalisme rejette.
    \item \textbf{Faux.} C'est l'inverse : le \cle{tempérament} est la part \cle{innée}, constitutionnelle, stable (base biologique) ; le \cle{caractère} est la part \cle{acquise}, modelée par l'éducation, l'environnement, les habitudes et la volonté.
\end{enumerate}
\end{corrige}

\begin{corrige}[Exercice 3 — Définitions attendues au Probatoire / Baccalauréat technique]
\begin{enumerate}
    \item \textbf{Personnalité :} Ensemble dynamique et structuré de traits psychologiques, de valeurs et de modes de réaction qui singularisent un individu. Contrairement à l'\cle{individualité} (fait brut d'être un indivisible, un « un » numérique, identique pour tous), la personnalité est une \cle{qualité} morale et existentielle, une \oe{}uvre de construction personnelle (originalité, responsabilité, unité de vie).
    \item \textbf{Personne :} Sujet de droit et de devoirs, être doué de \cle{raison}, de \cle{liberté} et de \cle{responsabilité}. Juridiquement : titulaire de droits subjectifs (capacité juridique). Moralement : \cle{fin en soi} (Kant), être de dignité, capable d'autonomie et d'engagement (Mounier). S'oppose à la « chose » (objet manipulable) et au simple « individu » (unité biologique).
    \item \textbf{Dignité (Kant) :} \cle{Valeur intrinsèque} et \cle{incomparable} de l'être rationnel. Ce qui est \emph{au-dessus de tout prix} et n'admet \cle{aucun équivalent}. Fondée sur l'\cle{autonomie} de la volonté (capacité de se donner sa propre loi morale) et sur la qualité de \cle{fin en soi}. Exige le respect absolu : ne jamais traiter l'humanité (en soi ou autrui) comme simple moyen.
    \item \textbf{Personnalisme (Mounier — communion) :} Doctrine qui place la \cle{personne} (être spirituel, libre, responsable, ouvert) au centre de la réalité. La personne ne se réalise que dans la \cle{communion} : relation d'\cle{amour}, de \cle{don de soi} et d'\cle{engagement} solidaire avec d'autres personnes. S'oppose à l'individualisme (repli sur soi) et au collectivisme (écrasement de l'individu par le groupe).
\end{enumerate}
\end{corrige}

\begin{corrige}[Exercice 4 — Association et complétion]
\begin{enumerate}
    \item \textbf{Correspondances :}
    \begin{itemize}
        \item Kant $\rightarrow$ \textbf{2.} La personne a une dignité, fin en soi, au-dessus de tout prix.
        \item Mounier $\rightarrow$ \textbf{1.} La personne se réalise dans la communion et l'engagement.
    \end{itemize}
    \item \textbf{Citation complétée :} « Tout a soit un \textbf{prix}, soit une \textbf{dignité}. » (\emph{Fondement de la métaphysique des mœurs}, II).
\end{enumerate}
\end{corrige}

\courssub{Je m'entraîne}

\begin{corrige}[Exercice 5 — Analyse des facteurs biologiques : le cas des jumeaux]
\begin{enumerate}
    \item \textbf{Facteur biologique commun :} Le \cle{patrimoine génétique} (ADN quasi identique, jumeaux monozygotes). Cela détermine le \cle{tempérament} de base (ex. : seuil de réactivité, morphologie, prédispositions physiologiques).
    \item \textbf{Facteurs non biologiques différenciants (au moins deux) :}
    \begin{itemize}
        \item \cle{Facteur social / professionnel :} Métiers différents (menuisier vs chauffeur taxi-brousse). L'apprentissage, la discipline du geste, la relation client, la gestion du temps, les risques quotidiens forgent des habitudes, des vertus (patience, prudence) ou des vices (impatience, prise de risque) distincts.
        \item \cle{Facteur culturel / relationnel :} Milieux de travail distincts, pairs différents, valeurs valorisées (précision/artisanat vs rapidité/rentabilité), modèles identitaires.
        \item \cle{Facteur psychologique / existentiel :} Choix personnels, interprétation subjective des événements, projets de vie, rencontres amoureuses ou amicales singulières.
    \end{itemize}
    \item \textbf{Thèse nuancée :} L'hérédité fournit le \cle{matériel de base} (tempérament, potentialités, limites biologiques), elle est une \cle{condition nécessaire} mais \cle{insuffisante}. La personnalité est l'\cle{intégration originale} de ce donné par le sujet dans une histoire singulière. Les jumeaux montrent qu'à \emph{donné biologique égal}, des \emph{existences différentes} produisent des \emph{personnalités différentes}. L'hérédité détermine le \emph{champ des possibles}, la liberté choisit et construit dans ce champ.
\end{enumerate}
\end{corrige}

\begin{corrige}[Exercice 6 — Facteurs sociaux et culturels : l'école et le quartier]
\begin{enumerate}
    \item \textbf{Classification :}
    \begin{center}
        \begin{tabularx}{\linewidth}{|l|X|X|}\hline
        \textbf{Influence} & \textbf{Facteur Social} & \textbf{Facteur Culturel} \\ \hline
        Langue fulfulde & & X (langue, tradition orale) \\ \hline
        Institution scolaire & X (structure, normes, diplôme) & X (savoirs, valeurs citoyennes, laïcité) \\ \hline
        Rites d'initiation & X (groupe d'âge, statut) & X (symboles, transmission valeurs traditionnelles) \\ \hline
        Groupe d'amis club numérique & X (réseau pair, appartenance) & X (culture tech, langage code, valeurs innovation) \\ \hline
        \end{tabularx}
    \end{center}
    \textit{Note :} La distinction est poreuse ; l'école est à la fois institution sociale et vecteur culturel.
    \item \textbf{Personnalité singulière :} Amina ne subit pas passivement ces influences. Elle les \cle{réapproprie} : elle \emph{traduit} la tradition peule dans le langage du code, elle \emph{critique} ou \emph{valorise} les rites à la lumière de la philosophie, elle \emph{choisit} ses engagements au club. Sa personnalité émerge de la \cle{synthèse originale} qu'elle opère entre ces mondes (ex. : créer une application en fulfulde pour préserver les contes). C'est une \cle{conquête}, non un déterminisme.
    \item \textbf{Erreur du « subir » :} Dire qu'elle « subit » son milieu revient à nier sa \cle{liberté} et sa \cle{responsabilité}, donc sa qualité de \cle{personne}. La personnalité implique une \cle{prise de distance} critique vis-à-vis des déterminismes (biologiques, sociaux, culturels). Amina \emph{fait} sa personnalité en répondant à son milieu ; elle en est l'\cle{auteur} (même si contrainte). Le déterminisme social nie la subjectivité agente.
\end{enumerate}
\end{corrige}

\begin{corrige}[Exercice 7 — Kant : Prix vs Dignité dans une situation concrète]
\begin{enumerate}
    \item \textbf{Distinction :} Le \cle{prix} (valeur marchande) s'attache aux \cle{moyens} : objets, travail, services. Ils sont \cle{fongibles} (remplaçables par un équivalent : 50 000 FCFA = une journée de travail). La \cle{dignité} s'attache à la \cle{fin en soi} : l'être humain comme sujet moral, libre et rationnel. Elle est \cle{infinie}, \cle{incomparable}, \cle{non fongible}. On ne peut « acheter » le consentement moral d'une personne ; on ne peut compenser la violation de ses droits par de l'argent.
    \item \textbf{Instrumentalisation :} Le patron propose un \cle{échange marchand} : argent contre violation du droit (repos légal). Il considère M. Kamga comme une \cle{ressource} (moyen) pour produire du profit, dont la volonté peut être achetée. Il nie son \cle{autonomie} (capacité de se donner sa propre loi, ici le respect du code du travail) et sa \cle{dignité} (valeur qui interdit de marchander ses droits fondamentaux). M. Kamga est réduit à une \cle{chose} à prix.
    \item \textbf{Autonomie de la volonté :} \textbf{Oui.} Refuser l'argent pour respecter la loi et sa propre valeur, c'est agir par \cle{devoir} (respect de la loi morale) et non par \cle{inclination} (besoin d'argent). C'est se donner à soi-même la loi (autonomie : \emph{auto-nomos}) : « Je ne me vends pas ». C'est affirmer sa qualité de \cle{législateur} dans le \cle{royaume des fins}, non de sujet passif aux offres marchandes.
\end{enumerate}
\end{corrige}

\begin{corrige}[Exercice 8 — Mounier : Personnalisme et engagement associatif]
\begin{enumerate}
    \item \textbf{Communion vs Collaboration :} La \cle{collaboration} technique vise une \cle{efficacité} externe (caniveaux propres) ; les personnes sont des \cle{rouages} interchangeables. La \cle{communion} (Mounier) est une \cle{relation d'esprit à esprit} : les élèves se \cle{reconnaissent} comme personnes, partagent une \cle{intention commune} (solidarité, soin du quartier, respect des aînés), s'engagent \cle{librement} et \cle{gratuitement}. L'acte transforme les agents (ils grandissent en humanité) autant que le monde. C'est un \cle{don de soi}.
    \item \textbf{« La personne est un être qui se donne » :} La personne ne se possède pas comme un objet ; elle \cle{s'accomplit} en \cle{se dépensant} pour autrui. En nettoyant pour les autres, sans salaire, les élèves \emph{sortent d'eux-mêmes} (éc-stase), réalisent leur \cle{vocation personnelle} (être-pour-autrui). Le don n'est pas une perte, c'est la \cle{manière d'être} de la personne : « Je suis ce que je donne. »
    \item \textbf{Opposition aux deux écueils :}
    \begin{itemize}
        \item \cle{Individualisme égoïste :} « Chacun pour soi », je ne nettoie pas devant chez moi, je méprise le bien commun. L'association brise cet isolement par l'\cle{ouverture}.
        \item \cle{Collectivisme :} Le groupe (parti, État, secte) impose l'action, l'individu s'efface, devient un numéro. Ici, l'engagement est \cle{libre}, \cle{personnel}, \cle{créateur} (ils ont \emph{créé} l'asso). La communion \emph{exalte} la singularité de chacun (compétences, idées) au lieu de l'écraser.
    \end{itemize}
    \textit{Illustration :} Un élève timide prend la parole pour organiser ; un fort physiquement porte les sacs ; un autre gère la communication. Chacun \emph{se donne} selon sa singularité pour le \emph{bien commun}.
\end{enumerate}
\end{corrige}

\courssub{Je me prépare au Bac}

\begin{corrige}[Exercice 9 — Dissertation type Baccalauréat Technique : « La personne humaine a-t-elle un prix ? »]
\textbf{Barème indicatif (sur 20) :} Analyse/Problématique (4) + Plan détaillé (6) + Introduction rédigée (5) + Culture philosophique / Exemples CM / Rigueur (5).

\begin{enumerate}
    \item \textbf{Analyse du sujet :}
    \begin{itemize}
        \item \textbf{Termes :} « Personne humaine » = être de raison, liberté, dignité (Kant), communion (Mounier). « Prix » = valeur marchande, équivalent monétaire, fongibilité, instrumentalisation. « Avoir un prix » = être réductible à une valeur d'échange, être traitée comme marchandise/moyen.
        \item \textbf{Tension :} Réalité sociale (travail salarié, traite, corruption, don d'organes rémunéré, GPA) où l'humain \emph{semble} avoir un prix VS Exigence morale/philosophique (Kant, Droits de l'Homme) : l'humain a une \cle{dignité}, est \cle{fin en soi}, \emph{au-dessus de tout prix}.
        \item \textbf{Problématique :} Si l'économie et la technique tendent à tout marchander (donnant un prix à l'humain), la philosophie peut-elle fonder une valeur absolue (dignité) qui résiste à toute évaluation marchande ? La dignité est-elle un leurre ou une conquête à défendre ?
    \end{itemize}

    \item \textbf{Plan détaillé (Thèse / Antithèse / Synthèse) :}

    \textbf{THÈSE : La personne humaine \emph{a} un prix (réalité du marché et de l'instrumentalisation).}
    \begin{itemize}
        \item \textbf{Arg 1 :} Le travail salarié : la force de travail est une marchandise (Marx). Exemple CM : \emph{Travail des enfants dans les plantations de cacao (Mbam-et-Inoubou)} : 500 FCFA/jour = prix explicite de l'enfant-moyen.
        \item \textbf{Arg 2 :} La corruption / népotisme : la fonction, la justice, le vote ont un prix. Exemple CM : \emph{Pot-de-vin pour un marché public, achat de consciences électorales} : la dignité du citoyen/fonctionnaire est monnayée.
        \item \textbf{Arg 3 :} La traite / esclavagisme moderne / GPA marchande : le corps, la procréation deviennent objets de contrat commercial. Exemple international : \emph{Mères porteuses en Ukraine/Inde, traite d'êtres humains en Libye}.
    \end{itemize}

    \textbf{ANTITHÈSE : La personne humaine n'a \emph{prix}, elle a une \emph{dignité} (fondement moral kantien).}
    \begin{itemize}
        \item \textbf{Arg 1 :} Distinction kantienne Prix/Dignité (\emph{Fondement}). La personne est \cle{fin en soi}, \cle{incomparable}. Aucun équivalent (argent, pouvoir, utilité) ne peut la remplacer.
        \item \textbf{Arg 2 :} Fondement : \cle{Autonomie} (législateur moral) et \cle{Raison}. Le prix suppose un \emph{intérêt} (inclination) ; la dignité suppose le \emph{respect} (devoir).
        \item \textbf{Arg 3 :} Conséquence juridique : Droits de l'Homme inaliénables (DUDH 1948, Constitution CM). Exemple CM : \emph{Refus du travail forcé, grève pour dignité (ex. personnel soignant), procès pour traite d'enfants gagnés par ONG}.
    \end{itemize}

    \textbf{SYNTHÈSE : La dignité est une \emph{exigence normative} qui doit \emph{réguler} la réalité marchande (Personnalisme Mounier).}
    \begin{itemize}
        \item \textbf{Arg 1 :} Le prix est une \cle{réalité factuelle} (nécessité économique, pouvoir) ; la dignité est une \cle{valeur régulative} (idéal moral, droit). Il faut \cle{socialiser} l'économie (salaire décent, protection sociale) pour que le prix du travail ne nie pas la dignité.
        \item \textbf{Arg 2 :} \cle{Mounier} : La personne se réalise dans la \cle{communion} (don gratuit, solidarité), non dans l'échange marchand. Exemple CM : \emph{Associations « Jeunes Solidaires », mutuelles de santé villageoises (tontines solidaires), don du sang bénévole} : actes de communion qui résistent à la marchandisation.
        \item \textbf{Arg 3 :} Vigilance : La technique (IA, biotech, data) crée de nouveaux marchés de l'humain (données, attention, génome). La dignité exige un \cle{cadre éthique et juridique} strict (RGPD, lois bioéthique, interdiction brevetage gène humain).
    \end{itemize}

    \item \textbf{Introduction rédigée :}
    \begin{quote}
    « L'homme n'a pas de prix, il a une dignité. » Cette assertion de Kant, dans le \emph{Fondement de la métaphysique des mœurs}, sonne comme un défi lancé à la réalité la plus triviale. Sur les marchés de Douala ou de Yaoundé, comme sur les places boursières mondiales, tout semble s'évaluer en francs CFA ou en dollars : la force de travail, le silence d'un témoin, la procréation via une mère porteuse, les données personnelles vendues aux géants du numérique. Le sujet « La personne humaine a-t-elle un prix ? » met aux prises deux ordres hétérogènes : l'ordre de l'\cle{échange}, où tout est fongible, mesurable, remplaçable, et l'ordre de la \cle{valeur absolue}, où l'être humain, en tant que sujet libre et raisonnable, refuse toute équivalence. La tension est vive : la personne \emph{subit} souvent d'avoir un prix (exploitation, traite, corruption), mais la philosophie et le droit affirment qu'elle \emph{ne peut en avoir} sans cesser d'être personne. Nous analyserons d'abord comment la réalité sociale et économique institue un prix de l'humain (Thèse), puis pourquoi la philosophie kantienne pose la dignité comme une valeur infranchissable (Antithèse), pour montrer enfin que la dignité n'est pas un donné statique mais une \cle{exigence de communion} (Mounier) qui doit sans cesse reconquérir le terrain de l'économie (Synthèse).
    \end{quote}
\end{enumerate}
\end{corrige}

\begin{corrige}[Exercice 10 — Explication de texte : Kant, \emph{Fondement de la métaphysique des mœurs}]

\begin{enumerate}
    \item \textbf{Analyse du texte :}
    \begin{itemize}
        \item[a)] \textbf{Thèse principale :} L'être humain (l'humanité en sa personne) possède une \cle{valeur absolue}, la \cle{dignité}, qui le distingue radicalement de toute \cle{chose} ayant un \cle{prix}. Cette dignité fonde le \cle{principe suprême de la morale} : l'interdiction de traiter l'homme comme simple moyen et l'obligation de le traiter toujours comme une fin en soi.
        \item[b)] \textbf{Distinction Prix / Dignité :}
        \begin{itemize}
            \item \cle{Prix} : Valeur \emph{relative}, \emph{marchande}, \emph{conditionnelle}. L'objet a un prix s'il peut être \emph{remplacé} par un équivalent (autre objet, somme d'argent). C'est la valeur des \cle{moyens} (utilité, rareté, désir).
            \item \cle{Dignité} : Valeur \emph{absolue}, \emph{intrinsèque}, \emph{inconditionnelle}. L'être qui a une dignité est \emph{au-dessus de tout prix}, \emph{irremplaçable} (n'admet point d'équivalent). C'est la valeur de la \cle{fin en soi}.
        \end{itemize}
        \item[c)] \textbf{Formule « fin en soi » :} Cela signifie que l'être rationnel (l'humanité) est la \cle{fin ultime} de toute action morale. On ne doit jamais l'utiliser \emph{uniquement} comme un instrument pour atteindre un autre but (simple moyen), sans respecter sa propre \cle{volonté}, sa \cle{liberté} et ses \cle{fins propres}. Tout maximé d'action doit pouvoir être universalisée comme loi pour un « royaume des fins » où chaque personne est législatrice et sujet.
    \end{itemize}

    \item \textbf{Discussion / Argumentation :}
    \begin{itemize}
        \item[a)] \textbf{Pourquoi l'humanité est une dignité ?} Parce que l'homme est un \cle{être rationnel} doué d'\cle{autonomie}. Il n'est pas poussé par des causes extérieures (instincts, désirs, contrainte physique) mais se \cle{donne sa propre loi} (loi morale). Cette \cle{liberté} (autonomie de la volonté) est la source de sa valeur infinie. Les choses (animaux, objets) n'ont que \emph{valeur relative} (prix) car elles sont \emph{hétérodéterminées} (moyens). L'homme, comme \cle{législateur universel} par sa raison, est la \cle{condition de possibilité} de toute valeur : sans sujet moral, rien n'a de prix.
        \item[b)] \textbf{Exemple contemporain :} La \emph{Gestation Pour Autrui (GPA) marchande} (ex. : cliniques en Ukraine, Géorgie, USA) ou \emph{Travail forcé / Traite d'enfants au Cameroun (plantations cacao, mines, travail domestique).}
        \begin{itemize}
            \item \textbf{Analyse :} Dans la GPA marchande, le corps de la femme, sa fonction procréative, l'enfant à naître deviennent des \cle{objets de contrat} : une somme d'argent (prix) contre un service de gestation (moyen). La femme est traitée comme un \cle{utérus à louer} (simple moyen pour le désir d'enfant du couple commanditaire), l'enfant comme un \cle{produit commandé}. Cela viole le texte : on utilise l'humanité (femme, enfant) comme \cle{moyen} (marchandise), niant leur \cle{dignité} (fin en soi). Au Cameroun, l'enfant dans la plantation est un \cle{moyen de production} (main-d'œuvre bon marché) au service du profit du planteur, au mépris de sa \cle{fin propre} (éducation, développement, jeu).
        \end{itemize}
        \item[c)] \textbf{La reconnaissance suffit-elle ? \textbf{Non.}} La reconnaissance \emph{juridique} (DUDH, Constitution) est une condition \cle{nécessaire} (cadre normatif, sanction), mais \cle{insuffisante} pour garantir le respect \emph{effectif} dans la réalité sociale.
        \begin{itemize}
            \item \textbf{Arguments :} 1) \emph{Faiblesse de l'État de droit :} Impunité des puissants, corruption de la justice (ex. : traite non jugée). 2) \emph{Pression économique :} La misère force à « vendre » sa dignité (travail dangereux, prostitution, vente d'organes, GPA) — la contrainte structurelle nie la liberté formelle. 3) \emph{Idéologie marchande / Technicienne :} L'utilitarisme, le transhumanisme, le capitalisme de surveillance réduisent l'humain à \cle{données}, \cle{capital humain}, \cle{ressource}. 4) \emph{Manque d'éducation :} Sans intériorisation de la dignité (éducation morale, philosophie), la loi reste lettre morte.
            \item \textbf{Nuance :} La reconnaissance juridique permet la \cle{résistance} (ONG, syndicats, lanceurs d'alerte, jurisprudence). C'est un \cle{levier} pour la conquête effective de la dignité (Mounier : la dignité se \emph{conquiert} dans la communion et l'engagement politique).
        \end{itemize}
    \end{itemize}
\end{enumerate}
\end{corrige}

\begin{corrige}[Exercice 11 — Étude de cas \& Question à argumentation construite : Travail des enfants et Hérédité]

\begin{enumerate}
    \item \textbf{Analyse philosophique du cas (Kant / Mounier) :}
    \begin{itemize}
        \item[a)] \textbf{Prix et instrumentalisation (Kant) :} La rémunération de \textbf{500 FCFA/jour} fixe explicitement un \cle{prix} à l'activité de l'enfant. L'enfant devient une \cle{marchandise} (force de travail) interchangeable. Le patron (ou le parent qui « loue » l'enfant) l'utilise comme un \cle{moyen} pour \emph{son} profit (entretien plantation, revenu familial). La \cle{fin propre} de l'enfant (son \cle{développement}, son \cle{droit à l'éducation}, son \cle{jeu}, sa \cle{construction personnelle}) est \emph{sacrifiée} au profit d'une fin étrangère. L'enfant est traité comme une \cle{chose} (ayant un prix) et non comme une \cle{personne} (ayant une dignité). Le consentement parental ne change pas la nature de l'acte : on ne peut disposer de la dignité d'autrui (ni même de la sienne, pour Kant).
        \item[b)] \textbf{Empêchement de l'édification de la personnalité (Mounier) :}
        \begin{itemize}
            \item \cle{Manque de communion scolaire :} L'école est le lieu privilégié de la \cle{socialisation secondaire}, de l'accès à la \cle{culture universelle}, de la rencontre avec des pairs et des maîtres qui permettent l'\cle{ouverture d'esprit} et la \cle{réflexivité}. L'enfant en est exclu.
            \item \cle{Aliénation :} Le travail pénible (machette, produits chimiques, charges lourdes) à 10 ans \emph{brise} le corps et l'esprit. L'enfant est \emph{aliéné} : son activité ne lui appartient pas, il ne s'y \cle{reconnaît} pas, il y subit la contrainte. Il ne \cle{se donne} pas librement ; il est \cle{pris}.
            \item \cle{Absence de don de soi libre :} La personnalité se forge dans le \cle{don gratuit} (engagement associatif, aide familiale choisie, étude). Ici, c'est de la \cle{contrainte} (survie, autorité parentale). L'enfant n'apprend pas la \cle{responsabilité} (répondre de ses actes librement) mais la \cle{soumission}.
        \end{itemize}
    \end{itemize}

    \item \textbf{Question à argumentation construite : « L'hérédité suffit-elle à expliquer la personnalité ? »}
    \begin{itemize}
        \item[a)] \textbf{Thèse de l'innéisme (Hérédité = Destin) :} La personnalité serait programmée dans les gènes (tempérament, intelligence, prédispositions psychiatriques, traits de caractère). \textbf{Exemple jumeaux :} Si l'hérédité suffisait, Paul et Pierre (Ex. 5) seraient identiques. \textbf{Limites :} Les jumeaux monozygotes élevés ensemble diffèrent ; élevés séparément, les différences s'amplifient. L'hérédité donne des \cle{potentialités} (plage de réaction), pas un \cle{devenir} figé. L'épigénétique montre que l'environnement active/inhibe les gènes.
        \item[b)] \textbf{Thèse de l'acquis (Environnement = Déterminisme social) :} La personnalité serait le produit pur de la société (famille, école, classe, culture, médias). « L'homme est fait par les situations » (Sartre inversé / Behaviorisme / Durkheim). \textbf{Limites :} Deux enfants dans la même famille/classe (ex. Amina et son frère hypothétique) deviennent des adultes différents. Le déterminisme social nie la \cle{liberté}, la \cle{créativité}, la \cle{résistance} (ex. résilience). Il ne rend pas compte de la \cle{singularité} irréductible.
        \item[c)] \textbf{Synthèse : La personnalité comme \cle{conquête} (Intégration originale par la liberté).} La personnalité n'est ni \emph{donnée} (innée), ni \emph{subie} (acquise), elle est \cle{faite}. Le sujet \cle{reçoit} des donnés (biologiques : tempérament, santé ; sociaux : milieu, langue, histoire) mais il les \cle{réapproprie}, les \cle{hiérarchise}, les \cle{transcende} par ses \cle{choix} et ses \cle{engagements}.
        
        \textbf{Illustration concrète : L'enfant de la plantation devenu instituteur}
        \begin{itemize}
            \item \textbf{Donnés initiaux :} \textbf{Koffi}, 12 ans, travaille dans la plantation de Mbam-et-Inoubou (donné biologique : force physique, santé fragile ; donné social : pauvreté, analphabétisme parental, norme travail enfant).
            \item \textbf{1. Rencontre / Événement déclencheur :} Un instituteur itinérant / ONG « Enfance Protégée » lui apprend à lire sous un manguier.
            \item \textbf{2. Choix libre (Liberté) :} Il \emph{décide} de fuir le travail des champs pour l'école, malgré la pression familiale (rupture).
            \item \textbf{3. Intégration originale :} Il mobilise sa \cle{mémoire orale} (culture peule/béti), sa \cle{connaissance du terrain} (biologie pratique) pour exceller en SVT/Géographie.
            \item \textbf{4. Engagement (Communion Mounier) :} Devenu instituteur, il \emph{se donne} en créant une école de brousse pour les enfants planteurs.
            \item \textbf{Conclusion :} Sa personnalité (instituteur engagé, passeur de culture) n'est \emph{ni} dans ses gènes, \emph{ni} dans la plantation. Elle est l'\oe{}uvre de sa \cle{liberté} intégrant ses blessures en \cle{force}.
        \end{itemize}
    \end{itemize}
\end{enumerate}
\end{corrige}

\newpage

\coursec{Fiche bilan}

\begin{popfigure}
\centering
\begin{tikzpicture}[
  mindmap,
  grow cyclic,
  every node/.style=concept,
  concept color=popblue,
  level 1/.append style={level distance=4.5cm, sibling angle=360/7},
  text=white,
  font=\small\bfseries,
  popblue/.style={concept color=popblue},
  poporange/.style={concept color=poporange},
  popgreen/.style={concept color=popgreen},
  poppurple/.style={concept color=poppurple},
  poppink/.style={concept color=poppink},
  popgold/.style={concept color=popgold},
  popdark/.style={concept color=popdark}
]
\node[popblue, align=center]{Personnalité\\et Personne}
  child[poporange] { node[align=center]{1. Définitions\\clés} }
  child[popgreen] { node[align=center]{2. Facteurs\\biologiques} }
  child[poppurple] { node[align=center]{3. Facteurs\\socio-culturels} }
  child[poppink] { node[align=center]{4. Kant :\\Dignité et Devoir} }
  child[popgold] { node[align=center]{5. Mounier :\\Personnalisme\\communautaire} }
  child[popdark] { node[align=center]{6. Distinction\\Personnalité /\\Personne} }
  child[popblue!70!popgreen] { node[align=center]{7. Méthode :\\application aux\\cas concrets} };
\end{tikzpicture}
\legende{Carte mentale de synthèse : Personnalité et Personne (1\textsuperscript{re} Tech). Au centre, la notion de l'unité ; autour, les sept points du cours : définitions, facteurs biologiques, facteurs socio-culturels, Kant, Mounier, distinction personnalité/personne, méthode.}
\end{popfigure}

\begin{aretenir}[Définitions essentielles]
\begin{itemize}
  \item \cle{Personnalité} : ensemble des traits psychologiques, affectifs, intellectuels et moraux qui singularisent un individu et le rendent reconnaissable ; c'est une structure dynamique, en devenir, qui se construit tout au long de la vie.
  \item \cle{Personne} : sujet de droits et de devoirs, être moral doté de dignité, fin en soi (Kant) ou être de communion et d'engagement (Mounier). La personne relève du plan juridique et moral, la personnalité du plan psychologique.
  \item \cle{Tempérament} : fond biologique, inné et relativement stable de la personnalité (humeur, réactivité, rythme de vie), lié à l'hérédité et à la physiologie.
  \item \cle{Caractère} : ensemble des tendances stables de conduite acquises par l'éducation, l'habitude et l'effort de la volonté ; il est le côté acquis et perfectible de la personnalité.
  \item \cle{Dignité} (Kant) : valeur intrinsèque et incomparable de l'être raisonnable ; ce qui a une dignité n'a pas de prix, ne peut être ni remplacé ni échangé contre un équivalent.
  \item \cle{Personnalisme} (Mounier) : philosophie qui affirme la primauté de la personne sur l'individu, la société et l'État ; la personne ne s'accomplit que dans la \emph{communion} avec autrui et l'\emph{engagement} dans l'action.
\end{itemize}
\end{aretenir}

\begin{aretenir}[Repères auteurs et thèses (à connaître)]
\begin{center}
\begin{tabularx}{\linewidth}{|>{\bfseries}l|X|X|}\hline
\rowcolor{popblueL}\textbf{Auteur} & \textbf{Thèse centrale} & \textbf{Concepts clés} \\\hline
Kant & L'être humain est une fin en soi ; la valeur morale réside dans la \cle{bonne volonté} qui agit par \cle{devoir}, selon l'impératif catégorique. & Dignité, autonomie, impératif catégorique, fin en soi, royaume des fins. \\\hline
Mounier & La personne n'est pas une substance close sur elle-même ; elle \cle{se fait} par l'ouverture aux autres (\cle{communion}) et par l'action dans l'histoire (\cle{engagement}). & Personnalisme communautaire, communion, engagement, personne contre individu. \\\hline
\end{tabularx}
\end{center}
\end{aretenir}

\begin{methode}[Analyser un cas concret : les facteurs de la personnalité]
\begin{enumerate}
  \item \textbf{Identifier} les traits de la personne décrite dans la situation (timidité, colère, sens du devoir, goût du travail en groupe, etc.).
  \item \textbf{Classer} chaque trait : relève-t-il de l'inné (tempérament, hérédité, physiologie) ou de l'acquis (caractère, éducation familiale, école, culture, groupe des pairs, médias) ?
  \item \textbf{Montrer l'interaction} des deux séries de facteurs : le biologique fournit un fond que le social oriente, corrige ou développe. Exemple : une timidité innée peut être surmontée grâce à la pratique du théâtre scolaire ou à la prise de parole en classe.
  \item \textbf{Conclure} : la personnalité n'est ni purement innée ni purement acquise ; elle est le produit d'une construction où le biologique et le socio-culturel s'articulent.
\end{enumerate}
\end{methode}

\begin{methode}[Dissertation : « La personne a-t-elle une valeur absolue ? »]
\begin{enumerate}
  \item \textbf{Introduction} : accroche (par exemple l'affirmation de la dignité humaine dans les grandes déclarations des droits de l'homme) ; définition des termes (\cle{personne}, \cle{dignité}, \cle{valeur absolue}) ; problématique (la valeur de la personne est-elle inconditionnelle ou dépend-elle des circonstances, de l'utilité, de la reconnaissance sociale ?) ; annonce du plan dialectique.
  \item \textbf{Thèse (Kant)} : la personne a une valeur absolue, inconditionnelle, parce qu'elle est un être raisonnable, fin en soi, doué d'autonomie ; elle a une dignité et non un prix.
  \item \textbf{Antithèse (critiques et Mounier)} : la valeur de la personne est relationnelle et historique ; elle est menacée par la réification (traiter l'homme comme une chose) et suppose la reconnaissance effective des autres ; sans communion ni engagement, la dignité reste abstraite.
  \item \textbf{Synthèse} : la dignité est une exigence normative universelle (Kant) qui ne se réalise concrètement que dans la vie commune et l'action partagée (Mounier).
  \item \textbf{Conclusion} : réponse nuancée à la question posée, avec ouverture (par exemple : quelles institutions garantissent aujourd'hui cette dignité au Cameroun ?).
\end{enumerate}
\end{methode}

\begin{methode}[Commentaire de texte philosophique]
\begin{enumerate}
  \item \textbf{Présenter} le texte : auteur, œuvre, thème, thèse défendue.
  \item \textbf{Expliquer} pas à pas, en suivant le mouvement du texte : dégager les concepts, les distinctions (par exemple dignité/prix, personne/individu) et les arguments.
  \item \textbf{Situer} le texte dans l'œuvre de l'auteur et dans son courant (déontologie kantienne, personnalisme mouniérien).
  \item \textbf{Discuter} : portée et limites de la thèse, objections possibles, actualité.
  \item \textbf{Conclure} par une évaluation personnelle et justifiée.
\end{enumerate}
\end{methode}

\begin{attention}[Erreurs fréquentes à éviter]
\begin{itemize}
  \item Confondre \cle{personnalité} (notion psychologique : ce qui singularise) et \cle{personne} (notion juridique et morale : ce qui fonde la dignité et les droits).
  \item Attribuer à Kant l'idée que la valeur de la personne dépend de ses mérites ou de son utilité : c'est exactement le contraire, la dignité est inconditionnelle.
  \item Réduire Mounier à un éloge de la vie en groupe : son personnalisme affirme la \emph{primauté de la personne} sur la société, et non l'écrasement de la personne par la collectivité.
  \item Opposer trop simplement inné et acquis : le programme attend l'analyse de leur \emph{interaction}.
  \item En dissertation, annoncer un plan dialectique puis rédiger un catalogue d'opinions sans confrontation des thèses.
\end{itemize}
\end{attention}

\begin{aretenir}[Checklist avant le Probatoire / Baccalauréat technique]
\begin{itemize}
  \item $\square$ Je sais définir \textbf{personnalité}, \textbf{personne}, \textbf{tempérament}, \textbf{caractère}, \textbf{dignité}.
  \item $\square$ Je distingue clairement les \textbf{facteurs biologiques} (inné, hérédité, physiologie) et les \textbf{facteurs socio-culturels} (famille, école, culture, groupe des pairs, médias).
  \item $\square$ Je sais illustrer l'interaction du biologique et du social par un exemple camerounais concret (rites d'initiation, éducation familiale, bilinguisme, vie scolaire).
  \item $\square$ Je restitue la thèse de \textbf{Kant} : la personne est une fin en soi, douée de dignité et d'autonomie ; l'impératif catégorique commande d'agir par devoir et de traiter l'humanité toujours comme une fin, jamais simplement comme un moyen.
  \item $\square$ Je restitue la thèse de \textbf{Mounier} : personnalisme communautaire, communion, engagement, distinction personne/individu, refus de l'individualisme.
  \item $\square$ Je compare Kant (valeur universelle et inconditionnelle de la personne) et Mounier (valeur existentielle, relationnelle et historique).
  \item $\square$ Je sais expliquer pourquoi la personne a une \textbf{dignité} et non un \textbf{prix}.
  \item $\square$ J'applique la distinction \textbf{personnalité (psychologique)} / \textbf{personne (juridique et morale)} à un cas pratique (mineur, malade, personne morale comme une entreprise).
  \item $\square$ Je structure une dissertation : introduction (accroche, définitions, problématique, plan), développement dialectique (thèse, antithèse, synthèse), conclusion (réponse et ouverture).
  \item $\square$ Je structure un commentaire de texte : présentation, explication détaillée suivant le texte, discussion critique, conclusion.
  \item $\square$ Je mobilise un vocabulaire précis : \emph{autonomie/hétéronomie, impératif catégorique/impératif hypothétique, réification, aliénation, communion, engagement}.
\end{itemize}
\end{aretenir}

\findecours

\end{document}
